\documentclass[11pt]{article}
\usepackage[T1]{fontenc}
\usepackage[utf8]{inputenc}
\usepackage{lmodern}
\usepackage[margin=1in]{geometry}
\usepackage{amsmath,amssymb,amsthm,mathtools,bm}
\usepackage{booktabs,array,longtable}
\usepackage{enumitem}
\usepackage{hyperref}
\usepackage{xurl}
\usepackage{seqsplit}
\newcommand{\OPHPaperReleaseID}{r2011}
\newcommand{\OPHPaperReleaseDate}{August 1, 2026}
\newcommand{\OPHPaperReleaseBanner}{%
\begin{center}
\small\textbf{Paper release:} \texttt{\OPHPaperReleaseID}\quad\textbf{Released:} \OPHPaperReleaseDate
\end{center}
}
\newcommand{\PaperReleaseID}{\OPHPaperReleaseID}
\newcommand{\PaperReleaseDate}{\OPHPaperReleaseDate}
\newcommand{\PaperReleaseBanner}{%
\begin{center}
\small\textbf{Paper release:} \texttt{\PaperReleaseID}\quad\textbf{Released:} \PaperReleaseDate
\end{center}
}


\hypersetup{colorlinks=true,linkcolor=blue,citecolor=blue,urlcolor=blue}
\setlength{\emergencystretch}{3em}
\setlength{\LTleft}{0pt}
\setlength{\LTright}{0pt}
\setlength{\tabcolsep}{5pt}
\renewcommand{\arraystretch}{1.08}

\newtheorem{theorem}{Theorem}[section]
\newtheorem{corollary}[theorem]{Corollary}
\newtheorem{proposition}[theorem]{Proposition}
\newtheorem{lemma}[theorem]{Lemma}
\newtheorem{definition}[theorem]{Definition}
\newtheorem{axiom}[theorem]{Axiom}
\newtheorem{assumption}[theorem]{Assumption}
\newtheorem{remark}[theorem]{Remark}

\newcommand{\SU}{\mathrm{SU}}
\newcommand{\U}{\mathrm{U}}
\newcommand{\ellbar}{\bar{\ell}}
\newcommand{\Conf}{\mathrm{Conf}}
\newcommand{\ophid}[1]{\path{#1}}
\newcommand{\decnum}[1]{\vcenter{\hbox{\parbox{0.68\linewidth}{\small\ttfamily\seqsplit{#1}}}}}

\title{Deriving the Particle Zoo from\\
Observer Consistency}
\author{B.\ M\"uller \and Alexander Osika \and Mario Poneder \and Kai Xue \and Maarten Antoine Visser \and David Matscheko}
\date{\OPHPaperReleaseDate}

\usepackage{microtype}
\begin{document}
\maketitle
\OPHPaperReleaseBanner

\begin{abstract}
This paper carries the Standard Model structure derived by
Observer-Patch Holography (OPH) into particle calculations. Complete
reversible port response and endogenous observer transport force the local
Standard Model gauge Lie algebra. A conditional finite matrix realization
gives the charge lattice and maximal faithful matter image
\((\mathrm{SU}(3)\times\mathrm{SU}(2)\times\mathrm{U}(1))/\mathbb Z_6\)
for its declared matter table. Source selection of the matrix current, matter
action, and global quotient, together with laboratory current
and flux attachment, physical seam and family attachment, four-dimensional
topology, scalar dynamics, and continuum quantum-field realization, are
premises of any stronger physical reading.

The selected branch gives exact charge quantization, a dimensionless
electroweak hierarchy, a zero Higgs naturality defect in its declared chart,
and conditional Higgs and top coordinates. The source chain emits the
electroweak chart coordinates \((80.330,\,91.119)\)~GeV. Separate strict
one-loop receipts give chart-local scalar zero checks for the \(W\) and \(Z\)
channels. Their external Standard Model fixture is not composed with the OPH
chart, and neither declared chart is identified with a physical resonance
sheet. A
positive Hermitian \(C_3\) face circulant obeys the exact identity
\(Q=1/3+(2/3)(|b|/a)^2\), hence \(Q=2/3\) exactly when
\(|b|/a=1/\sqrt2\); a finite tracial construction supplies that balance under
its stated premises, and under the balance and mass-ordering premises the
measured electron and muon masses fix the tau mass inside a 72-eV enclosure
centered at \(1776.969027\)~MeV, \(0.43\) standard uncertainties from
measurement. Conditional on the declared Clebsch pattern, one-loop
transport gives
\(m_s/m_d=(m_\mu/m_e)_{\mu_U}/9=22.9743\). Against the two 2024
Flavor Lattice Averaging Group light-quark ratios it is \(12.8\%\) to
\(15.2\%\) high, rejecting that route.
The reciprocal quark and weighted-cycle neutrino candidates also fail their
tests. The six absolute quark masses are non-identifiable from the stated
source data, and absolute lepton and hadron masses are empirical inputs rather
than source-only results. No nonzero pole mass is claimed as a source-only prediction.
\end{abstract}

\section*{What This Paper Contributes}

The Standard Model gauge paper proves that complete reversible port response
and endogenous observer transport force the local Standard Model gauge Lie
algebra.
Its conditional matrix and matter realization supplies the hypercharge
lattice, a three-color carrier, the common \(\mathbb Z_6\) kernel, and the
maximal faithful matter image
\((\SU(3)\times\SU(2)\times\U(1))/\mathbb Z_6\).
It also supplies
the distinct conditional compact-gauge/Tannaka reconstruction route, a
conditional generation window \(3\le N_g\le5\), and the connection
and metric carrier roles. Under the single-complete-object, faithful-action,
and operational-cost premises, the exact screen-band theorem selects rank
three inside that window. The declared finite simulator recovers the same
projector at its lowest positive generator frequency. Tensoring that response
band with the declared generation table gives a conditional complex
rank-\(45\) candidate. Chirality and the diagonal \(\mathbb Z_6\) action are
properties of the table. A separate finite source-domain receipt checks the
declared operator \(D_\sigma\otimes I_{45}\) and conditional inheritance of a
positive dimensionless gap. The source does not select this matter action or
bridge the twelve-port Spin packet to the local operator domain.
Identification with physical matter poles and continuum fields is a
separate premise.
Classical massless
modes require explicit action, background, and phase premises; quantum poles
require a stronger particle certificate. This paper studies the downstream
particle readout after the local closure coordinate is supplied.

The displayed \(\mathbb Z_6\) quotient is the maximal faithful image of the
declared representation. It is not a source-selected physical global form.
Two structural routes feed that downstream calculation. Incidence expresses
\(J\) as a polynomial in adjacency, while target-blind port readback derives
the signed response. For the finite-current theorem, write
\(\mathfrak g=D(\mathbb R^{12})\subset\mathfrak u(H)\). Faithfulness gives
dimension twelve, and commutator closure makes \(\mathfrak g\) compact
reductive. Transitivity gives one fixed line for the alternating group on five
letters, \(A_5\). Endogenous proper-carrier transport makes this action inner.
It therefore fixes the centre pointwise. The centreless alternative is
\(\mathfrak{su}(2)^4\), whose factorwise fixed dimensions are zero or three
and cannot sum to one. The centre has dimension one, and compact-simple
classification splits the eleven-dimensional semisimple part as \(3+8\).
Thus
\[
\mathfrak g\cong
\mathfrak u(1)\oplus\mathfrak{su}(2)\oplus\mathfrak{su}(3).
\]
This argument assumes no ambient continuous gauge group.
The released matrix current and matter chain are conditional. Transportable sectors plus
Tannaka reconstruction give a separate conditional compact group. The
target-free matter reduct does not select the generation count inside the
window. Under separate complete-band and cost-order premises, the screen-band
theorem selects rank three. Tensoring that response band with the declared
generation table gives a conditional rank-\(45\) candidate. A separate local
receipt checks a declared tensor-identity operator and conditional gap
inheritance. It neither selects the matter action nor transports the
twelve-port Spin packet to the local domain. Matter-pole identification,
continuum Spin/locality, the physical seam choice, the third persistence leg,
and laboratory identification are premises of any physical reading.
Completeness for extra light sectors
is a separate attachment premise. The intertwiner between the two current
objects is likewise a premise.

A single screen-cell coordinate \(P_\star\) feeds the electromagnetic,
electroweak, and mass branches. The exterior matter witness fixes four weak
doublets. Conditional on the shared physical load carrier, the selected branch
emits \(v/E_\star\) with Higgs naturality defect \(\epsilon_H=0\); a weak scale
in GeV also requires an independently closed \(E_\star\). Each declared pixel
map has a machine-certified interval proof of a unique fixed point. On the
target-anchored empirical-closure branch, the measured charged-lepton triple
lies inside every outward-rounded arithmetic enclosure. These are diagnostic
intervals rather than prospective prediction intervals. The branch inverts exactly at the witness: one
anchor-gap value closes all three masses on the measured triple at once,
inside the certified band. The scheme bridge is therefore a declared
confirm-or-refute target. Absolute
charged-lepton and first-principles hadron masses lie outside the theorem,
and empirical endpoint data are identified wherever a transport bridge uses
them.

\section{Introduction}

Observer-Patch Holography (OPH) asks a concrete follow-up question after the gauge structure is fixed.
Can one local screen constant organize the particle spectrum, or are the masses and mixings a
collection of unrelated inputs?

The Standard Model gauge paper separates the local Lie algebra from its
matrix and matter realizations. Complete reversible response and endogenous
transport force
\[
\mathfrak{su}(3)\oplus\mathfrak{su}(2)\oplus\mathfrak u(1).
\]
Conditional on the separately declared matrix current, rank-15 matter table,
and descent packet, the maximal faithful matter image is
\[
\frac{\SU(3)\times\SU(2)\times\U(1)}{\mathbb Z_6},
\qquad
N_c=3.
\]
The exact hypercharge lattice and color count belong to this conditional
realization. Separate intrinsic-CP and weak-sector premises give the window
\(3\le N_g\le5\), which does not by itself select the generation count.
The graph carrying the alternating group on five letters, \(A_5\), together
with anomaly cancellation and the target-free matter reduct does not
select a count. A strengthened screen interface does: among single complete
faithful bands in the window, the exact Laplacian costs
\(5-\sqrt5<6<5+\sqrt5\) select the rank-three band uniquely. The declared
unitary simulator recovers the same residue at its lowest positive generator
frequency. Tensoring this mode with the declared generation table gives a
conditional rank-\(45\) candidate. Chirality and the diagonal
\(\mathbb Z_6\) action come from that table. A separate local-domain receipt
checks a declared tensor-identity operator and conditional gap inheritance.
The source does not select the matter action or bridge the twelve-port Spin
packet to that domain. Physical family interpretation requires matter-pole,
continuum, seam-selection, persistence, and laboratory receipts.
The paper also supplies the exact
hypercharge lattice and a conditional exterior representation witness. If a
trace-balanced block carrier
\(V=C\oplus W\), with dimensions \(3+2\) and hypercharges
\((-1/3,1/2)\), is physically selected, then
\(\Lambda^2V\oplus\Lambda^4V\) branches exactly as
\(Q\oplus u^c\oplus e^c\oplus d^c\oplus L\). It has the three
one-Higgs invariant lines, cancels all five gauge and mixed anomalies, and
contains four weak doublets per generation. The determinant balance, Spin
lift, non-vacuum exterior package, and center/deck descent are
receipt-checked under the response and matter contracts. The scalar scan fixes
compatible charges and Yukawa channels, not \(H=W\) or scalar multiplicity.
Laboratory identification of the current and flux sectors, exclusion
of extra light sectors, attachment of the \(A_5\) face representation to
physical families, scalar attachment/dynamics, and continuum QFT are
separate premises.
The presentation-invariant normal-form framework of Ref.~\cite{observableNormalForms} isolates the
same-source confluence, cross-source boundary-identification, liveness, and refinement criteria
used upstream. It quotients hidden coordinates, labels, worker layout, and ancillas when they leave
observer-visible records unchanged. Visible port incidence, topology, response maps, and repair laws
remain physical branch data. The framework supplies no particle-sector selector or numerical
particle output; every such claim below is conditional on the specialist gauge branch and the
particle-specific receipts.

The particle branch is a one-fixed-point forward reconstruction problem with an
explicit closure matrix. It propagates one dimensionless pixel ratio \(P\)
through the spectrum. The
construction does not emit photon, gluon, or graviton particle masses from symmetry labels. It records
their conditional classical carrier modes together with the quantum-particle certificates those modes consume as premises, while the Higgs
candidate is on its declared quantitative surface. The hierarchy/naturality bridge is closed on its
selected branch. The restricted quark source-spread non-identifiability theorem, the common-scale
rejection of its reciprocal-ray candidate, and the weighted-cycle neutrino branch are
visible as obstruction or comparison surfaces. Target-anchored quark mass textures, the mixed-convention
formula diagnostics, and compare-only absolute neutrino attachments are withheld from public
prediction tables. The \(P\)-closure root,
the electroweak \(W/Z\) rows, charged leptons, and hadrons obey the
sector-specific boundaries stated below.

\subsection{Structural particle content and its attachment boundary}

On the declared branch, particle carrier roles are constrained in three
separate layers:
\begin{enumerate}[leftmargin=2em]
\item overlap consistency on the observer-patch network;
\item the Axiom~1/Axiom~2 fixed-space theorem for the local Standard Model
gauge Lie algebra, followed by the conditional matrix current, rank-15 matter
table, anomaly balance, and tensor descent;
\item the conditional rank-three screen selection, its abstract tensor product
with the generation table, and the separate declared local operator context.
\end{enumerate}
The first layer says what kinds of local data can be compared consistently across neighboring
patches. The second forces the local Lie algebra. Its declared matrix and
matter realization gives the exact charge lattice, color carrier, common
\(\mathbb Z_6\) kernel, and maximal faithful image conditionally.
The third separates the finite carrier theorem from its declared physical
completion.
The result is an exact contract-conditional finite carrier and one-generation
matter packet. The selected response band and generation table give a
conditional rank-\(45\) candidate. The separate local-domain branch verifies
the declared tensor-identity operator without selecting it from the source.
A physical particle interpretation requires source binding, laboratory
current identification, matter-pole and continuum attachment, interaction,
scalar attachment/dynamics, and QFT receipts.

The transportable-sector/Tannaka route is a separate conditional
compact-group reconstruction. Physical source binding of the finite current
and equality with the Tannaka current are premises of any physical
identification. The conditional rank-\(45\)
candidate does not become a physical family without source-selected local
action, cross-domain Spin transport, matter-pole, continuum, seam-selection,
persistence, and symmetry-descent receipts.

\emph{Recovering Observer Spacetime and Einstein Dynamics from Overlap Consistency and Deriving Standard Model Gauge Structure from Observer Overlap Consistency}~\cite{ophEinstein,ophGauge} is therefore central to this paper. Complete
reversible response and endogenous transport force its local gauge Lie algebra
\[
\mathfrak{su}(3)\oplus\mathfrak{su}(2)\oplus\mathfrak u(1).
\]
The separately declared matrix and matter realization gives the maximal
faithful image and color count
\[
\frac{\SU(3)\times\SU(2)\times\U(1)}{\mathbb Z_6},
\qquad
N_c=3.
\]
The same conditional packet contains one scalar-doublet channel and a rank-15
internal matter witness. The charge lattice and color count use the stated matrix and matter
premises. The quotient is the maximal faithful image of that representation,
while selection of the physical global form is a separate premise. Separate intrinsic-CP
and weak-sector premises give \(3\le N_g\le5\). The target-free matter reduct
leaves the three-generation value unselected. On the strengthened screen
interface, exact cost
minimization selects the canonical rank-three band and the declared simulator
realizes its response residue. Its tensor product with the declared
generation table is a conditional rank-\(45\) candidate. A separate finite
local-domain receipt checks a declared tensor-identity operator. It does not
select that action or identify the local domain with the Spin packet. Physical
matter-pole and continuum attachments require separate derivations.

On that selected one-Higgs branch the scalar carrier has a canonical local screen-chart model.
Let \(C_{\rm EW}\cong\mathbb{CP}^1\) be the support-visible electroweak chart and fix the positive
Hopf line-bundle convention by the neutral component condition \(Q(\phi^0)=0\). Borel--Weil gives
\[
H_{\rm OPH}=H^0(C_{\rm EW},\mathcal O(1))\cong\mathbb C^2,
\]
the first nontrivial holomorphic section space. With OPH's hypercharge and \(\mathbb Z_6\)
normalization this is exactly the \((1,2)_{1/2}\) Higgs carrier. Projectivization classifies a
nonzero section direction as a point of \(\mathbb P(H_{\rm OPH})\cong\mathbb{CP}^1\), but it
forgets the scalar hypercharge phase. For the lower-component vacuum vector
\[
\phi_0=\frac{v}{\sqrt2}\binom{0}{1},\qquad v\ne0,
\]
one has
\[
e^{i\alpha T_3}e^{i\beta Y}\phi_0
=e^{i(\beta-\alpha)/2}\phi_0.
\]
Consequently \([\phi_0]\) has the projective two-torus stabilizer
\(\U(1)_{T_3}\times\U(1)_Y\), modulo the inherited finite center, whereas the vector
\(\phi_0\) is fixed only when \(\beta=\alpha\), locally. Its connected stabilizer is therefore the
electromagnetic diagonal \(\U(1)_Q\), generated by \(Q=T_3+Y\). Thus the projective geometry
explains the carrier-ray classification, while the chosen nonzero vector supplies the
symmetry-breaking statement. This explains the representation, charge, and symmetry-breaking
geometry of the one-Higgs slot. It does not derive \(m_H\), the quartic, \(v\), or
Coleman--Weinberg dynamics; the
weak scale, Higgs/top quantitative surface, and \(\epsilon_H=0\) hierarchy/naturality closure are
the OPH quantitative branch.

This also delimits the phrase ``and no others.'' The finite OPH claim concerns
the low-energy packet implied by the stated response, matter, and descent
contracts. Additional connected gauge generators,
additional light Higgs multiplets, extra light chiral families, or low-energy supersymmetric
partners are absent from the declared completion because that completion
declares them absent. This is a declaration, not source-derived
exhaustion of physical sectors. Propagating gauge-particle existence requires
the complete primitive data, action, phase, and quantum-pole premises below.

The individual families then emerge for different reasons. The realized electromagnetic, color,
and dynamical-metric branches identify connection and metric carrier roles. Explicit Maxwell,
perturbative pure-Yang--Mills, and pure-Einstein quadratic actions then yield classical transverse
or TT massless modes on their stated backgrounds and phases. Quantum photon, gluon, or graviton
states require the separate physical-Hilbert-space, pole-residue, and asymptotic/phase receipt. The weak bosons arise when
the electroweak gauge sector is propagated through its quantitative closure branch. Quarks and leptons arise from the
conditional chiral matter packet together with the rank-three screen
selection. Their physical three-family interpretation requires the
conditional rank-\(45\) tensor candidate to be identified with physical
matter-pole residues and continuum fields, with the local matter action and
cross-domain Spin transport selected. Their
family splittings are then read from the deeper overlap-transport and excitation machinery, not
from a second arbitrary postulate that says ``copy the family three times and assign masses by
hand.'' Once color is realized and confined, stable hadrons are composite readout channels of the
quark/gluon sector.

The role of the pixel constant \(P\) is also important to state clearly. Here, \(P\) does not
decide whether an electromagnetic particle pole exists, whether there are three colors, or whether
the candidate family band is physically attached. The electromagnetic representation role,
the color carrier, and the declared generation completion belong to the finite packet before the
quantitative closure step. What \(P\)
does is set the shared quantitative scale on which receipt-certified content is read out
numerically. So the logic of the paper is:
\[
\begin{aligned}
\text{observer consistency}
&\to
\text{realized gauge/matter branch} \\
&\to
\text{representation roles}
\to
\text{conditional action-level carrier modes} \\
&\to
\text{common scale }P
\to
\text{receipt-gated quantitative spectrum}.
\end{aligned}
\]
That is how the derivation explains why a universe with the realized branch and the
shared pixel scale exhibits this particle zoo.

\subsection{Logical basis and reading rule}

The particle derivation uses the same OPH basis as the SM/GR derivation in
Ref.~\cite{ophEinstein,ophGauge}. Its three axioms are the oriented
twelve-port observer screen, observer agreement, and
conditional maximum randomness. On that basis, the
SM/GR derivation in Ref.~\cite{ophEinstein,ophGauge} supplies the structural chain used
here: the local Standard Model gauge Lie algebra forced by complete response
and endogenous transport; the conditional maximal faithful matter image,
hypercharge lattice, and color triplet \(N_c=3\) under explicit matrix and
matter contracts; a separate
receipt-conditional Tannaka
classification; and the conditional generation window. Under the named
complete-band and operational-cost premises, an exact screen theorem selects
the rank-three response band, and the declared unitary response places its
residue at the lowest positive generator frequency. Extra-light-sector
completeness, laboratory current identification, and physical attachment of
that finite band to three matter families are premises the chain does not
supply.
Refs.~\cite{ophconsensus,ophmicrophysics} supply the patch-net, repair, measurement, and
regulated screen language for flavor transport and observer-facing readout.

The quantitative particle side adds one local closure variable. The common pixel ratio \(P\),
reported by the synthesis paper's incomplete outer/inner declared map~\cite{ophsynthesis}, feeds the forward
electroweak map and its descendants. The public empirical comparison branch displays
\[
\alpha^{-1}(0)=137.035999177(21),
\qquad
\alpha(0)\simeq0.00729735256433,
\qquad
P_C\simeq1.6309682094.
\]
The comparison pixel \(P_C\) is defined from the measured endpoint; it is not a source-root
coordinate.
The computation has a fixed source order: golden-ratio entropy balance gives
\(\varphi\), boundary Gaussian normalization supplies the \(\sqrt{\pi}\) width,
a trial \(P\) feeds the source map through unification, running, and electroweak
anchoring, and Ward-projected electromagnetic transport gives the Thomson
endpoint used by the outer/inner pixel fixed point.
The declared numerical map fixes a unique root on its certified interval. The map is
incomplete because no derivation identifies its numerical endpoint with the physical
electromagnetic readout. The
first-principles diagnostic trunk records the certified coordinates
\[
\begin{aligned}
P_{\mathrm{fwd}}&=1.630972095858897\ldots,\\
\alpha_{\mathrm{root}}^{-1}
&=136.994835177413\ldots.
\end{aligned}
\]
The root is the interval-certified unique fixed point of the declared numerical map.
The approximately \(0.041\) inverse-alpha difference to the public Thomson endpoint is a residual
of that incomplete map. It neither derives a QCD/hadronic contribution nor identifies a physical
endpoint relation.
The QCD-free hierarchy witness is the cleaner first-principles stress test. Combining the source/root
value with the finite-screen unified gauge-width contribution evaluated at the CODATA-derived
comparison pixel, \(\alpha_U(P_C)\), gives the mixed-provenance no-hadron diagnostic
\(A_{\alpha_U}^{\mathrm{fp}}=137.0359595136\ldots\), below the Thomson endpoint. It is
mixed-provenance comparison bookkeeping rather than a source-only fine-structure prediction. The certified
self-consistent gauge-width fixed point is \(\alpha^{-1}=137.035660136946577\ldots\);
the mixed diagnostic combines the inner value at \(P_{\mathrm{fwd}}\) with
\(\alpha_U\) at the CODATA-derived comparison pixel and is excluded from the physical result.
The detailed endpoint table appears in Section~\ref{sec:particle-input-ledger}.
The separate proposed
cosmic record-closure target
\(\mathfrak F_{r,0}(D_\star)=\{D_\star\}\), with
\(N_{\mathrm{CRC}}=\log D_\star\), defined through
the correctable code of the reachable public record atoms under globally coupled checkpoint
kernels, belongs to the cosmological-capacity branch and
supplies no particle theorem. The Newton coupling uses the separate
selected no-\(G\) scale certificate \(\gamma_\star=\ell_\star\nu_{\mathrm{Cs}}/c\), equivalently
\(B_\star=3\pi/\ell_\star^2\), so the particle branch does not derive \(G\) by back-solving the
pixel ratio. A fixed-\(D=24\) public checkpoint packet, carrier
representation, whole-fiber scalarization, and finite extension and refinement
receipts are exact. A bounded all-rung counterfamily has incompatible zero
sets under shared base, positivity, carrier, and executable finite controls.
The exact arithmetic establishes nonidentifiability for that bounded
completion class. Universal all-rung membership in the complete A1--A3
capacity-source contract and an executable-to-Lean bridge are absent. Direct
\(N\) is not evaluable on the incomplete source antecedent, and the stronger
source-class verdict is open. A positive direct result requires completion of
that antecedent, one physical slack zero, and a universe-level carrier
attachment. Its identification
with the electroweak bridge is conditional on a
positive, unital, refinement-natural identification of the screen load with
the electroweak load; the bridge value
\(N_0=3.5321315434\ldots\times10^{122}\) on the source-forward branch is
about \(6.6\) percent above the weighted Planck base-\(\Lambda\)CDM capacity
\(N_\Lambda=3.3129270981\ldots\times10^{122}\). Under the separately
declared one-class finite-presence attachment, this becomes
\(3.2920978773\ldots\times10^{122}\); the distinct Poisson carrier gives
\(3.3000722254\ldots\times10^{122}\). Their respective residuals are
\(-0.63\) and \(-0.39\) percent. Exact positive compositional completions of
the same local datum select neither global action nor a blocked-event
semantics. The named-law branch is therefore not evaluable on this source
class. Both comparisons are retrospective. If the missing
physical identifications are supplied, the pair \((P_\star,N_{\mathrm{CRC}})\) determines
dimensionless curvature products such as
\(\Lambda_\star a_{\mathrm{cell}}=3\pi P_\star/N_{\mathrm{CRC}}\), not the SI scale product
\(\Lambda_\star N_{\mathrm{CRC}}\) by itself. The rounded
\(N_\Lambda\simeq3.313\times10^{122}\) display is a cosmological central-value label, not the
high-precision input for \(G\). Structural carriers, quantitative outputs, conditional calculations, and
continuation calculations therefore inherit the declared theorem premises, with
particle-facing quantitative burden carried by the local declared-map root \(P_\star\).
The public fine-structure display row is a comparison endpoint conditioned on the measured
Thomson value. The first-principles diagnostic row and empirical hadron closure
rows are separate support records.
Section~\ref{sec:particle-input-ledger}
records that boundary in full.

\section{P-Closure and the Reverse-Engineering Strategy}

The reverse-engineering claim is easiest to state plainly. In
ordinary phenomenological use, the Standard Model does not derive the observed particle masses and
mixings from one common microscopic constant. It is usually presented with \(19\) free parameters
in the minimal massless-neutrino theory, and with at least \(26\) once neutrino masses and mixing
are included, depending on neutrino-sector conventions. OPH organizes
that situation around one universal pixel ratio \(P\) from the outer/inner closure relation
described in the synthesis paper. Whether any parameter compression is realized follows
from counting the declared settings against the independent landed basins. The same \(P\) must drive all
downstream bosonic, quark, lepton, neutrino, and hadronic branches.

One clarification matters. The broader self-closure formulation uses the
universe-level equation
\[
N=\log M_0(\mathfrak U_N),
\]
where \(M_0\) is the multiplicative size of the largest correctable public
record code. Its selector-free finite form is stable global screen capacity
\[
\mathfrak F_{r,0}(D_\star)=\{D_\star\},
\qquad N_{\mathrm{CRC}}=\log D_\star,
\]
for a universe supplied with logarithmic capacity \(N\). The synthesis paper
specifies the finite readback directly by
\(M_0(q)=\alpha(G_q)\), the independence number of the compound confusability
graph of the reachable public record atoms under the globally coupled
checkpoint family. Under whole-fiber scalarization and a faithful capacity
carrier it is deflationary; under confusability-reflecting capacity extension it
is monotone, so top-down iteration reaches the greatest fixed point on a
declared finite chain. The fixed-\(D=24\) packet is exact inside its declared
finite source category. A bounded all-rung counterfamily has incompatible
fixed sets under shared base agreement, positivity, the carrier bound, and
executable finite controls. This establishes nonidentifiability for the
bounded completion class. It does not establish universal all-rung membership
in the complete A1--A3 capacity-source contract or an executable-to-Lean
bridge. Direct \(N\) is not evaluable on the incomplete source antecedent, and
the stronger source-class verdict is open. A positive result requires a
complete source antecedent, one physical slack zero, and a universe-level
carrier attachment.
The electroweak identification additionally requires the common
screen/electroweak load-carrier hypothesis. The conditional bridge value
\(N_0=3.5321315434\ldots\times10^{122}\) is about \(6.6\) percent above the
weighted Planck base-\(\Lambda\)CDM capacity
\(N_\Lambda=3.3129270981\ldots\times10^{122}\), a propagated \(2.4\) to
\(2.5\) one-dimensional sigma. The conditional finite-presence and Poisson
branches instead lie \(0.63\) and \(0.39\) percent below that coordinate,
respectively. Exact positive composition countermodels select neither global
action from the finite source. Their physical carriers are distinct and both
comparisons are retrospective. Read through
\(\Lambda_\star\ell_\star^2=3\pi/N\), the same bridge is the zero-dial relation
\(\Lambda\ell_P^2=3e^{-6\pi/(P\alpha_U(P))}\) only after the independent
capacity, horizon, and common-load receipts. Neither reserve factor is part of
the common-load theorem itself. The same exponent
runs the weak hierarchy of this paper. The
particle-spectrum derivation studied here uses the local declared-map root \(P_\star\) in place of
a long particle-by-particle parameter list.

Observation is allowed to supply a branch hint. That is a consequence of the OPH picture, where
the universe is a closed fixed structure and internal observers can read approximate coordinates
from experiments. The proof standard is stricter: the exact \(P_\star\) used by the particle
branch is computed by the self-referential numerical pixel equation, and the root is unique on the
declared interval. A further derivation is required before that root is a physical endpoint.
The pixel equation and its uniqueness statement are carried by the shipped
interval contraction certificate: the closure map is a self-map with a
derivative bound certified by interval arithmetic on the declared interval. The
domain-global statement is discharged as well: the companion domain-global certificate bounds
\(\sup|g'|<1\) on every piece of a 256-piece subdivision of the declared numerical domain
(\(\alpha^{-1}\in[100,200]\), both readout maps, empty exceptional set), so each declared readout
map has exactly one fixed point on that numerical domain. This certificate does not provide the
missing physical endpoint relation.

\subsection{Why a single \(P\) matters}

The structural theorems determine an exact contract-conditional finite
representation packet, but not \emph{what physical particle world} realizes it
or what absolute scale in GeV it occupies. The overlap,
modular, gauge, anomaly, and admissibility arguments give the declared Standard Model packet,
the exact hypercharge pattern, and a three-color carrier. The generation
window is conditional; separate band premises select rank three. Tensoring the
selected response band with the declared generation table gives a conditional
rank-\(45\) candidate. A distinct local-domain receipt checks a declared
tensor-identity operator and conditional gap inheritance without
source-selecting the action or transporting the Spin packet. Physical
matter-pole and continuum interpretation is a separate premise. The
carrier-mode theorem is an
additional action-level result and does not emit zero-GeV quantum particles. These structural
results do not by themselves tell us the numerical values
of the \(W\) boson mass, the Higgs boson mass, or the up-quark mass.

The role of the pixel closure is to supply one common quantitative scale
variable for the entire downstream spectrum. In the synthesis paper, the
fine-structure branch asks for the nonzero detuning of a holographic screen cell
such that the cell's outer geometric displacement from perfect self-similar
equilibrium equals the electromagnetic observation scale emitted by the
encoded branch under the common physical carrier identification. That
identification is a separate receipt.
The outer side of the closure is
\[
P=\varphi+\alpha_{\mathrm{in}}(P)\sqrt{\pi}.
\]
The first-principles computation is a five-step source chain. First, the
golden-ratio entropy balance of the local screen cell supplies
\(\varphi=(1+\sqrt5)/2\). Second, boundary maximum-entropy normalization fixes
the width of the leading electromagnetic detuning to \(\sqrt{\pi}\). Third, a
trial \(P\) is sent through the source map
\[
M_U(P)=E_P e^{-2\pi}P^{1/6},\qquad E_{\mathrm{cell}}(P)=\frac{E_P}{\sqrt P},
\]
followed by the heat-kernel closure
\[
\bar\ell_{\mathrm{SU}(2)}(t_2(P))+\bar\ell_{\mathrm{SU}(3)}(t_3(P))=\frac{P}{4},
\]
which selects \(\alpha_U(P)\) and the running family
\(\alpha_i(m_Z;P)\). Fourth, electroweak mixing gives the source anchor
\[
a_0(P)=\alpha_2^{-1}(m_Z;P)+\frac53\alpha_1^{-1}(m_Z;P).
\]
Fifth, Ward-projected \(\U(1)_Q\) transport gives the Thomson endpoint
\[
A_T(P)=T_Q(a_0(P),F_{\mathrm{src}}(P))=\alpha_{\mathrm{em}}^{-1}(0;P),
\]
and the cell closes when
\[
P=\varphi+\frac{\sqrt{\pi}}{A_T(P)}.
\]
The numerical solve is therefore a root problem for
\[
H(P):=P-\varphi-\frac{\sqrt{\pi}}{A_T(P)}.
\]
The branch interval is localized from the observer-facing data, the solver evaluates \(A_T(P)\)
from the declared source map at each candidate point, and the accepted value is the unique zero of
\(H\) on that interval. The measured endpoint can locate the interval. It cannot replace the root
condition.
This proves uniqueness only for the declared numerical map. The derivation connecting its root to
the observer-supporting physical electromagnetic endpoint is missing.
The CODATA-conditioned comparison readout
\(\alpha^{-1}(0)=137.035999177(21)\) gives
\[
P\simeq1.6309682094.
\]
The same equation is the local ruler for the downstream particle rows. The five-layer
OPH diagnostic trunk emits the certified source-side diagnostic point
\[
\begin{aligned}
P_{\mathrm{fwd}}&=1.630972095858897\ldots,\\
\alpha_{\mathrm{root}}^{-1}
&=136.994835177413\ldots.
\end{aligned}
\]
The root is an incomplete declared-map output, not a physical fine-structure
prediction.
The detailed table records how the diagnostic trunk, endpoint residual, and
source-spectral payload fit together.
A separate hardware note reports an optical-cavity check of the same declared
geometry. It is a non-discriminating engineering check and supplies no
physical endpoint.
Once the diagnostic map coordinate is set, the particle question becomes whether all downstream readouts
can be written as
\[
X_j = G_j(P)
\]
with no new sector-specific constants inserted by hand. This paper studies exactly that downstream
map.

\subsection{How the electroweak branch works}

In the implementation used here, the quantitative electroweak branch is the first major readout from \(P\). Its algebraic
output family is
\[
\{M_W^{\mathrm{chart}},\ M_Z^{\mathrm{chart}},\
\alpha_{\mathrm{em}}^{-1}(q^2),\ \sin^2\theta_W(q^2),\ v\},
\]
evaluated from one declared input \(P\) on the printed
running/matching/threshold/scheme package. The superscript
\(\mathrm{chart}\) marks a mass-chart coordinate, not a certified complex
pole. The construction
is straightforward in spirit: start from \(P\), build the electroweak running family, select the physical carrier point on that family, and read off the \(W\) boson and \(Z\) boson pair from that selected point. The Higgs/top critical stage
inherits that same electroweak core and does not introduce a new
free-input sector. On the declared Ward-projected transport branch, the low-energy electromagnetic row is represented by the Thomson endpoint
\[
\alpha_{\mathrm{Th}}^{-1}(P)=\lim_{q^2\to 0}\alpha_{\mathrm{em}}^{-1}(q^2;P)
\]
of that same electromagnetic transport family.
The forward logical order on this branch is:
\[
\begin{aligned}
P
&\longmapsto
\bigl(M_U(P),E_{\mathrm{cell}}(P)\bigr)
\longmapsto
\alpha_U(P)
\longmapsto
\bigl(t_U(P),t_{\mathrm{tr}}(P)\bigr)\\
&\longmapsto
\bigl(t_2(P),t_3(P),v_{\mathrm{chart}}(P)\bigr)
\longmapsto
\alpha_i(\mu_\ast;P).
\end{aligned}
\]
Here \(v_{\mathrm{chart}}(P)\) is the mass-chart coordinate. It is not a
renormalized vacuum expectation value in any fixed scheme, and no receipt
identifies it with one.
The forward transmutation certificate records that the same runtime basis reconstructs the unified diffusion parameter \(t_U(P)=4\pi^2\alpha_U(P)\) and transmutation exponent \(t_{\mathrm{tr}}(P)=2\pi/((N_c+1)\alpha_U(P))\) as the pixel-closure solve itself. The runtime subgraph reads measured electroweak data only for validation. No frozen provenance receipt excludes target dependence or selects one repair law.

The electroweak branch is the place where the named pixel coordinate organizes the weak-sector rows. The
reference-fitted inverse adapter displays
\[
M_W=80.3625~\mathrm{GeV},
\qquad
M_Z=91.1879~\mathrm{GeV}.
\]
The selected-carrier chart, the runtime-target-separated formula candidate, and
the reference-fitted coherent repair diagnostic are recorded separately in the support
table below.

\subsection{What the electroweak branch fixes}

Under the declared electroweak running, matching, threshold, and scheme premises, the pixel variable \(P\) supplies the
source basis
\((\alpha_U,\alpha_{2,m_Z},\alpha_{Y,m_Z},\eta_{\mathrm{source}},v)\) used by the
candidate chart. The displayed value law is an exact implication of the
five-part quotient-path certificate stated below, but the finite carrier does
not emit that certificate. The candidate law evaluates the mass-side pair
\((W,Z)\) together with the running-family anchor \((a_0,s_0,v)\). The
electromagnetic row is physically read only after Ward projection to the unbroken \(\U(1)_Q\)
channel; its low-energy value is the \(q^2\to0\) endpoint of the same transport family. The
derivation also includes two explicit comparison calculations with distinct provenance: the
exact selected-carrier chart and the reference-fitted coherent repair diagnostic. The same support
logic applies to the flavor branches, whose premises differ. The electroweak
quantitative branch states its source payload, same-scheme remainder, and
interval certificate explicitly.

For empirical comparison, a documented piecewise
\(e^+e^-\to\mathrm{hadrons}\) compilation, built from resonance parameters and perturbative QCD,
evaluates the subtracted dispersion integral to
\(\Delta\alpha_{\mathrm{had}}^{(5)}(m_Z)=0.027609\pm0.000112\).
Inserting it into the endpoint map with the frozen source anchor and lepton transport packet gives
\(\alpha^{-1}=136.3827548175\) on
\([136.3670480603,136.3984651934]\), at \(P=1.6310415204\), on the
empirical-closure row class.
The comparison to the measured endpoint sits in an explicitly compare-only block: the measured
value lies outside the interval, and the certified same-scheme anchor gap is
\([0.6198609041,0.6505569679]\) inverse-alpha units. The anchor \(a_0(P)=\alpha_{\mathrm{em}}^{-1}(m_Z^2;P)\) is the
one-loop renormalization-group value run from the OPH unification scale to \(m_Z\); the physical
five-flavor on-shell value is \(128.939\), and the difference from the OPH anchor \(128.308\) is
\(0.631\), at the lower edge of the certified gap. The gap records transport missing from the
one-loop anchor; it is not evidence by itself that a source chain closes. A source-only bridge
that closes this gap reduces to the OPH hadronic spectral measure,
which needs a working hadron construction. The empirical endpoint is a
comparison calculation and does not supply that source-only bridge.
The exact exterior weak multiplicity is four. Its use as
the transmutation factor \(\beta_{\mathrm{EW}}=4\) is conditional on the
common screen/electroweak load-carrier identification; overloaded
\(\beta\)-ratios appear only on benchmark readouts and are not part of the theorem contract.

\section{Results at a Glance}

The particle derivation carries the local pixel scale into the directly comparable rows below.
Detailed scope statements are collected in the support sections that follow.

\begin{table}[h]
\centering
\scriptsize
\setlength{\tabcolsep}{3pt}
\begin{tabular}{@{}>{\raggedright\arraybackslash}p{0.30\textwidth}>{\raggedright\arraybackslash}p{0.37\textwidth}>{\raggedright\arraybackslash}p{0.25\textwidth}@{}}
\toprule
Chain & OPH output & Role \\
\midrule
\(P\)-closure and fine structure & certified incomplete-map root \(136.994835177413\ldots\) (certified source-root row); mixed diagnostic \(137.0359595136\ldots\); certified gauge-width fixed point \(137.035660136946577\ldots\); empirical endpoint \(136.3827548175\in[136.3670480603,136.3984651934]\); measured \(137.035999177(21)\) & distinct support classes: incomplete declared-map output, mixed diagnostic excluded from physical output, empirical hadron closure with anchor gap \([0.6198609041,0.6505569679]\), and measurement \\
Classical carrier modes & two transverse Maxwell modes; \(2\dim G\) perturbative pure-Yang--Mills modes; two pure-Einstein TT modes & conditional quadratic-action results; quantum particles require separate Hilbert-space, pole, and phase premises \\
Electroweak \(W/Z\) & source-only running/chart coordinates \((80.330,\,91.119)\)~GeV; strict one-loop consumer and separate scalar analytic receipts; candidate value-law coordinates \((80.3770000154,\allowbreak\,91.1879780779)\)~GeV & the analytic receipts exclude zeros on declared principal-sheet boxes and isolate, for each of \(W\) and \(Z\), one simple scalar zero with derivative and scalar-residue balls in its declared lower-half pole box on a channel-specific algebraic chart. They identify neither chart with the physical resonance sheet and prove no unique continuation, sign bridge, full-matrix Laurent residue, current amplitude, independent numerical replay, or W/Z mass comparison \\
Electroweak hierarchy\slash naturality & \(\begin{gathered}v/E_\star\ \text{on each named pixel branch}\\N_{\mathrm{EW}}(P_C)\simeq3.53235\times10^{122}\\N_0(P_{\mathrm{fwd}})\simeq3.53213\times10^{122}\\ \mathcal B_{\mathrm{EW}}=0,\ \epsilon_H=0\end{gathered}\) & selected dimensionless bridge; identifying either bridge coordinate with cosmic capacity requires a source-derived direct public-record closure with one selected physical capacity and the typed same-load bridge. The source-forward uncorrected value is \(6.6\) percent above the weighted Planck base-\(\Lambda\)CDM coordinate. The separately premised finite-presence and Poisson candidates lie \(0.63\) and \(0.39\) percent below it. Exact compositional countermodels select neither action. Both reserve comparisons are retrospective; the horizon--record identification and \(E_\star\) are separate premises \\
Higgs/top relation & double-criticality branch \((m_H,m_t)=(125.77,\,172.63)\)~GeV at the declared boundary-scale candidate \(E_\star e^{-\pi}P^{-1/6}\) (two loops), with \(m_H=125.72\)~GeV on the fit-free curve at the measured top; target-anchored declared-surface fit \((125.1995304097,\allowbreak\,172.3523553288)\)~GeV kept separate; no source-only physical mass emitted & the criticality family has zero continuous parameters over the boundary scale. Its scale selection assumes that the boundary record has exactly two equal-depth anchor registers, one geometric and one transmutational, together with the declared Gaussian/affine record and port-additive repair premises. The target-anchored fit carries no predictive force. This branch does not supply the source root, physical scale, rigidity, provenance, uncertainty enclosure, or complex pole \\
Charged leptons & target-anchored witness values withheld; finite eight-path digital carrier and empirical transport intervals retained as diagnostics & the model proves schema satisfiability and a central record dilation, not physical source selection; the intervals use target-anchored ratios, measured \(\alpha\), and empirical transport \\
Selected-frame quarks & public numeric rows withheld; reciprocal-ray product \(0.835323\) and held-out error \(21.56\%\) at \(M_Z\) & common-scale data reject the reciprocal-ray candidate; the generic interface has six scalars and requires a source-derived flavor-orbit selector \\
Neutrino absolute attachment & compare-only absolute masses withheld; scale-free and mixing comparison rows retained & weighted-cycle comparison branch \\
Hadrons & first-principles strong-binding descent plus empirical \(e^+e^-\to\mathrm{hadrons}\) closure surface & OPH hadron construction for first-principles masses; measured hadron data for empirical closure rows \\
Affine event-record stitching & certified affine event supports with separately typed conditioned-frame balls, followed by cross-boundary token continuation or \(\mathrm{AMBIGUOUS}\) & conditional observer-record theorem; not a species, mass, or coupling input \\
\bottomrule
\end{tabular}
\end{table}

\begin{table}[h]
\centering
\tiny
\setlength{\tabcolsep}{2pt}
\begin{tabular}{@{}>{\raggedright\arraybackslash}p{0.14\textwidth}>{\raggedright\arraybackslash}p{0.11\textwidth}>{\raggedright\arraybackslash}p{0.19\textwidth}>{\raggedright\arraybackslash}p{0.20\textwidth}>{\raggedright\arraybackslash}p{0.27\textwidth}@{}}
\toprule
Sector & Theorem scope & Derivation & Output(s) & Caveat / theorem boundary \\
\midrule
Carrier roles and modes & structural group/content theorem plus conditional action theorem & explicit transverse/TT quadratic kernels with no derived particle mass & action, background, and phase are hypotheses; quantization, positive-residue pole, and asymptotic/deconfinement conditions are separate \\
Electroweak bosons & running/chart calculation, strict one-loop consumer, scalar analytic receipts, conditional quotient transport, and comparison adapter & source-only zero-selector law; selected carrier \(\rightarrow\) two-coordinate chart; complete renormalized packet \(\rightarrow\) strict scalar consumer; finite quotient-path certificate \(\rightarrow\) conditional value law & source-only coordinates \((80.330,\,91.119)\)~GeV and candidate value-law coordinates \((80.3770000154,\allowbreak\,91.1879780779)\)~GeV are chart outputs; the analytic receipts concern an external fixture that is not composed with either chart & principal-sheet zero exclusion and one simple channel-specific lower-half algebraic-chart zero for each of \(W\) and \(Z\) do not identify a physical resonance sheet or supply a unique continuation, sign bridge, full-matrix Laurent residue, current amplitude, independent replay, or W/Z mass comparison \\
Electroweak hierarchy\slash naturality & selected-branch dimensionless theorem & local \(P\to\alpha_U\to v/E_\star\) hierarchy branch\newline \(\rightarrow\) exterior weak multiplicity \(3+1=4\)\newline \(\rightarrow\) mathematical capacity bridge and RG/Higgs square; conditional screen branch: 12 ports and an oriented 24-slot register & \(N_{\mathrm{EW}}(P_C)=3.5323546226929906511\ldots\times10^{122}\), \(N_0(P_{\mathrm{fwd}})=3.532131543418936\ldots\times10^{122}\), \(\mathcal B_{\mathrm{EW}}=0\), \(\epsilon_H=0\) & Supports the dimensionless hierarchy and naturality identities on the named branch. The 24-slot register does not produce four weak loads.\newline Equality with cosmic capacity follows from self-reference only after a source-derived direct public-record closure supplies the cosmic reading and a typed common-load bridge identifies it with the screen and electroweak readings. The finite-presence and Poisson factors require distinct extra premises and give retrospective comparison coordinates. An independent physical \(E_\star\), the public Thomson endpoint, theorem-level \(W/Z\), and the other listed mass rows are separate surfaces. \\
Higgs/top stage & double-criticality family from the gauge sector plus a target-anchored fit & criticality law \(\lambda=0,\ \beta_\lambda=0\) at one source scale \(\rightarrow\) fully constrained \((m_H,m_t)\) family \(\rightarrow\) declared boundary-scale candidate & no source-only mass; double-criticality branch \((m_H,m_t)=(125.77,\,172.63)\)~GeV, with \(m_H=125.72\)~GeV on the fit-free curve at the measured top; target-anchored fit \((125.1995304097,\allowbreak\,172.3523553288)\)~GeV & A physical mass requires the declared Gaussian/affine record and port-additive repair premises, the two equal-depth boundary-record anchors, a source root and physical scale, finite quotient-path selection, same-scheme running, target-independent rigidity, provenance, uncertainty enclosure, and a complex pole.\newline The companion top coordinate is not a separate public top-mass prediction. \\
Quark family & common-scale rejection plus restricted source non-identifiability theorem & one-scheme dimensionless Yukawas \(\rightarrow\) reciprocal-ray test\newline \(\rightarrow\) six-scalar generic interface and flavor-selector boundary & no public numeric quark row; at \(M_Z\), \(\rho_u\rho_d=0.835323\) and the endpoint-granted held-out error is \(21.56\%\) & The reciprocal-ray candidate is physically rejected across the tested running scales.\newline The \((\mathbb R_{>0})^2\) theorem is a restricted lower-bound obstruction.\newline Mixed-convention template and target-conditioned cumulant residuals are diagnostics.\newline A flavor-orbit selector, quark--Higgs carrier, and common-scale source transport are absent \\
Charged leptons & conditional construction plus exact target-anchored witness & shared excitation dictionary \(\rightarrow\) ordered charged carrier \(\rightarrow\) exact centered readback \(\rightarrow\) determinant-line lift on theorem-level physical charged data & exact centered readback; same-family target-anchored charged triple withheld from public prediction tables & public charged masses are not emitted from \(P\).\newline The theorem branch does not emit a theorem-level sector-isolated charged determinant exponent vector.\newline It does not attach a source-side determinant character to the physical charged determinant line.\newline The determinant-line lift and algebraic mass readout apply only on theorem-level physical charged data \\
Neutrinos & target-informed conditional candidate plus auxiliary checks & template family transport \(\rightarrow\) same-label scalar certificate \(\rightarrow\) weighted-cycle candidate \(\rightarrow\) compare-only bridge and absolute attachment \(\rightarrow\) shared-basis identity & declared scale-free candidate and PMNS/Majorana comparison rows; absolute mass attachments withheld from public prediction tables & exact linear algebra conditional on the declared template and selectors.\newline The candidate fails the NuFIT 6.1 correlated \((\sin^2\theta_{23},\delta_{\mathrm{CP}})\) profile, and neither the scale-free point nor the bridge invariant is a prediction or theorem~\cite{nufit61} \\
Hadrons & strong-binding construction absent; empirical comparison contract & source-derived hadronic spectral backend contract, including quotient ensemble, source QCD parameter map, Ward current accounting, spectral exports, and empirical \(e^+e^-\to\mathrm{hadrons}\) payload schema & no first-principles hadron masses; empirical comparison rows are separate from first-principles OPH results & First-principles hadron prediction requires a working OPH hadron construction, such as GLORB/Echosahedron, with the Ward-projected two-current spectral measure, higher-point/transition spectral sectors, same-scheme remainder, and systematics. The empirical comparison uses separate \(e^+e^-\to\mathrm{hadrons}\) input and cannot establish the first-principles theorem \\
Affine event-record stitching & conditional certificate theorem & event-manifold affine chart and observer clock atlas\newline \(\rightarrow\) affine event support for each descended token and clock slice\newline \(\rightarrow\) separately typed conditioned-frame ball where required\newline \(\rightarrow\) positive event-location gap\newline \(\rightarrow\) overlap descent\newline \(\rightarrow\) real transverse interface-crossing germs\newline \(\rightarrow\) common-chart sector/gauge transport\newline \(\rightarrow\) ID-independent one-to-one assignment gap\newline \(\rightarrow\) coarse/fine contraction check & nonbranching event-record paths, rays, or isolated records when the location and stitch certificates pass; otherwise \(\mathrm{AMBIGUOUS}\), \(\mathrm{REJECTED}\), or \(\mathrm{INTERACTION}\) \(\mathrm{REQUIRED}\) & This conditional theorem concerns observer-visible records on the event-manifold branch. It does not derive particle species, masses, gauge charges, scattering amplitudes, or geodesic motion. Repeated implementation IDs, nearest-neighbor fits, frame coordinates used as event positions, and file-boundary coincidences are inadmissible evidence. \\
\bottomrule
\end{tabular}
\end{table}

The theorem rows and target-anchored rows describe different objects. The theorem table carries a
two-modulus quark non-identifiability result and a
source-domain charged no-go boundary. The target-anchored rows contain
same-family witnesses and diagnostic calculations. The electromagnetic, color, and metric carrier
roles have conditional classical-mode receipts but no derived quantum mass rows; the \(W/Z\)
numbers are convention-dependent chart and adapter coordinates; and the
Higgs boson sits on the declared Higgs/top critical surface.

The following sections use the premises and caveats stated in the table.

\paragraph{Non-hadron results.}
The non-hadron results split by branch as follows:

\begin{table}[h]
\centering
\scriptsize
\setlength{\tabcolsep}{3pt}
\begin{tabular}{@{}>{\raggedright\arraybackslash}p{0.14\textwidth}>{\raggedright\arraybackslash}p{0.17\textwidth}>{\raggedright\arraybackslash}p{0.31\textwidth}>{\raggedright\arraybackslash}p{0.30\textwidth}@{}}
\toprule
Result class & Exact output(s) & Exact derivation & Conditions \\
\midrule
Carrier roles and classical modes & no photon/gluon/graviton particle-mass output & axioms \(\rightarrow\) realized electromagnetic/color/dynamical-metric roles; additional action/background/phase receipt \(\rightarrow\) transverse/TT quadratic modes & quantum-particle receipt not supplied; confined color has no free asymptotic-gluon claim \\
Electroweak chart and analytic consumer & no nonzero source-only mass output & source tuple \(\rightarrow\) exact two-coordinate chart; complete external renormalized packet \(\rightarrow\) strict scalar pole consumer and bounded analytic checks; finite quotient-path certificate \(\Rightarrow\) candidate value law; inverse adapter separate & the external fixture is not composed with the OPH chart. The receipts do not identify a physical sheet or full matrix amplitude, and no source-only pole pair or mass comparison follows \\
Hierarchy\slash naturality bridge & exact dimensionless bridge residual and zero Higgs naturality defect & axioms \(\rightarrow P\)-fixed source branch \(\rightarrow\) exact exterior weak multiplicity four \(\rightarrow\) mathematical \(N_{\mathrm{EW}}\) bridge map \(\rightarrow\) source-to-Higgs settled-form square; conditional screen branch: 12 ports and an oriented 24-slot register & selected-branch theorem for \(v/E_\star\) and the naturality identities; the register count does not construct the weak load; equality with cosmic capacity requires the direct public-record fixed point and the common screen/electroweak load-carrier map; it does not supply a physical \(E_\star\) or derive the mass rows \\
Higgs/top inverse slice & no nonzero source-only mass output & gauge core \(\rightarrow\) declared Jacobian \(\rightarrow\) exact inverse slice & target-conditioned benchmark only; the conditional coordinates inherit every open physical premise \\
Charged exact witness & exact same-family charged triple withheld & axioms \(\rightarrow\) shared excitation dictionary \(\rightarrow\) ordered charged carrier \(\rightarrow\) exact centered readback \(\rightarrow\) closed quadratic readout theorem \(\rightarrow\) same-family exact witness & same-family target-anchored comparison witness only; theorem branch carries exact centered readback plus the closed common-shift no-go. The theorem does not provide a \(\widehat C_e\) lift that emits the physical scalar \(\mu_{\mathrm{phys}}(Y_e)\) \\
Quark physical boundary & numeric predictions withheld; common-scale reciprocal-ray candidate rejected; mixed-convention formula residuals are target-anchored & six dimensionless Yukawas in one scheme and scale \(\rightarrow\) reciprocal-ray rejection \(\rightarrow\) six-scalar interface and selector no-go & at \(M_Z\), \(\rho_u=1.110889\), \(\rho_d=0.751941\), product \(0.835323\); even with four endpoints granted, the held-out miss is \(21.56\%\). The two-spread theorem is a restricted non-identifiability result \\
Neutrino candidate & declared scale-free comparison point; absolute mass attachments withheld & template family transport \(\rightarrow\) same-label scalar certificate \(\rightarrow\) target-informed weighted-cycle law \(\rightarrow\) compare-only bridge and absolute attachment \(\rightarrow\) shared-basis identity & rejected by the NuFIT 6.1 correlated profile; target dependence, source closure, basis orientation, and absolute normalization preclude a prediction~\cite{nufit61} \\
\bottomrule
\end{tabular}
\end{table}

The values and benchmark checks are
\[
\alpha^{-1}(0)=137.035999177(21),
\qquad
P\simeq1.6309682094,
\]
\[
\begin{aligned}
N_{\mathrm{EW}}(P_C)
&=
3.5323546226929906511\ldots\times10^{122},\\
N_0(P_{\mathrm{fwd}})
&=
3.532131543418936\ldots\times10^{122},
\qquad
\mathcal B_{\mathrm{EW}}=0,
\qquad
\epsilon_H=0.
\end{aligned}
\]
The source-forward bridge value is about \(6.6\) percent above the weighted
Planck base-\(\Lambda\)CDM capacity
\(N_\Lambda=3.31292709806038\ldots\times10^{122}\). The conditional
finite-presence and Poisson carriers give
\(3.292097877326465\ldots\times10^{122}\) and
\(3.300072225377652\ldots\times10^{122}\), respectively. Their comparisons
are retrospective. Exact countermodels obey the declared positive
composition and cut-count regrouping laws while producing different global
effects, so the finite source selects neither branch. Physical equality requires an executed public-record
fixed point plus the typed common screen/electroweak load-carrier map. No
source-only nonzero particle mass follows. The selected-carrier, value-law,
inverse-adapter, and conditional Higgs coordinates appear only in the technical comparison
sections below. Target-anchored charged-lepton
and quark witness values, including the mixed-convention quark target packet, appear only in
the exact-fit comparison calculations. Compare-only absolute neutrino attachments are withheld
from public prediction tables. The neutrino comparison surface keeps representative splitting
checks such as
\[
\begin{aligned}
\Delta m_{21}^2&=7.488059465106851\times10^{-5}\,\mathrm{eV}^2,\\
\Delta m_{31}^2&=2.5123118727618473\times10^{-3}\,\mathrm{eV}^2,\\
\Delta m_{32}^2&=2.4374312781107786\times10^{-3}\,\mathrm{eV}^2.
\end{aligned}
\]
The table states each derivation chain and the caveat that prevents a stronger theorem claim.

\paragraph{How to read mismatches and exact hits.}
Exact numerical agreement is not by itself a proof of a blind prediction. The provenance
record classifies the \(W/Z\) row as target-used frozen-reference reproduction, the charged-lepton
triple as an empirically anchored current-family witness, and the mixed-convention quark packet as a
target-anchored selected-frame comparison rather than a public source-only prediction. The
same record classifies the hierarchy/naturality formula as a zero-residual
identity on its declared selected branch, with \(\epsilon_H=0\); the strict
source-root and physical-scale certificates are separate premises. Where a
row does not match, or where it matches only under constrained premises, the explanations are
as follows. The source/root certificate gives
\(\alpha_{\mathrm{root}}^{-1}=136.994835177413\ldots\) and
\(P_{\mathrm{fwd}}=1.630972095858897\ldots\), but the declared map is
incomplete. The certified gauge-width map gives
\(\alpha^{-1}=137.035660136946577\ldots\), with gauge-width residual
\(2.5\times10^{-6}\) relative to the measured \(137.035999177(21)\). The no-hadron packet
that combines different pixel coordinates is no fixed point and is excluded from physical
output. Source-only closure requires Ward-projected hadronic spectral transport. The separate
endpoint-accounting diagnostic is
\[
\Delta^{\mathrm{fp}}_{\mathrm{H,cal}}
=\alpha_U(P_C)\,C_{24,Q},
\qquad
C_{24,Q}
=1.0009647859732326253849511140702475\ldots.
\]
The \(C_{24,Q}\) factor is endpoint accounting rather than a source-emitted theorem. At the public endpoint pixel value
\[
P\simeq1.6309682094
\]
the endpoint residual package requires
\[
\Delta_{\mathrm{source}}(P)=0.041465861005223389053448715357314044\ldots.
\]
That scalar belongs to the source-side hadronic spectral transport and
scheme-remainder map. The declared empirical route would use a separately labeled
\(e^+e^-\to\mathrm{hadrons}\) spectral input. No same-scheme integrated endpoint payload is emitted
here. The first-principles transport row excludes an inserted comparison endpoint. The
electroweak support record names the residual map and the
RG/matching/threshold/scheme packet.
An observer inside the branch measures the dressed Thomson coupling, not the
undressed source diagnostic \(1/136.994835\ldots\), because a zero-momentum
electromagnetic measurement sees the Ward-projected \(\U(1)_Q\) current after
charged-lepton vacuum polarization, confined-quark/hadron spectral transport,
and same-scheme finite endpoint matching have been included.
The electroweak \(W/Z\) row is therefore a chart and prescription comparison, not a physical mass
benchmark. Charged leptons carry a source-domain
no-go because the determinant trace-lift attachment from the electroweak
descendants of \(P\) to physical charged data is absent. The auxiliary
direct-top PDG row differs from the theorem coordinate by \(0.20673301656674425~\mathrm{GeV}\),
or \(0.28458848947515303\) combined standard deviations, because Q007TP and Q007TP4 are distinct
extraction codomains; the direct-top response kernel has a source-domain no-go boundary. Neutrino absolute masses are
not directly measured, and PMNS-angle residuals are visible comparison tension outside the
theorem branch. First-principles hadron masses have no emitted prediction because the
required Ward-projected hadronic spectral measure must come from a working OPH
hadron construction. Empirical hadron closure rows carry a separate support class.

\paragraph{Local unification surface.}
The bosonic rows carry one cross-branch statement. The same pixel input \(P\), fixed on the
synthesis-paper outer/inner closure relation, fixes the
electroweak and criticality trunk
\[
P
\longmapsto
\alpha_U(P)
\longmapsto
\bigl(t_U(P),t_{\mathrm{tr}}(P)\bigr)
\longmapsto
v_{\mathrm{chart}}(P)
\longmapsto
\left(M_W,M_Z\right).
\]
\[
\sigma_{\mathrm{crit}}
=
\alpha_U(P)\cos(2\theta_{W0})/\sqrt{\pi}
\longmapsto
\left(m_H,m_t\right).
\]
Here \(\alpha_U(P)\) is the branch value selected by the same forward
pixel-closure solve, so
the bosonic trunk contains no inverse electroweak readback of the internal transmutation data.
On the companion gravity side, the same pixel law packages
\[
\ellbar_{\mathrm{SU(2)}}(t_{2,\mathrm{run}})
\;+\;
\ellbar_{\mathrm{SU(3)}}(t_{3,\mathrm{run}})
=
P/4.
\]
The gravity normalization itself is emitted by the selected scale certificate
\(\gamma_\star=\ell_\star\nu_{\mathrm{Cs}}/c\), equivalently
\(B_\star=3\pi/\ell_\star^2\). With
\(a_{\mathrm{cell}}=P\ell_\star^2\), the Newton area law gives
\(G_{\mathrm{geom}}=\ell_\star^2\), so the pixel cancels.
The local unification surface separates one explicit familiar-unit readout package from three
mathematical surfaces: the \(W/Z\) comparison, the conditional
declared-surface Higgs/top coordinate,
and the gravity-side readout with its strict classical-regime clause.
On the gravity side, the stated local extension surface uses the lifted product presentation of the
realized quotient branch and identifies
\[
\ellbar_{\mathrm{shared}}
=
\ellbar_{\mathrm{SU(2)}}(t_{2,\mathrm{run}})
\;+\;
\ellbar_{\mathrm{SU(3)}}(t_{3,\mathrm{run}}).
\]
On that same surface the pixel law fixes \(\ellbar_{\mathrm{shared}}=P/4\),
and the local SI readout is
\[
G_{\mathrm{SI}}=\frac{c^3\ell_\star^2}{\hbar}
\]
from the selected scale certificate. The \(\chi_\nu\) protected-reserve branch
uses the same public endpoint convention as
\[
P_\chi=P_{\mathrm C}=1.630968209403959\ldots,
\qquad
\epsilon_{\rm res}=\frac{P_\chi}{24},
\]
after the same-collar shared-edge budget and one-class \(\mathbb Z_6\) trace
receipts pass. This does not make \(P/4\) a primitive Hilbert-space dimension:
\(P/4\) is the local screen-cell entropy budget in the particle and hierarchy
branches, while \(P_\chi/24\) is the shared scalar-reserve density used by the
\(\chi_\nu\) collar branch. On that declared extension surface the same
familiar-unit package reads
\[
L_{\mathrm{loc}}=\sqrt{a_{\mathrm{cell}}}\,\widehat L(P),
\qquad
t_{\mathrm{loc}}=\frac{\sqrt{a_{\mathrm{cell}}}}{c}\,\widehat T(P),
\]
\[
E_{\mathrm{loc}}=\frac{\hbar c}{\sqrt{a_{\mathrm{cell}}}}\,\widehat E(P),
\qquad
\Theta_{\mathrm{loc}}=\frac{\hbar c}{k_B\sqrt{a_{\mathrm{cell}}}}\,\widehat\Theta(P),
\]
with dimensionless \(\widehat L,\widehat T,\widehat E,\widehat\Theta\). Thus, at fixed \(P\), the
local ruler is \(\sqrt{a_{\mathrm{cell}}}\); seconds are that ruler divided by the structural
Lorentz output \(c\); and GeV and Kelvin are downstream familiar-unit displays of the inverse
local ruler through \(\hbar\) and \(k_B\). On that declared extension
surface the local scale-readout bundle is
\[
\begin{aligned}
c&=299792458\,\mathrm{m/s},\\
G&=6.674299995910528\times10^{-11}\,\mathrm{m^3\,kg^{-1}\,s^{-2}},
\end{aligned}
\]
No nonzero particle-mass coordinate is a source-only prediction. The \(W/Z\) inverse adapter
and the conditional Higgs/top coordinates are compare-only quantities discussed in their technical
sections. The structural Lorentz output is the common invariant null cone and speed
\(c_\star\). The decimal \(299792458\,\mathrm{m/s}\) is exact by the SI definition of the metre,
not a predicted magnitude. The paper keeps the
\(G\) row as an exact scale-readback value, with no literal zero-difference identity against the
rounded benchmark \(6.6743\times10^{-11}\).

The detailed inventory below records the source-derived theorems, conditional
continuations, and target-anchored comparisons sector by sector.

\subsection{Detailed particle inventory}

\tiny
\setlength{\tabcolsep}{2pt}
\begin{longtable}{@{}>{\raggedright\arraybackslash}p{0.135\textwidth}>{\raggedright\arraybackslash}p{0.145\textwidth}>{\raggedright\arraybackslash}p{0.135\textwidth}>{\raggedright\arraybackslash}p{0.215\textwidth}>{\raggedright\arraybackslash}p{0.265\textwidth}@{}}
\toprule
Family & Particle & Theorem scope & OPH value or strongest derived benchmark & Premise boundary \\
\midrule
\endfirsthead
\toprule
Family & Particle & Theorem scope & OPH value or strongest derived benchmark & Premise boundary \\
\midrule
\endhead
\bottomrule
\endfoot
Conditional carrier modes & electromagnetic & action/phase receipt & two transverse classical \(k^2=0\) modes; \(\mu_{Q,\mathrm{quad}}^2=0\) only in the displayed Maxwell action & physical Hilbert space, positive-residue pole, and stable asymptotic photon receipt absent \\
Conditional carrier modes & color & perturbative/pre-confinement action receipt & \(2\dim G\) transverse classical \(k^2=0\) modes; \(\mu_{\mathrm{YM},\mathrm{quad}}^2=0\) only in the displayed pure-Yang--Mills expansion & positive-residue pole and deconfined asymptotic colored sector absent; no free gluon is claimed on the confining branch \\
Conditional carrier modes & tensor & action/background receipt & two Einstein transverse-traceless classical \(k^2=0\) modes; \(\mu_{\mathrm{EH},\mathrm{TT}}^2=0\) only on the pure-Einstein branch & metric quantization, positive physical Hilbert space, positive-residue pole, and asymptotic/EFT particle interpretation absent \\
Electroweak branch & \(W\) boson & no source-only physical mass & strict scalar consumer; principal-sheet zero exclusion and one simple channel-specific lower-half algebraic-chart zero for each of \(W\) and \(Z\) & neither chart is identified with the physical resonance sheet; no unique continuation, sign bridge, full-matrix Laurent residue, current amplitude, independent replay, or composition with the OPH chart is certified \\
Electroweak branch & \(Z\) boson & no source-only physical mass & same strict consumer and bounded scalar analytic checks & same boundary as the \(W\) row \\
Higgs/top critical stage & Higgs boson & no source-only physical mass & no OPH-native pole & requires an independent source root and physical scale, finite quotient-path selection, same-scheme running and threshold transport, target-independent criticality rigidity, a complex pole, and an uncertainty enclosure \\
Quark family & top quark & separate target-anchored comparison coordinate & withheld from public prediction table & cross-section pole-mass extraction coordinate; it is not a sixth entry in one common running-mass chart and is not selected by the source spread laws \\
Quark family & up quark & reciprocal-ray candidate rejected & withheld from public prediction table & common-scale dimensionless Yukawa test fails; a source-derived flavor-orbit selector is absent \\
Quark family & down quark & reciprocal-ray candidate rejected & withheld from public prediction table & same common-scale rejection and six-scalar interface boundary \\
Quark family & strange quark & held-out reciprocal-ray failure & withheld from public prediction table & part of the \(19.9\%\) to \(21.6\%\) held-out failure across tested scales \\
Quark family & charm quark & held-out reciprocal-ray failure & withheld from public prediction table & stored GeV-valued matrices are mixed-convention mass textures, not physical Yukawas \\
Quark family & bottom quark & reciprocal-ray candidate rejected & withheld from public prediction table & common-scale source transport, Higgs normalization, and flavor-orbit selection absent \\
Charged leptons & electron & continuation gap & \(n/a\); exact same-carrier centered readback exists once a charged source pair is emitted, but the absolute scale is blocked by the closed common-shift no-go; benchmark target \(g_e^\star=0.0457789\), equivalently \(\Delta_e^{\mathrm{abs},\star}=3.00398633\) & close branch-generator splitting, then emit the lift whose descended scalar is \(\mu_{\mathrm{phys}}(Y_e)\); within that lift the exact smaller forcing object is the physical identity-mode equalizer, after which \(\widetilde C_e(Y_e)=\widehat C_e(Y_e)+\mu_{\mathrm{phys}}(Y_e)\,\mathbf 1\), \(A_{\mathrm{ch}}(Y_e)=\mu_{\mathrm{phys}}(Y_e)\), and the readouts \(g_e\), \(\Delta_e^{\mathrm{abs}}\) follow \\
Charged leptons & muon & continuation gap & \(n/a\); same exact centered-readback / common-shift-no-go frontier as the electron row & same \(\widehat C_e^{\mathrm{cand}} \rightarrow\) branch-generator splitting \(\rightarrow\) lift \(\rightarrow \mu_{\mathrm{phys}}(Y_e) \rightarrow A_{\mathrm{ch}} \rightarrow g_e\) closure chain, with the physical identity-mode equalizer beneath the descended scalar \\
Charged leptons & tau lepton & continuation gap & \(n/a\); same exact centered-readback / common-shift-no-go frontier as the electron row & same \(\widehat C_e^{\mathrm{cand}} \rightarrow\) branch-generator splitting \(\rightarrow \mu_{\mathrm{phys}}(Y_e) \rightarrow A_{\mathrm{ch}} \rightarrow g_e\) closure chain, with the physical identity-mode equalizer beneath the descended scalar \\
Neutrino comparison & declared \(f\)-basis weighted-cycle matrix & rejected target-informed candidate & no flavor-particle mass row; the frozen declared-basis coordinate gives \(\theta_{12}=34.2259^\circ\), \(\theta_{23}=49.7228^\circ\), \(\theta_{13}=8.68636^\circ\), \(\delta=305.581^\circ\), \(J=-0.02753\), and, under its declared normal-ordering labels, \(\Delta m_{21}^2/\Delta m_{32}^2=0.03072111\) & the upstream flavor kernel is a hand-written template; the physical charged basis and mass-label rule are premises the template does not supply; the target-informed selector precludes a blind-prediction claim; the NuFIT 6.1 correlated profile rejects the candidate; Majorana and absolute-attachment values are compare-only \\
Hadrons & proton & strong-binding construction absent; empirical comparison contract emitted & no first-principles emitted prediction & first-principles prediction requires a working source-derived hadronic backend, such as GLORB/Echosahedron, plus Ward-projected two-current, higher-point, and transition spectral exports, same-scheme remainder, and production systematics; empirical closure rows use separate \(e^+e^-\to\mathrm{hadrons}\) input \\
Hadrons & neutron & strong-binding construction absent; empirical comparison contract emitted & no first-principles emitted prediction & same OPH construction premise, with isospin-resolved hadron systematics in the construction support class \\
Hadrons & neutral pion proxy & strong-binding construction absent; empirical comparison contract emitted & no first-principles emitted prediction & same OPH construction premise; local stable-channel surrogate output is not a paper prediction \\
Hadrons & \(\rho(770)^0\) proxy & strong-binding construction absent; empirical comparison contract emitted & no first-principles emitted prediction & same OPH construction premise, plus finite-volume resonance extraction in the construction support class \\
\end{longtable}
\normalsize

\section{Source Values, Screen Architecture, and Theorem Premises}\label{sec:particle-input-ledger}

The particle derivation uses the three axioms, the local pixel closure, the
regulated screen realization, and the imported OPH theorems stated below.
The conditional screen-capacity branch is used only in results that state it
as a premise.

\subsection{Canonical OPH basis}

The particle-spectrum derivation uses the three OPH axioms shared with the
spacetime and gauge derivation papers~\cite{ophEinstein,ophGauge}. They are
restated here for local readability because every particle-family chapter
depends on them either directly or through the structural OPH chain imported
from those papers.

% Shared three-axiom basis panel. Included by papers that state the full
% axiom basis. Canonical machine source: claims/axiom_registry.yaml;
% reviewed reference: docs/AXIOM_REFERENCE.md. Do not hand-copy a competing
% formal block into a paper; include this fragment.

\paragraph{Axiom 1 (oriented observer-patch federation and spherical support).}
There exists an observer patch net on an oriented spherical screen: at every
finite resolution each local carrier has twelve primitive boundary ports
forming the vertices of an oriented triangular boundary with 30 edges and 20
faces, combinatorially the boundary of an icosahedron, and carriers join
through typed seams and coherent triple overlaps, refine to an oriented
spherical support, and expose local state, readback, records, repair moves,
and checkpoints. Their complete infinitesimal port response is represented
faithfully on a finite-dimensional unitary response space and is closed under
the commutators generated by ordered response composition. Formally, for
every regulator \(r\) in a
directed system there is a typed object
\(\mathfrak N_r=(\mathcal P_r,\mathcal A_r,\mathcal R_r,\mathcal I_r,
\mathcal U_r,\mathcal C_r,N_r,S_r,b_r)\):
a finite patch/overlap category with an isotone local algebra net and
central record algebras; a federation of finite carriers, each with twelve
primitive pairwise-orthogonal central port projections summing to one and a
boundary packet \(K=(P,E,F,o)\) with \(|P|=12\), \(|E|=30\), \(|F|=20\),
degree-five ports, five-cycle links, and coherently oriented two-face edges,
isomorphic to the icosahedral boundary complex with no preferred labels;
seam algebras with unital restrictions and coherent triple-overlap cocycles
forming the nerve \(N_r\); a finite edge-midpoint refinement support
\(S_r\) realized orientation-preservingly in \(S^2\) with mesh tending to
zero; and a source-bound degree-one bridge \(b_r:N_r\to S_r\), all commuting
with refinement. Each carrier also has
\(V_{r,i}=\mathbb R[P_{r,i}]\), a finite-dimensional complex Hilbert space
\(H_{r,i}\), and an injective real-linear map
\[
D_{r,i}:V_{r,i}\longrightarrow\mathfrak u(H_{r,i}).
\]
The image \(\mathfrak g_{r,i}=D_{r,i}(V_{r,i})\) is closed under
commutators and complete for the declared quotient-visible infinitesimal
port response. The primitive port probes span \(V_{r,i}\), no additional
public response direction is omitted, and
\(-\operatorname{Tr}(D_{r,i}(v)D_{r,i}(w))\) is positive definite. The
response construction is natural under admissible presentation equivalence
and refinement. The axiom does not select a particular response table,
inverse-port law, Lie type, global group, particle interpretation, or metric
content.

\paragraph{Axiom 2 (observer agreement).}
Observers operating on the screen agree on the meaning of the data they
jointly interpret. Formally, the interpretation map
\(\mathcal J_r:\mathsf{Data}_r\to\mathsf{Meaning}_r\) from
observer-accessible data to operational meanings is natural with respect to
every visible overlap restriction, recharting, seam translation,
higher-overlap map, federation map, and refinement map: every declared
data-access diagram for accepted public data commutes after interpretation,
and \(\mathcal J_r\circ c_{s\to r}=C_{s\to r}\circ\mathcal J_s\) across
resolutions. The axiom constrains accepted shared data only, and the domain
it quantifies over is supplied by the Axiom 1 interfaces rather than by the
axiom itself.

Accepted reversible overlap transports form a groupoid \(\mathcal O_r\).
For a complete carrier chart \(o\), its closed paths
\(\operatorname{Hol}_r(o)\) map onto the orientation-preserving incidence
automorphisms \(\operatorname{Aut}^{+}(K_{r,i})\). For every proper carrier
automorphism \(a\), Axiom 2 supplies a closed path \(\gamma_a\) and one
projective implementer
\([U_a]\in\operatorname{PU}(H_{r,i})\) satisfying
\[
\operatorname{Ad}_{[U_a]}D_{r,i}(v)=D_{r,i}(a\cdot v).
\]
Writing
\[
G_{D,r,i}^{0}
=
\left\langle
\exp(tD_{r,i}(v)):t\in\mathbb R,\ v\in V_{r,i}
\right\rangle^{0},
\]
the implementer is endogenous:
\[
[U_a]=[g_a c_a],
\qquad
g_a\in G_{D,r,i}^{0},
\quad
c_a\in C_{U(H_{r,i})}(\mathfrak g_{r,i}).
\]
Closed paths compose projectively and the construction commutes with
refinement. A unitary lift need not be unique. Scalar phases and centralizer
actions are permitted because they act trivially on the response algebra.
An independent spectator carrying a matching projective action cannot
replace the response-generated factor or contribute a public response
direction. The closed path, port action, and implementer must arise from the
same accepted source history in an executable realization.

The two response clauses support an abstract target-free compact Lie-type
theorem. They do not certify that an executable producer emitted a particular
current fixture. Such a certification must reconstruct the port generators,
commutator, and projective implementers from ordered source histories without
a named current model upstream. Axiom 2 does not imply global state
extension, repair termination, confluence, unique normal forms, record
durability, Byzantine safety, dissemination, or commutation of state
optimization with refinement; raw mismatch and repair precede acceptance.

\paragraph{Axiom 3 (conditional maximum randomness).}
Everything that observer agreement leaves unconstrained is maximally random.
Formally, at finite regulator \(r\) a state is a compatible family of local
normalized states on the accessible algebra net; \(\mathcal K_r\) is the
nonempty convex set of such families satisfying the finite observer-visible
constraints supplied by Axioms 1 and 2 through an A1-generated constraint
grammar with a factorization theorem; and, relative to an exact reference
family \(\tau_r\), an A1-generated observer cover \(\mathcal G_r\), and
strictly positive exact weights \(w_{r,P}\) constructed from quotient-visible
A1 data by a rule natural under admissible presentation equivalence, the
realized state is the
information projection. The cover restriction map is injective on
\(\mathcal K_r\); equivalently, at smooth points no nonzero feasible
tangent is invisible to the whole cover. The realized state is
\[
\rho_r=\operatorname*{arg\,min}_{\rho\in\mathcal K_r}
\sum_{P\in\mathcal G_r} w_{r,P}\,D(\rho_{r,P}\Vert\tau_{r,P}).
\]
The same projection rule applies separately to each declared optimizer
object type (ontic state, inference state, or transition distribution), and
statements about one type do not transfer to another.
This is equivalent to weighted local entropy maximization when every local
reference density is identity-proportional in its declared trace. "Random"
means least informative relative to the
declared reference and the complete agreement constraint set. The axiom
selects one state inside one fixed feasible space. It does not compare
unrelated state spaces or field contents, does not imply optimizer
compatibility across refinement, and does not imply recovery, alignment, or
mixing; each such statement carries its own theorem, interface, or
countermodel.

The complete formal basis, including its constraint map and non-implications,
is given in the
\href{https://github.com/FloatingPragma/observer-patch-holography/blob/main/docs/AXIOM_REFERENCE.md}
{canonical three-axiom reference}.


% Shared closure-principle panel. Included by papers that state the full
% axiom basis, immediately after THREE_AXIOM_BASIS.tex. This states a
% global principle on the axiom class, not a fourth axiom; every surface
% counting the basis counts exactly three axioms. Reviewed reference:
% docs/AXIOM_REFERENCE.md.

\paragraph{The closure principle.}
The axioms describe the observer screen. The global closure principle states
that the simulating and the simulated description are one system. Every quantity that
has both a construction-side reading and a readback-side reading must take
the same value once a typed bridge proves that the two readings denote one
invariant. The equality is forced by self-identity. Constructing that bridge
and the return map remains part of the physics. Existence, uniqueness, and
stability are separate determinacy tests on the resulting equation. The name
given to such an equation carries no mathematical weight: after the typed
identification, unequal readings would describe two systems rather than the
single self-referential universe.

The local screen-grain proposal
\[
P=\varphi+\sqrt{\pi}/A_T(P),
\]
with \(A_T(P)\) the inverse coupling returned by the declared source map,
has one interval-certified fixed point on its physical domain; its
laboratory attachment awaits a same-scheme hadronic spectral transport.
The displayed map's certified fixed point sits at $P=1.6309720959$ and
returns $A_T=136.9948352$; the certified gauge-width variant's fixed
point sits at $P=1.6309682414$ and returns $137.0356601$. Through the
same outer equation the measured 2022 CODATA coupling
$\alpha^{-1}=137.035999177(21)$ reads back the grain value
$P=1.6309682094$. The comparisons stand at $3.0\times10^{-4}$ relative
for the displayed map and at $2.5\times10^{-6}$ for the gauge-width
variant, as diagnostics of the declared source maps.

The direct global proposal is
\[
N=\log M_0(\mathfrak U_N).
\]
The complete declared finite source class does not select a unique cosmic value, so
any successful completion needs an additional source law. A separate
common-load branch starts from
\[
N_0=\pi\exp\!\left(\frac{6\pi}{P\,\alpha_U(P)}\right).
\]
On the finite collar branch, the declared total reserve expectation \(P/4\)
and six-class equidistribution give presence probability \(P/24\) for each
declared class. If one class is physically selected as the blocked event, its
scalar-weighted presence receipt is discharged, and that normalized
collar-survival factor is proved to act on the global capacity, the
corresponding candidate is
\[
N_{\mathrm{pres}}=N_0\left(1-\frac{P}{24}\right),
\qquad
\ln\frac{N_{\mathrm{pres}}}{\pi}
=
\frac{6\pi}{P\,\alpha_U(P)}
+\ln\left(1-\frac{P}{24}\right).
\]
The exponential alternative \(N_{\mathrm{Pois}}=N_0e^{-P/24}\) requires a
separate mean-count or continuum carrier. Neither global attachment is
derived. The inherited screen/electroweak bridge premises, physical
common-load bridge, physical \(\mathbb Z_6\) seam action, one-class reserve
selection, scalar-weighted receipt, reserve-to-capacity map, and
horizon-record identification are open. On the source-forward numerical
branch,
\[
N_0=3.5321\times10^{122},
\]
so the finite-presence candidate is
\(3.2921\times10^{122}\) and the exponential candidate is
\(3.3001\times10^{122}\). The full weighted Planck
base-\(\Lambda\)CDM chain gives the comparison coordinate
\(3.3129\times10^{122}\). The respective residuals are
\(-0.63\) percent and \(-0.39\) percent. Both comparisons were exposed before
the branch choice and carry no predictive weight.


The first axiom's spherical support is The Standard Model gauge paper's
screen-first global-support branch. On the producer branch, \(S^2\) becomes
available only after spherical-incidence, mesh, cross-ratio, normalization,
refinement, and carrier-to-support receipts establish it. The resulting
support screen is distinct from both the local twelve-port carrier boundary
and the federation overlap nerve.

Two families of premise interfaces enter beside the axioms, each named at
the results that consume it. The gravity-adjacent chapters consume the
edge-center, collar-recovery, generalized-entropy, and modular-normalization
interfaces of the spacetime paper; these are classified derivation targets
and physical identifications with their own countermodels, not axioms. The
particle chapters consume declared sector completions: the generation count
inside the conditional window, scalar attachment and multiplicity, and the
completeness statement for extra light sectors are physical attachment
premises, and the charged-lepton and quantitative selector branches state
their declared premises where they appear. The transportable-sector/Tannaka
reconstruction is a separate conditional classification route.

\subsection{Regulated screen architecture and patch language}

The SM/GR derivation in Ref.~\cite{ophEinstein,ophGauge}
states the axioms in algebraic language. Ref.~\cite{ophmicrophysics} supplies one concrete
regulated realization of that language: a federation of finite patch carriers with echosahedral
multi-port interfaces, recurrent toroidal subchannels, exposed overlap packets, record algebras,
and local repair instruments. The formal observer patch is the bounded access-and-record structure.
One echosahedral carrier is a candidate primitive physical realization on the homogeneous branch;
the carrier alone does not satisfy the observer definition. This regulated architecture is
important for the particle derivation for three reasons.

First, it makes the screen picture operational. A patch is a finite local
algebra in a bounded carrier, and an overlap is a boundary-visible algebra with declared shared
observables. Second, it makes the edge-sector
language concrete. The particle derivation uses edge sectors, transport, refinement,
fusion, and overlap data to build the gauge and flavor branches; Ref.~\cite{ophmicrophysics} shows
how those objects can be realized at fixed cutoff by finite gauge-aware patch carriers.
Third, it clarifies the measurement interface. The observer-facing measurement package is a
theorem-bearing fixed-cutoff statement about a central record algebra embedded
in a declared finite algebra-state representation, Born probabilities for its
event projectors, and L\"uders conditioning on that representation. The
quantum representation is not derived from the commuting record indicators.

Three geometric objects govern this reference implementation. The local carrier boundary has the
twelve-port icosahedral incidence certified on the echosahedral lineage. The federation screen is the
finite federation together with its overlap nerve. The support screen is the observer-visible
\(S^2\) chart obtained only when global incidence, mesh, cross-ratio, normalization, and refinement
receipts hold. A federation of local icosahedral carriers can have a nonspherical nerve, so local
\(A_5\) symmetry does not force a global \(S^2\) support.

The framework is presentation-invariant and carrier-sensitive. A different hidden implementation
is physically equivalent when its observer-visible quotient agrees. A different visible port
graph, topology, record process, or response law may select a different branch. Hardware evidence
enters only through a public OPH evidence bundle with stable hashes and verifier receipts.

Physical phase locking is a candidate producer for coherent overlap comparison. To support that
interpretation it must produce the accepted repair relation, transactional confluence, public
records, and controlled noise bounds. No current theorem identifies phase locking with consensus
confluence, modular flow, or an observer clock. Separate support-visible
continuum theorems supply the geometric and gauge limits.

\subsection{Quantitative inputs and where they enter the particle derivation}

The quantitative derivation developed here uses one incomplete declared-map coordinate and one
conditional capacity coordinate together with the
selected no-\(G\) scale certificate:
\begin{align}
P &\equiv a_{\mathrm{cell}}/\ell_\star^2,\\
D &\equiv\dim\mathcal H_{{\rm cap},r,D},
\qquad N_{\mathrm{CRC}}\equiv\log D_{\mathrm{CRC}},\\
\gamma_\star&\equiv\frac{\ell_\star\nu_{\mathrm{Cs}}}{c},
\qquad
B_\star\equiv\frac{3\pi}{\ell_\star^2}.
\end{align}

These closure values have distinct downstream dependencies.

\paragraph{Pixel area \(P\).}
The pixel area is the declared local UV-area ratio. In this derivation it feeds the heat-kernel /
edge-law input stage and, through that route, the forward electroweak surface. It is the
common upstream numerical variable for the electroweak, Higgs/top, flavor, quark, and
hadron-facing branches. The particle paper therefore treats \(P=P_\star\) as the inherited
root of the incomplete outer/inner declared map
\[
P_\star=\varphi+\frac{\sqrt{\pi}}{A_T(P_\star)}
\]
from the synthesis paper~\cite{ophsynthesis}, not as a quantity derived again inside the
particle-spectrum argument itself.

\paragraph{Conditional screen capacity \(N_{\mathrm{CRC}}\).}
The screen capacity is not part of the local recovered-core gravity/gauge derivation. It enters
the separate screen-capacity branch through the stable cosmic record-closure target
\[
\mathfrak F_{r,0}(D_\star)=\{D_\star\},
\qquad N_{\mathrm{CRC}}=\log D_\star,
\qquad
\Lambda_{\mathrm{CRC}}\ell_\star^2=\frac{3\pi}{N_{\mathrm{CRC}}},
\qquad
\Lambda_{\mathrm{CRC}}=\frac{3\pi}{G N_{\mathrm{CRC}}}.
\]
The synthesis paper specifies the physical active-readback target directly.
At finite cutoff, compatible reachable public record atoms and the globally
coupled checkpoint family define a confusability graph \(G_q\) and
\[
M_0(q)=\alpha(G_q),
\qquad
\mathfrak F_{r,\varepsilon}(D)=
\{M_\varepsilon(q):q\in\widetilde\Omega_{r,D}\}.
\]
Only a whole-fiber singleton is scalar. A faithful capacity-carrier
representation makes that scalar exact map deflationary; a
confusability-reflecting capacity embedding makes it monotone, and fixed-\(D\) refinement eventually
stabilizes. A source-derived fixed-cutoff public checkpoint packet,
carrier representation, and whole-fiber scalarization exist at \(D=24\).
A bounded target-clean all-rung counterfamily has different exact zero sets
under shared base agreement, positivity, the carrier bound, and executable
finite controls. The exact result is nonidentifiability on that bounded
completion class. Universal all-rung membership in the complete A1--A3
capacity-source contract and the executable-to-Lean bridge are open. Direct
\(N\) is not evaluable on its incomplete source antecedent. A positive direct
result requires completion of that antecedent, a unique physical zero, and a
universe-level carrier attachment.
Identifying that capacity with de Sitter entropy requires the
horizon--record identification; identifying it with the electroweak bridge
requires the common screen/electroweak load-carrier identification. The
source-forward bridge value
\(3.5321315434\ldots\times10^{122}\) is about \(6.6\) percent above the
weighted Planck base-\(\Lambda\)CDM comparison coordinate
\(3.3129270981\ldots\times10^{122}\).
On the declared one-class branch, physically selecting that class and
attaching its scalar-weighted finite presence-survival factor gives
\(3.2920978773\ldots\times10^{122}\); the exponential mean-count branch gives
\(3.3000722254\ldots\times10^{122}\). Exact neutral and multiplicative
completions of the same local datum satisfy the declared positive composition
and regrouping laws, then disagree globally. The finite source selects no
action or blocked-event semantics, so the named-law branch is not evaluable
and its horizon branch has no capacity object. Both comparisons are
retrospective. The last equality is therefore a conditional scale-certified
display; it does not determine an SI curvature scale from \(P_\star\) and
\(N_{\mathrm{CRC}}\) alone.
An independently frozen discrimination scale \(\rho_{\rm op}\) is a test of
the direct code capacity through
\(\log M_0-\pi/\rho_{\rm op}^2\), not its definition. A diagonal
terminal-state count or its argmax does not construct the off-diagonal map.
In the particle context, this matters only for the cosmological-capacity discussion that
frames why local null data do not determine the
cosmological constant. The implemented branch uses the static-patch normalization
\[
N_{\mathrm{patch}}=\left(\frac{r_{\mathrm{dS}}}{\ell_P}\right)^2\simeq1.05\times10^{122},
\qquad
N_{\mathrm{scr}}=\pi N_{\mathrm{patch}}\simeq3.313\times10^{122}.
\]
This capacity bookkeeping supplies no neutrino result; the weighted-cycle neutrino derivation is
a separate branch.

\paragraph{Selected scale certificate.}
The gravity scale is a separate observation-located certificate. On the declared branch,
\[
\gamma_\star=\frac{\ell_\star\nu_{\mathrm{Cs}}}{c},
\qquad
B_\star=\frac{3\pi}{\ell_\star^2},
\qquad
G_{\mathrm{SI}}=\frac{c^3\ell_\star^2}{\hbar}.
\]
If the declared-map root acquires its missing physical identification, the local pixel
\(P_\star\) identifies \(a_{\mathrm{cell}}=P_\star\ell_\star^2\) and
\(\ellbar_{\mathrm{shared}}=P_\star/4\). In the Newton area law \(P_\star\) cancels, so this paper
does not treat \(G\) as a particle-sector consequence of \(P_\star\).
\(\ellbar_{\mathrm{shared}}\) is a shared-cut density, not the logarithm of an autonomous cell
Hilbert factor. The particle branch uses it for cell/edge consistency; it does not select a fixed
primitive observer capacity.

\paragraph{Why these are the declared closure/readback values.}
One of the discipline constraints of the OPH derivation is that particle-sector freedom should not
be hidden in a large uncontrolled parameter set. This paper exposes the incomplete declared-map
root \(P_\star\), the proposed \(N_{\mathrm{CRC}}\) correctable-public-record closure
coordinate, and the
selected scale certificate. The same shared-coordinate discipline applies to
every conditional branch. For the particle-spectrum branch, the nontrivial
common quantitative burden sits on \(P_\star\): \(N_{\mathrm{scr}}\) enters the separate
capacity/cosmology side and \(\gamma_\star\), equivalently \(B_\star\), enters the separate
gravity-scale side, whereas the
reverse-engineering claim for masses and couplings is that one shared declared-map coordinate should
drive the spectrum instead of a long sector-by-sector input list.

\subsection{Technical theorem data carried into the particle-spectrum derivation}

This paper uses the same theorem premises as
\emph{Recovering Observer Spacetime and Einstein Dynamics from Overlap Consistency and Deriving Standard Model Gauge Structure from Observer Overlap Consistency}~\cite{ophEinstein,ophGauge}, together with
the regulator-language clarification integrated into the consensus and microphysics
appendices. This separation matters only
insofar as it separates branch-internal statements from declared inputs.

On the finite identity-reference branch, the third axiom and the declared
local constraint grammar give the regulator-scale local Gibbs form, the
quasi-local propagation bound, and the endpoint-control estimate for bounded
intervals. Refinement stability is supplied by the separate optimizer-pushforward
and refinement-tail interface. In regulator language, one works with finite
local Hilbert spaces and a boundary-fixed overlap action on cut data. Beyond
those regulator premises, the Standard Model gauge paper supplies the support-visible
Bisognano--Wichmann (BW)
scaling theorem for the Lorentz/null-modular/Einstein branch. Transportability and the fixed-cutoff bosonic sector category are theorem-produced in The Standard Model gauge paper; the refinement/fiber
ladder and cofinal admissible compact-gauge witness additionally require its explicit compact-gauge refinement receipt. On the central branch the matching gluing
obstruction is the combined zero-obstruction transport criterion: vanishing triangle defect plus
trivial represented holonomy of the strictified edge \(1\)-cocycle. It introduces no second theorem-side input. Any
local-Gibbs or collar-mixing language on that branch is fixed-cutoff recoverability/support
control; it is not the source of the compact-gauge witness.
For the BW/geometric side, the target is the support-visible extracted geometric subnet.
The Standard Model gauge paper proves the needed scaling theorem by combining
regularized support-visible modular transport, weak-\(*\) extraction of the
cap pair in the Gelfand--Naimark--Segal representation, support-readable modular covariance,
BW framing, held-out oriented cross-ratio rigidity, and independently normalized geometric
\(2\pi\)-Kubo--Martin--Schwinger (KMS) convergence with wrong-scale controls. The unregularized full-algebra common-floor
route is deliberately not claimed, because off-support directions may collapse without affecting
observer-facing matrix elements. Bare finite consensus does not by itself produce the finite
cap-normal BW certificate used by that theorem. The theorem also consumes an
independent mixed Gelfand--Naimark--Segal common-comparison premise on the
same tower. Local maximum entropy and refinement do not produce that premise.

For the particle-spectrum derivation, this has a practical consequence.
Structural statements such as compact-gauge reconstruction, the
contract-conditional Standard Model quotient, the exact hypercharge lattice,
the color count \(N_c=3\), and the declared generation completion are not floating
independently of the declared input and theorem premises. Nor do they imply a physical family
attachment. The additional classical carrier modes
live on explicit quadratic action/background/phase receipts, and particle poles require the
stronger quantum receipt. These statements live on a
controlled chain whose scope conditions need to be visible in this paper,
especially for quantitative closure, continuation, and
nonperturbative sectors.

\subsection{Imported core theorem packages}

This paper uses the following OPH theorem packages without re-proving them in full.

\begin{enumerate}[leftmargin=2em]
\item \textbf{Patch-net and overlap language.} Ref.~\cite{ophconsensus} organizes overlap repair,
fixed-point language, cycle holonomy, and gauge-as-gluing in a form that is especially useful
for flavor transport and observer-centric interpretation.
\item \textbf{Screen-side regulated architecture.} Ref.~\cite{ophmicrophysics} gives the
federated echosahedral patch-carrier model, the regulated patch-net embedding theorem, and the
fixed-cutoff edge heat-kernel / Casimir theorem.
\item \textbf{Measurement interface.} Ref.~\cite{ophmicrophysics}, together with the integrated
measurement appendices and the flagship paper \emph{From Observer Consensus to Standard Physics}~\cite{ophsynthesis}, supplies the fixed-cutoff central-record package and, under a separately
declared algebra-state and two-wing representation, the Born-rule,
L\"uders-conditioning, and Bell/CHSH theorems used when
the measurement sections explain how observations pick out definite observer-accessible outcomes and how
declared two-wing quantum laws obey the Tsirelson upper bound. No Bell violation is derived from
the commuting record indicators or from the repair law.
\item \textbf{Forced local Standard Model Lie algebra, conditional matter
image, and conditional compact reconstruction.}
\emph{Recovering Observer Spacetime and Einstein Dynamics from Overlap Consistency and Deriving Standard Model Gauge Structure from Observer Overlap Consistency}~\cite{ophEinstein,ophGauge}
and the dedicated gauge-group fragment prove that complete reversible
response and endogenous proper-carrier transport force
\[
\mathfrak{su}(3)\oplus\mathfrak{su}(2)\oplus\mathfrak u(1).
\]
The separately declared matrix current and rank-15 matter table, together with
anomaly balance and tensor descent, give the maximal faithful image
\[
\frac{\SU(3)\times\SU(2)\times\U(1)}{\mathbb Z_6},
\qquad
N_c=3.
\]
The charge lattice and \(N_c=3\) come from the stated matrix and matter
premises. The quotient is exact for the declared representation, while its
selection as the physical global gauge group is open. Separate intrinsic-CP
and weak-sector premises give \(3\le N_g\le5\). The same
Standard Model gauge paper gives the transportable-sector/Tannaka construction as a
separate conditional compact-group route. Physical source binding of the
finite current and identification with the Tannaka current require the open
producer, intertwiner, and laboratory attachment.
\item \textbf{Conditional topic notes.} The integrated branch appendices carry useful material
on supersymmetry, the finite-quotient baryogenesis source theorem and its open record-generator branch, proton stability, generation structure, and Yukawa hierarchy,
but those sections have to be read with their stated claim boundaries intact. They are sources
for conditional analyses, not recovered-core theorems.
\end{enumerate}

This reading rule governs the whole paper. Structural gauge and carrier results can be used
as structural results. The electroweak branch is a declared quantitative calibration branch. The
Higgs/top critical stage is a conditional downstream calculation on its declared running, matching, and threshold surface. The quark,
charged-lepton, neutrino, and hadron chapters state conditional or comparison
results because their physical premises are incomplete.

\section{Gauge Architecture, Standard Model Structure, and Structural Carriers}

The particle-spectrum derivation starts with the forced local gauge Lie
algebra
\[
\mathfrak{su}(3)\oplus\mathfrak{su}(2)\oplus\mathfrak u(1).
\]
The declared matrix current and rank-15 matter table, determinant balance,
Spin lift, and tensor descent then give the conditional maximal faithful
matter image, charge lattice, and carrier roles before any mass readout. The overlap/transport obstruction calculus and
transportable-sector/Tannaka reconstruction form a separate conditional
classification route. The target-free matter reduct leaves the generation
count unselected. The strengthened finite screen interface selects rank three
under its single-object, faithful-action, and operational-cost premises. The
physical family interpretation and the extra-sector completeness statement
are open. This logical
split matters for this paper. The gauge branch does not fix the carriers'
quantum-particle masses; the action-level and quantum spectral gates below decide whether a
classical mode or particle row exists.

On the declared packet, the gauge/content result used throughout this paper is
\[
\frac{\SU(3)\times\SU(2)\times\U(1)}{\mathbb Z_6},
\qquad
3\le N_g\le 5,
\qquad
N_c=3.
\]
\emph{Recovering Observer Spacetime and Einstein Dynamics from Overlap Consistency} and
\emph{Deriving Standard Model Gauge Structure from Observer Overlap Consistency}~\cite{ophEinstein,ophGauge}, the longer main derivation surface, and the
dedicated gauge-group proof fragment agree on the split. The conditional
finite receipts fix the maximal faithful image, hypercharge lattice, and
\(N_c=3\) from their stated matrix-current and matter premises. They do not
select the physical current, matter action, or global quotient. Intrinsic CP
capability together with weak-sector UV completeness gives
\(3\le N_g\le5\). The exact screen-band theorem selects
rank three on the strengthened finite interface, and the declared Laplacian
simulator realizes the corresponding lowest-positive-frequency residue.
Witten parity is a consistency check. The finite local-domain receipt
checks the separately declared operator \(D_\sigma\otimes I_{45}\) and
conditional gap inheritance. It does not construct a source-selected local
matter carrier or bridge the twelve-port Spin packet to that domain. These
steps do not identify physical matter poles or exclude extra light sectors.

The exterior substrate is an explicit algebra. The ten-row table is its
finite component census. For the supplied carrier \(V=\mathbb C^3\oplus\mathbb
C^2\), the exterior basis is indexed by all \(32\) subsets of the five
coordinate modes. Their color and weak degrees have multiplicities
\(\binom3c\binom2w\). Removing the vacuum degree \((0,0)\) and the top degree
\((3,2)\) leaves exactly the ten component bidegrees used in the finite matter
scan. On every row, the integer charge is \(-2c+3w\), fermionic parity is
\(c+w\bmod2\), and charge conjugation is complement in the three color and
two weak modes. Exterior creation squares to zero and distinct creation
generators anticommute. These are algebraic exclusion statements on the
supplied exterior carrier. They do not select that carrier as physical
matter or establish continuum spin--statistics.

A supplied projective sector partition gives a precise operational
composition with this exterior table. Partition pinching identifies two
finite matrices exactly when every matrix in the sector-preserving commutant
has the same trace statistic on them. Partition averaging is the distinct
commutative public-record readout; both readouts erase cross-sector corners.
If an explicit map assigns each exterior row to a nonzero partition sector
and supplies the corresponding central weight, substituting those labels
preserves the same diagonal \(\mathbb Z_6\) weight kernel. A separate supplied
chiral anomaly-free mask is forced to one of the conjugate parity sectors.
The theorem defines no group action on the projectors and does not attach the
mask to a physical sector action. The maps are premises. It derives no
physical partition, matter action, parity representative, or laboratory
charge.

The Axiom~1/Axiom~2 finite-current route, its conditional matrix and matter
continuation, and the transportable-sector/Tannaka route are separate chains.
Their equality as one
physical gauge-current object and physical source binding of the first route
are open. The rank-three response band therefore remains a candidate family
fiber. Its abstract rank-\(45\) tensor product with the generation table does
not attach that band to three physical families.

For the particle zoo this structural branch fixes three carrier roles, not three particle masses.
The Standard Model gauge paper's carrier-mode theorem then supplies the following conditional action-level
statements. The Maxwell action on the ordinary unbroken electromagnetic vacuum has two transverse
classical modes. The pure Yang--Mills action about a flat connection has \(2\dim G\) perturbative
transverse modes, but the confined color phase has no asserted gauge-invariant asymptotic gluon.
The pure two-derivative Einstein--Hilbert action about Minkowski space has two classical
transverse-traceless modes. In every case a quantum particle additionally requires a positive-
energy physical Hilbert space, a positive-residue physical two-point pole, and the appropriate
asymptotic or deconfinement receipt.

\begin{table}[h]
\centering
\small
\begin{tabular}{@{}>{\raggedright\arraybackslash}p{0.13\textwidth}>{\raggedright\arraybackslash}p{0.12\textwidth}>{\raggedright\arraybackslash}p{0.12\textwidth}>{\raggedright\arraybackslash}p{0.23\textwidth}>{\raggedright\arraybackslash}p{0.26\textwidth}@{}}
\toprule
Carrier role & Theorem scope & Action-level output & Physical degrees of freedom & Quantum-particle requirements \\
\midrule
Electro\-magnetic connection & conditional Maxwell mode & transverse \(k^2=0\) kernel & two transverse classical modes & photon requires quantization and a positive-residue physical pole \\
Color connection & conditional perturbative pure-Yang--Mills mode & one transverse \(k^2=0\) kernel per generator & \(2\dim G\) perturbative modes & deconfined BRST/LSZ receipt required; no free-gluon claim in confined QCD \\
Metric perturbation & conditional pure-Einstein mode & \(D^{\mathrm{TT}}\propto i\Pi^{\mathrm{TT}}/(k^2+i0)\) & two classical TT modes & graviton requires metric quantization, physical Hilbert space, and pole receipt \\
\bottomrule
\end{tabular}
\end{table}

A separate screen-side construction sharpens the boundary of the
electromagnetic row~\cite{ophmicrophysics}. The complete thirty-seam source
maps onto the record lattice
\[
D_6=\left\{z\in\mathbb Z^6:\sum_i z_i\equiv0\pmod2\right\},
\]
whose response-metric completion is three-dimensional. Under proper-carrier
naturality and a unique strictly
positive normalized minimizer, the sixty directed seam moves carry equal
weight and define an exact homogeneous Dirichlet operator. Its plane-wave
symbol has a basis-free momentum-orthogonal fiber of real dimension two at
every nonzero momentum. For a supplied nonzero coordinate dilation, if that
spatial symbol is declared to be an oscillator stiffness, the resulting
conditional mode obeys
\[
A_T''+\Lambda_a(k)A_T=0,\qquad
\omega_a(k)^2=\Lambda_a(k).
\]
The spatial symbol vanishes at the momentum origin, so the nonnegative
square-root branch satisfies \(\omega_a(0)=0\).
The projector, rank-two fiber, scalar degeneracy, oscillator equation, and
zero mode are exact mathematical statements. Gauge redundancy, Gauss law,
the Maxwell action, the physical electromagnetic group, a physical clock,
and a photon pole are absent from this theorem.

The physical electromagnetic diagonal
\(\U(1)_Q\), generated by \(Q=T_3+Y\), also requires the certified
\((\mathbf1,\mathbf2)_{1/2}\) carrier and a selected nonzero neutral
lower-component vacuum. The projective weak-doublet calculation above
identifies the stabilizer of that supplied vacuum vector. A source-selected
scalar action and its physical vacuum are separate requirements. Only after that selection,
together with physical position, clock, frame, propagation, readout,
quantization, and pole data, can the rank-two seam oscillator be tested for
identity with the electromagnetic carrier.

The possible phenomenological routes use different inputs. A direct
time-of-flight calculation is a one-particle propagation test. Electron and
positron dispersion and pair-production cross sections are outside that
calculation. It needs a physical frequency map with a controlled remainder,
position and clock attachment, frame and boost law, cosmological transport,
emission model, detector timing, nuisance treatment, and exclusive readout.
A pair-production threshold or opacity calculation additionally needs the
physical electron and positron actions, a common energy-momentum conservation
law, pair kinematics and interactions, background-radiation transport, cross
sections, and source composition. The finite seam operator supplies neither
physical comparison on its own.

Before physical identification, the complete positive cosine symbol has an
exact all-momentum bound. Its nonnegative auxiliary root \(\Omega\) obeys
\[
|\Omega(k)-\Omega(p)|\le |k-p|
\]
in the selected Euclidean carrier chart. This is a carrier-coordinate
contraction. It supplies no physical position, frequency, clock, field,
wave-packet, frame, scale, or detector map. If the seam symbol is physical
photon frequency squared, its complete cosine form is bounded above by
\(|k|^2\). Positive-mass charged leptons with
Lorentz-invariant positive-energy dispersion then exclude photon decay into an
electron-positron pair under additive conservation. With the soft background
photon Lorentz invariant at leading order, use
\(E_i^2=p_i^2+m_i^2+\delta_{i,2}E_i^4\), with
\(i\in\{\gamma,+,-\}\), \(m_\gamma=0\), and \(m_+=m_-=m_e\). A leading
head-on, collinear threshold
with independent hard-photon, positron, and electron dimension-six
coefficients depends at fixed positron energy share \(0<x<1\) on
\[
\delta_{\gamma,2}-x^3\delta_{+,2}-(1-x)^3\delta_{-,2}.
\]
The equal-share readout has rank one and a two-dimensional coefficient fiber.
The photon-only threshold translation therefore depends on the separate
Lorentz-invariant-lepton premise; the screen theorem does not select it. On
that branch, equal sharing is the unique global optimizer of the leading
head-on, collinear residual. The general independent-lepton and full
anisotropic optimizations are open.

These conditional classical modes should not be confused with quantum-particle mass rows. The
group and gravity branches identify the carrier roles; the displayed actions and phases supply
their classical propagation; and only a separate quantum spectral receipt can establish them as
photon, gluon, or graviton particles. The boson sections address the \(W\), \(Z\), and Higgs
rows on their own electroweak and Higgs/top quantitative stages.

\section{Electroweak Calibration and the Higgs/Top Critical Stage}

The bosonic calculation separates source outputs from comparison coordinates.
The electroweak calibration supplies the \(W\) boson and \(Z\) boson chart,
and the criticality calculation supplies conditional Higgs and top
coordinates. Neither calculation emits an OPH-native particle pole.

\subsection{Electroweak chart on the quantitative branch}

The electroweak derivation consists of a single-\(P\) running family followed by a
reduced two-scalar carrier, a selector on that carrier, an exact carrier mass
chart, and a conditional quotient-transport theorem beyond the selected
carrier. In practical terms, the construction starts from
the declared pixel input \(P\), builds the running electroweak family, reduces it to the selected
carrier, reads the \(W\) boson/\(Z\) boson pair from the selected point on that carrier,
and then evaluates the mass pair together with the Ward-projected electromagnetic transport family.
On this branch, a source-only prediction must respect that ordering: first
certify the shared pixel input \(P\) on the declared
running/matching/threshold/scheme surface, thereby fixing the source basis
\[
(\alpha_U,\alpha_{2,m_Z},\alpha_{Y,m_Z},\eta_{\mathrm{source}},v),
\]
with \(\alpha_U(P)\), equivalently \(t_U(P)\) and \(t_{\mathrm{tr}}(P)\), fixed by the
same forward pixel-closure solve,
and only then emit the electroweak transport family from that source basis.
The stated theorem does not contain this full chain. Its \(W/Z\) rows are
data-comparison adapter coordinates rather than OPH-native poles.

The calibration couplings are read at
the source scale \(91.5883371732~\mathrm{GeV}\), a coordinate of the source
solve, rather than at the measured \(Z\) target. Evaluation of the conditional
pole map in the raw pre-carrier basis omits the compact-carrier hypercharge
step. The selector-consistent conditional evaluation lands the neutral
coordinate near \(90.66~\mathrm{GeV}\), about half a GeV below the converted
reference target. The raw-basis and selector-consistent numbers are
conditional diagnostics of the imported prescription. Neither is an
OPH-native pole.

The same quantitative branch has one exact golden-ratio benchmark. If one writes
\[
x(C):=\frac{S_{\mathrm{gen}}(C)}{S_{\mathrm{bulk}}(C)}
=
1+\frac{\langle L_C\rangle}{S_{\mathrm{bulk}}(C)}
\]
for the total/bulk/edge hierarchy and imposes exact self-similar balance
\[
\frac{S_{\mathrm{gen}}(C)}{S_{\mathrm{bulk}}(C)}
=
\frac{S_{\mathrm{bulk}}(C)}{\langle L_C\rangle},
\]
then \(x\) obeys \(x^2-x-1=0\), so the unique positive equilibrium point is
\[
x=\varphi:=\frac{1+\sqrt5}{2}.
\]
Equivalently the order parameter \(A_\varphi(x)=x-1-\frac1x\) vanishes there. The synthesis paper
records the unique numerical root of the incomplete outer/inner declared map. The point of
the equilibrium theorem here is that the proximity to \(\varphi\) has a structural reason. Exact
balance is too symmetric to support durable records, structure, and dynamics, so the declared map
assigns a small
equilibrium-breaking detuning away from that balance point.

The compare-only \(W/Z\) adapter rows in the final bundle are
\[
m_W = 80.3625~\mathrm{GeV},
\qquad
m_Z = 91.1879~\mathrm{GeV}.
\]
These numbers sit on a target-conditioned validation surface and do not form
a target-free prediction theorem.
The exact selected-carrier chart gives
the local pair
\[
m_W^{\mathrm{carrier}} = 80.38629169244275~\mathrm{GeV},
\qquad
m_Z^{\mathrm{carrier}} = 91.18290444674243~\mathrm{GeV},
\]
so the mathematics distinguishes the selected-carrier chart from the
candidate value-law pair
\[
\bigl(80.3770000154,91.1879780779\bigr)\,\mathrm{GeV}.
\]
It also distinguishes the conditional quotient-transport implication from the
measured-reference inverse repair surface.
The rounded pair \((80.377,91.1879780919)\,\mathrm{GeV}\) is a boundary alias
of the candidate, not the inverse adapter. None is a physical pole theorem.

\begin{theorem}[Measured-pair reparametrization (calibration): coherent electroweak inverse fit]
Let the emitted running/core electroweak basis be
\[
Q_{\mathrm{run}}(P)=
\bigl(\alpha_{Y,m_Z}(P),\alpha_{2,m_Z}(P),v_{\mathrm{chart}}(P),\eta_{\mathrm{source}}(P)\bigr),
\]
and write
\[
\tau_Y(\tau_2)
=
-\frac{\tau_2+2\eta_{\mathrm{source}}}{1+4\tau_2^2},
\qquad
n_{\mathrm{fiber}}(\tau_2)
=
1+\frac{\alpha_{Y,m_Z}\tau_Y(\tau_2)+\alpha_{2,m_Z}\tau_2}
{\alpha_{Y,m_Z}+\alpha_{2,m_Z}}.
\]
Freeze one authoritative target pair
\[
\mathcal T^\dagger=(M_W^\dagger,M_Z^\dagger),
\qquad 0<M_W^\dagger<M_Z^\dagger,
\]
and define the charged anchor
\[
\tau_{2,W}^\dagger
=
\frac{M_W^{\dagger\,2}}{\pi v^2\alpha_{2,m_Z}}-1,
\qquad
\delta\alpha_2^\dagger
=
\frac{M_W^{\dagger\,2}}{\pi v^2}-\alpha_{2,m_Z}.
\]
Then the fiber-parallel hypercharge leg is
\[
\tau_Y^\dagger=\tau_Y(\tau_{2,W}^\dagger),
\qquad
n_{\mathrm{fiber}}^\dagger=n_{\mathrm{fiber}}(\tau_{2,W}^\dagger),
\]
\[
M_{Z,\mathrm{fiber}}^\dagger
=
v\sqrt{\pi(\alpha_{Y,m_Z}+\alpha_{2,m_Z})\,n_{\mathrm{fiber}}^\dagger},
\]
and the orthogonal neutral-shear scalar is
\[
\delta M_Z^\perp=M_Z^\dagger-M_{Z,\mathrm{fiber}}^\dagger,
\]
equivalently
\[
\delta n^\dagger
=
\frac{(M_Z^\dagger+M_{Z,\mathrm{fiber}}^\dagger)\,\delta M_Z^\perp}
{\pi v^2(\alpha_{Y,m_Z}+\alpha_{2,m_Z})},
\qquad
\delta\alpha_Y^\perp
=
\frac{(M_Z^\dagger+M_{Z,\mathrm{fiber}}^\dagger)\,\delta M_Z^\perp}{\pi v^2}.
\]
The fiber-parallel hypercharge motion is
\[
\delta\alpha_Y^\parallel
=
\alpha_{Y,m_Z}\,
\frac{8\eta_{\mathrm{source}}(\tau_{2,W}^\dagger)^2-\tau_{2,W}^\dagger}
{1+4(\tau_{2,W}^\dagger)^2}.
\]
Therefore the unique coherent repair package is
\[
\Sigma_{EW}^\dagger
=
\bigl(\delta\alpha_2^\dagger,\delta\alpha_Y^\parallel,\delta\alpha_Y^\perp\bigr),
\]
equivalently \((\tau_{2,W}^\dagger,\delta n^\dagger)\). If the selected carrier anchor is the
\(\tau_2=0\) fiber point, so that
\[
\alpha_{2,*}=\alpha_{2,m_Z},
\qquad
\alpha_{Y,*}=\alpha_{Y,m_Z}(1-2\eta_{\mathrm{source}}),
\]
then the coherent validation couplings are
\[
\alpha_{2,\dagger}=\alpha_{2,*}+\delta\alpha_2^\dagger,
\qquad
\alpha_{Y,\dagger}=\alpha_{Y,*}+\delta\alpha_Y^\parallel+\delta\alpha_Y^\perp,
\]
and they emit one coherent quintet
\[
M_W^\dagger=v\sqrt{\pi\alpha_{2,\dagger}},
\qquad
M_Z^\dagger=v\sqrt{\pi(\alpha_{Y,\dagger}+\alpha_{2,\dagger})},
\]
\[
\alpha_{\mathrm{em},\dagger}^{-1}
=
\frac{\alpha_{Y,\dagger}+\alpha_{2,\dagger}}
{\alpha_{Y,\dagger}\alpha_{2,\dagger}},
\qquad
\sin^2\theta_{W,\dagger}
=
\frac{\alpha_{Y,\dagger}}{\alpha_{Y,\dagger}+\alpha_{2,\dagger}}.
\]
The target-conditioned electroweak validation law is one unique coherent
value-emission law on top of the closed fiber map.
\end{theorem}

For the reference-fitted inverse-adapter comparison,
\[
\begin{aligned}
\alpha_{2,\star}&=0.03377843630219015,\\
\alpha_{Y,\star}&=0.009682831911900495,\\
\eta_{\mathrm{source}}&=0.022147000871961295,
v&=246.76711732749683~\mathrm{GeV},\\
M_W^\dagger&=80.3625~\mathrm{GeV},\\
M_Z^\dagger&=91.1879~\mathrm{GeV},
\end{aligned}
\]
and the law emits
\[
\begin{aligned}
\tau_{2,W}^\dagger&=-0.0005918464744071317,\\
\delta\alpha_2^\dagger&=-1.9991648436440412\times10^{-5},\\
M_{Z,\mathrm{fiber}}^\dagger&=91.16822266259587~\mathrm{GeV},\\
\delta M_Z^\perp&=19.6773374041328~\mathrm{MeV},\\
\delta\alpha_Y^\parallel&=5.9969727526906094\times10^{-6},\\
\delta\alpha_Y^\perp&=1.8756950557734103\times10^{-5},\\
\delta n^\dagger&=4.2716772021678714\times10^{-4}.
\end{aligned}
\]
So the coherent validation couplings are
\[
\begin{aligned}
\alpha_{2,\dagger}&=0.03375844465375371,\\
\alpha_{Y,\dagger}&=0.00970758583521092,
\end{aligned}
\]
and the same coherent branch emits
\[
M_W^\dagger=80.3625~\mathrm{GeV},
\qquad
M_Z^\dagger=91.1879~\mathrm{GeV},
\]
\[
\begin{aligned}
\alpha_{\mathrm{em},\dagger}^{-1}&=132.6344411210187,\\
\sin^2\theta_{W,\dagger}&=0.22333729871365943.
\end{aligned}
\]
These two carrier-side scalars belong to the target-conditioned validation
surface. They are not the physical electromagnetic readout on the
Ward-projected transport branch.
That inverse-fit law provides an exact coherent calibration surface: it is a
measured-pair reparametrization. On that
target-conditioned surface the reference \(W/Z\) pair is hit exactly because the basis is solved from
that pair; the hit is the defining condition, and the branch is calibration.
It carries no predictive force.
The forward running/chart solve gives
\((80.330,\,91.119)\)~GeV. The Particle Data Group (PDG) comparison values are
mass-dependent-width Breit--Wigner parameters, while complex-pole masses define another
convention. As a stated convention diagnostic under the complex-pole conversion of the
PDG 2026 reference values, the energy-pole coordinates are
\((M_W,M_Z)=(80.3411410,91.1623040)\)~GeV, whereas the older
\(\sqrt{\operatorname{Re}s}\) coordinates are
\((80.3340218,91.1537725)\)~GeV. These two definitions must not be mixed.
The running/chart coordinates carry no OPH theory covariance, and the
physical readout contract is open. No comparison to physical resonance
parameters is defined for them.
The candidate forward value law has a stronger conditional theorem: it follows
uniquely from a finite quotient-path certificate. The certificate is not
emitted by the three axioms or by the existing carrier artifacts.

\begin{definition}[Finite quotient-path certificate]\label{def:particle-d10-qt}
After quotienting gauge representatives, port labels, scheduler data, and
other hidden presentation coordinates, require:
\begin{enumerate}[leftmargin=2em]
\item the real, CP-even, color-singlet, charge-preserving response through two
returns is exactly \(\mathbb Rq_2\oplus\mathbb Rq_n\), with no third scalar or
charged-neutral off-block term;
\item the primitive response is bilinear in
\((\eta_{\mathrm{source}},\alpha_U)\), one primitive event has fixed unit
response normalization, and the quotient probability measure gives each of
the four transmutation slots weight \(1/4\), yielding
\(\lambda_{\mathrm{EW}}=\eta_{\mathrm{source}}\alpha_U/4\);
\item explicit carrier path lists, carrying one factor of
\(\eta_{\mathrm{source}}\) per return and uniform color-orbit measure \(1/3\),
have primitive, one-return, and two-return incidences
\((1,2/3,1)\) for \(q_2\) and \((1,4/3,2)\) for \(q_n\), while the diagonal
\(\mathbb Z_6\) projector in the declared response representation has
normalized trace \(1/6\), source amplitude \(\rho_{\mathrm{EW}}\), and subtracts
\((\rho_{\mathrm{EW}}/6)\eta_{\mathrm{source}}^2\) from both degree-two
characters;
\item mismatch descent fixes the charged sign as contracting and the neutral
sign as uplifting, and the parallel hypercharge fibre is the least-norm
minimizer with Gram factor \(1+4\tau_2^2\) and residual pairing
\(\tau_2+2\eta_{\mathrm{source}}\);
\item the listed paths exhaust the admissible class: deeper paths,
alternative central weights, and mixed terms are absent, quotient-trivial, or
separated by a positive target-independent selector gap.
\end{enumerate}
\end{definition}

\begin{theorem}[Conditional quotient-transport value law]\label{thm:particle-conditional-d10-qt}
Let the source-only electroweak basis be
\[
Q_{\mathrm{src}}(P)=
\bigl(\alpha_U(P),\alpha_{2,m_Z}(P),\alpha_{Y,m_Z}(P),
\eta_{\mathrm{source}}(P),v_{\mathrm{chart}}(P)\bigr),
\]
and define
\[
\rho_{\mathrm{EW}}
:=
\frac{\alpha_{2,m_Z}-\alpha_{Y,m_Z}}{\alpha_{2,m_Z}+\alpha_{Y,m_Z}},
\lambda_{\mathrm{EW}}
:=
\frac{\eta_{\mathrm{source}}\alpha_U}{4}.
\]
Assume that every entry is emitted on the same source branch, that
\(\eta_{\mathrm{source}}=\rho_{\mathrm{EW}}\alpha_U\), and that no measured
\(W/Z\), measured \(v\), fitted electroweak parameter, or calibrated proxy is
an ancestor of the tuple. Hence
\(\lambda_{\mathrm{EW}}=\eta_{\mathrm{source}}^2/(4\rho_{\mathrm{EW}})\).
Assume also the finite certificate of Definition~\ref{def:particle-d10-qt} and the
physical-domain inequalities
\(\alpha_2'>0\), \(\alpha_Y'>0\).
Then the beyond-selected-carrier transport map is uniquely
\[
\tau_2^{\mathrm{exact}}
=
-\lambda_{\mathrm{EW}}
\left(
1+\frac23\eta_{\mathrm{source}}
+
\left(1-\frac{\rho_{\mathrm{EW}}}{6}\right)\eta_{\mathrm{source}}^2
\right),
\]
\[
\delta n^{\mathrm{exact}}
=
\lambda_{\mathrm{EW}}
\left(
1+\frac43\eta_{\mathrm{source}}
+
\left(2-\frac{\rho_{\mathrm{EW}}}{6}\right)\eta_{\mathrm{source}}^2
\right),
\]
with fiber hypercharge transport
\[
\tau_Y^{\mathrm{fiber}}
=
-\frac{\tau_2^{\mathrm{exact}}+2\eta_{\mathrm{source}}}
{1+4(\tau_2^{\mathrm{exact}})^2}.
\]
The coherent transport shifts are
\[
\delta\alpha_2=\alpha_{2,m_Z}\tau_2^{\mathrm{exact}},
\qquad
\delta\alpha_Y^{\parallel}
=
\alpha_{Y,m_Z}
\frac{8\eta_{\mathrm{source}}(\tau_2^{\mathrm{exact}})^2-\tau_2^{\mathrm{exact}}}
{1+4(\tau_2^{\mathrm{exact}})^2},
\]
\[
\delta\alpha_Y^{\perp}
=
(\alpha_{2,m_Z}+\alpha_{Y,m_Z})\,\delta n^{\mathrm{exact}}.
\]
With
\[
\alpha_{2,*}=\alpha_{2,m_Z},
\qquad
\alpha_{Y,*}=\alpha_{Y,m_Z}(1-2\eta_{\mathrm{source}}),
\]
the coherent transport couplings are
\[
\alpha_2'=\alpha_{2,*}+\delta\alpha_2,
\qquad
\alpha_Y'=\alpha_{Y,*}+\delta\alpha_Y^{\parallel}+\delta\alpha_Y^{\perp},
\]
and they emit one coherent mass chart
\[
M_W^{\mathrm{chart}}=v\sqrt{\pi\alpha_2'},
\qquad
M_Z^{\mathrm{chart}}=v\sqrt{\pi(\alpha_2'+\alpha_Y')},
\qquad
v=v_{\mathrm{chart}}(P).
\]
\[
a_0(P)
:=
\frac{\alpha_{2,m_Z}(P)+\alpha_{Y,m_Z}(P)}
{\alpha_{2,m_Z}(P)\alpha_{Y,m_Z}(P)},
\qquad
s_0(P)
:=
\frac{\alpha_{Y,m_Z}(P)}
{\alpha_{2,m_Z}(P)+\alpha_{Y,m_Z}(P)}.
\]
Thus the certificate fixes the candidate mass pair together with the
running-family anchor \((a_0,s_0,v)\) from the source basis. The physical electromagnetic row
is read from the Ward-projected unbroken \(\U(1)_Q\) channel in
Theorem~\ref{thm:ward_u1q_transport}.
\end{theorem}

\begin{proof}
The first condition reduces every admissible response to the two displayed
coordinates. The second fixes their common primitive activity. The third
gives the charged and neutral path characters, including the common
\(\mathbb Z_6\) subtraction, and the fourth fixes their signs. The fibre functional
\[
\mathcal J_Y(t)
=\frac12(1+4(\tau_2^{\mathrm{exact}})^2)t^2
+(\tau_2^{\mathrm{exact}}+2\eta_{\mathrm{source}})t
\]
is strictly convex, so its unique minimizer is
\(-(\tau_2^{\mathrm{exact}}+2\eta_{\mathrm{source}})
 /(1+4(\tau_2^{\mathrm{exact}})^2)\).
Substitution gives the repaired couplings and the one-Higgs mass chart. The
exhaustiveness condition
excludes every output-changing admissible deformation, which proves uniqueness
on the certified class.
\end{proof}

Using the same declared electroweak validation basis
\[
\begin{aligned}
\alpha_{2,m_Z}&=0.03377843630219015,\\
\alpha_{Y,m_Z}&=0.010131601067241624,\\
\alpha_U&=0.04112498041477454,\\
\eta_{\mathrm{source}}&=0.022147000871961295,\\
v&=246.76711732749683~\mathrm{GeV},
\end{aligned}
\]
the conditional quotient-transport law evaluates to
\[
M_W^{\mathrm{chart}}=80.37700001539531~\mathrm{GeV},
\qquad
M_Z^{\mathrm{chart}}=91.18797807794321~\mathrm{GeV},
\]
\[
\begin{aligned}
a_0&=128.30576920234813,\\
s_0&=0.23073542347506173.
\end{aligned}
\]
This specialization verifies the candidate value-law algebra. It does not
derive the finite quotient-path certificate, establish that this calibration
tuple belongs to the source-only pixel branch, or identify \(W,Z\) as complex
propagator poles. A proof requires the finite path lists,
quotient canonicalizer, exact incidence and central-trace checks, fibre Gram
calculation, deformation enumeration or positive selector gap, and a same-branch
target-independence proof.

\begin{proposition}[Local nondegeneracy of the two-output mass chart]
Define
\[
n_{\mathrm{fib}}(\tau_2)
:=
1+\frac{\alpha_2\tau_2+\alpha_Y\tau_Y(\tau_2)}
{\alpha_2+\alpha_Y}.
\]
On the physical domain
\(1+\tau_2>0\) and
\(n_{\mathrm{fib}}(\tau_2)+\delta n>0\),
\[
\det
\frac{\partial(M_W^{\mathrm{chart}},M_Z^{\mathrm{chart}})}
{\partial(\tau_2,\delta n)}
=
\frac{M_W^{\mathrm{chart}}M_Z^{\mathrm{chart}}}
{4(1+\tau_2)\,[n_{\mathrm{fib}}(\tau_2)+\delta n]}
>0.
\]
Holding \((\alpha_2,\alpha_Y,v,\eta)\) fixed, the two-output chart is locally
nondegenerate. This determinant neither proves the two-coordinate response
condition nor excludes an additional
physical response channel. Within the displayed two-coordinate chart, the
missing content is a source selector rather than another fit coordinate.
\end{proposition}

\begin{proof}
\(M_W^{\mathrm{chart}}\) is independent of \(\delta n\), while
\[
\frac{\partial M_W^{\mathrm{chart}}}{\partial\tau_2}
=\frac{M_W^{\mathrm{chart}}}{2(1+\tau_2)},
\qquad
\frac{\partial M_Z^{\mathrm{chart}}}{\partial\delta n}
=\frac{M_Z^{\mathrm{chart}}}{2[n_{\mathrm{fib}}(\tau_2)+\delta n]}.
\]
The determinant is the product of these diagonal entries.
\end{proof}

\begin{theorem}[Ward-projected \texorpdfstring{\(\U(1)_Q\)}{U(1)_Q} transport law and Thomson-limit electromagnetic readout]\label{thm:ward_u1q_transport}
Let the source-only running family be fixed by the forward pixel solve on the declared
running/matching/threshold/scheme surface, with source basis
\[
Q_{\mathrm{src}}(P)=
\bigl(\alpha_U(P),\alpha_{2,m_Z}(P),\alpha_{Y,m_Z}(P),
\eta_{\mathrm{source}}(P),v_{\mathrm{chart}}(P)\bigr).
\]
Assume the realized electric-charge branch, so that
\[
Q=T_3+Y
\]
on the realized low-energy branch and color-singlet physical states carry integer \(Q\)-charge.
Assume further that the post-selector electroweak transport object is the populated quotient-local
kernel
\[
K^{\mathrm{EW}}_{\mathrm{src}}(q^2;P)
=
\bigl\{\Pi_{AA}(q^2;P),\Pi_{AZ}(q^2;P),\Pi_{ZZ}(q^2;P),\Pi_{WW}(q^2;P)\bigr\}
\]
on the declared physical observable algebra, together with a Ward projector
\[
\mathcal W_Q:
K^{\mathrm{EW}}_{\mathrm{src}}(q^2;P)\longrightarrow \Pi_Q(q^2;P)
\]
onto the unbroken electromagnetic channel satisfying at \(q^2=0\)
\[
\mathcal W_Q[\Pi_{AZ}(0;P)]=0,
\qquad
\mathcal W_Q[\Pi_{AA}(0;P)]\neq 0.
\]
Assume the abelian projected edge-sector probabilities obey the \(\U(1)\) heat-kernel law
\[
p_n(q^2;P)\propto e^{-t_Q(q^2;P)\lambda_n},
\qquad
\lambda_n=n^2,
\]
equivalently
\[
g_Q^2(q^2;P)=\frac{t_Q(q^2;P)}{2\pi},
\]
and that the scalar readout uses the same source family, declared
renormalization scheme, and originating electroweak kernel throughout.
Then the unique electromagnetic coupling readout on the Ward-projected branch is
\[
\alpha_{\mathrm{em}}^{-1}(q^2;P)=\frac{8\pi^2}{t_Q(q^2;P)}.
\]
If
\[
a_0(P):=\alpha_{\mathrm{em}}^{-1}(m_Z^2;P),
\]
then the unique zero-normalized affine scalar on the same source-locked family is
\[
\delta_\alpha(q^2;m_Z^2;P)
:=
\frac{t_Q(m_Z^2;P)}{t_Q(q^2;P)}-1,
\]
so that
\[
\alpha_{\mathrm{em}}^{-1}(q^2;P)
=
a_0(P)\bigl(1+\delta_\alpha(q^2;m_Z^2;P)\bigr).
\]
Hence the Thomson-limit electromagnetic readout on that same
Ward-projected electromagnetic kernel:
\[
\alpha_{\mathrm{Th}}^{-1}(P)
:=
\lim_{q^2\to0}\alpha_{\mathrm{em}}^{-1}(q^2;P)
=
a_0(P)\,\frac{t_Q(m_Z^2;P)}{t_Q(0;P)}.
\]
Consequently, the low-energy electromagnetic row on this branch is read as the Thomson endpoint of
the electroweak transport family instead of as a mass-chart byproduct.
\end{theorem}

\begin{proof}
The realized electric-charge branch fixes the charge operator by
\(Q=T_3+Y\), and the
quotient-protected charge theorem fixes the physical color-singlet \(Q\)-lattice to be integral.
The abelian edge-sector theorem supplies the heat-kernel extraction rule
\[
p_n\propto e^{-t_Q\lambda_n},
\qquad
g_Q^2=\frac{t_Q}{2\pi},
\]
so the Ward-projected \(Q\)-channel determines a unique transport time \(t_Q(q^2;P)\). Because
the Ward projector kills \(A\) - \(Z\) mixing at \(q^2=0\), the selected channel is the physical
unbroken electromagnetic channel. Then
\[
\alpha_{\mathrm{em}}=\frac{g_Q^2}{4\pi}=\frac{t_Q}{8\pi^2},
\]
hence
\[
\alpha_{\mathrm{em}}^{-1}=\frac{8\pi^2}{t_Q}.
\]
Evaluating at \(q^2=m_Z^2\) defines the source-locked anchor \(a_0(P)\), and dividing by the same
expression at general \(q^2\) yields
\[
\delta_\alpha=\frac{t_Q(m_Z^2)}{t_Q(q^2)}-1.
\]
The provenance-equality clause forces one running-family source, one frozen scheme, and one origin
kernel for all scalar readouts, so no mixed-family or \(Z\)-only surrogate is promotable. The
Thomson formula is the \(q^2\to0\) limit of the same identity.
\end{proof}

\begin{corollary}[Conditional Ward-projected Maxwell normalization]\label{cor:ward_maxwell_normalization}
On the same Ward-projected \(\U(1)_Q\) branch, assume in addition the low-energy Maxwell action with
the displayed nonzero kinetic normalization. Then the electromagnetic field strength and current obey
\[
F_Q=dA_Q,\qquad dF_Q=0,\qquad d{*}F_Q=g_Q^2(q^2;P){*}J_Q,
\]
with
\[
g_Q^2(q^2;P)=\frac{t_Q(q^2;P)}{2\pi},
\qquad
\alpha_{\mathrm{em}}^{-1}(q^2;P)=\frac{8\pi^2}{t_Q(q^2;P)}.
\]
At the Thomson endpoint, this is the Maxwell normalization used by the public
fine-structure row.
\end{corollary}

\begin{proof}
The compact reconstruction supplies the unbroken abelian electromagnetic connection direction
\(\U(1)_Q\) with charge generator \(Q=T_3+Y\), but does not supply the Maxwell kinetic action.
On an abelian factor the compact-gauge
curvature restricts to \(F_Q=dA_Q\), so \(d^2=0\) gives \(dF_Q=0\). Varying the quadratic
electromagnetic action in the OPH coupling convention gives \(d{*}F_Q=g_Q^2{*}J_Q\). The theorem
above identifies the same branch coupling as \(g_Q^2=t_Q/(2\pi)\), hence
\(\alpha_{\mathrm{em}}^{-1}=8\pi^2/t_Q\).
\end{proof}

\begin{corollary}[Source-only admissibility criterion on the Ward-projected electroweak branch]
The Ward-projected electroweak branch supports a source-only fine-structure row only when the populated
Ward-projected transport kernel satisfies
\[
\lim_{q^2\to0}\alpha_{\mathrm{em}}^{-1}(q^2;P)=\alpha_{\mathrm{Th}}^{-1}(P),
\]
on the same source-locked family and without any separate endpoint value feeding that branch.
Concretely, the evidence packet must contain the source-derived Ward-current
spectral measure, source Jacobi kernel, electromagnetic normalization,
same-scheme finite remainder, no-target-leak dependency graph, full pixel-map
contraction interval, and a prediction interval narrower than the declared
direct-measurement interval. Without those objects, the public endpoint is an empirical
hadron-closure row.
\end{corollary}

This criterion is the two-current marginal of the larger source-derived hadronic spectral
backend. It is adequate for the Thomson endpoint and HVP branch only if the same source law also
records the QCD quotient ensemble, source parameter map, Ward current normalization, systematics
accounting, and no-target-leak DAG. Full hadronic precision claims require the higher-point and
transition spectral exports from that same backend rather than reusing the scalar fine-structure
residual.

The endpoint table records the source-locked anchor from the runtime candidate,
\[
a_0(P)=128.3079654732862482099611087417567\ldots,
\]
and the OPH-plus-empirical hadron-closure value
\[
\alpha_{\mathrm{emp}}^{-1}(0)=136.3827548174946,
\qquad
\alpha_{\mathrm{emp}}^{-1}(0)\in[136.3670480603,136.3984651934].
\]
The source derivation contains no built-in reference inverse-\(\alpha\) value.
The measured endpoint
\(137.035999177(21)\) is excluded as a source-only closure-solve input and is a separate
comparison coordinate.
The source-side calculation emits the certified value
\(\alpha_{\mathrm{root}}^{-1}=136.994835177413\ldots\), so its difference
from the endpoint readout is an endpoint-and-matching comparison packet.
This certificate establishes a unique root of the incomplete declared numerical map. It does not
derive a relation between that root and the physical Thomson endpoint.

\begin{remark}
The empirical closure row is a Thomson-limit inverse fine-structure evaluation using external
hadron data. The source anchor
\(a_0(P)=128.3079654732862482099611087417567\ldots\) is the electroweak-scale running-family value at
\(m_Z^2\), while the measured row is separate. The 2022 CODATA/NIST reference value is
\(\alpha^{-1}=137.035999177(21)\)~\cite{codata2022}.
\end{remark}

The measured-reference inverse repair surface provides diagnostic validation
and carries no predictive force.
The companion carrier-side pair
\[
\begin{aligned}
\alpha_{\mathrm{em},\dagger}^{-1}&=132.6344411210187,\\
\sin^2\theta_{W,\dagger}&=0.22333729871365943
\end{aligned}
\]
is mass-chart bookkeeping on that target-conditioned validation surface. It
is not the physical electromagnetic
readout on the Ward-projected \(\U(1)_Q\) branch.

The electroweak quantitative branch contains an underdetermination theorem for
the quadratic repair family, a color-balanced descent, a candidate mass
emitter, and the Ward-projected electromagnetic transport theorem above the
source-locked running-family anchor. All of them use one source family, one
declared renormalization scheme, and one originating electroweak kernel. The
forward transmutation calculation reconstructs the same \(\alpha_U(P)\),
\(t_U(P)\), and \(t_{\mathrm{tr}}(P)\) as the pixel-closure solve without
reading them back from measured couplings. This separation does not establish
target independence. This construction supplies neither the finite
quotient-path certificate nor a physical pole.

\subsection{Hierarchy interpretation on the local transmutation branch}

The selected branch certifies a dimensionless electroweak
hierarchy/naturality relation. The large
cosmic horizon ratio and the electroweak hierarchy belong to different OPH branches.
A value \(N_{\mathrm{CRC}}\) would be the output of a positive direct
correctable-public-record branch. The fixed-cutoff packet and bounded
counterfamily do not emit such a value. Its de Sitter interpretation and
electroweak identification would require the horizon--record identification
and the common screen/electroweak load-carrier identification, respectively. The local pixel/transmutation readout is the ratio
\[
\frac{v}{E_\star}
=
P_\star^{-1/2}
\exp\!\left[-\frac{2\pi}{4\alpha_U(P_\star)}\right].
\]
The integer \(4\) has an exact representation-theoretic witness on the
selected exterior package: three color copies of the weak doublet \(Q\) plus
one lepton doublet \(L\). Identifying this multiplicity with the
transmutation coefficient is conditional on an intertwiner between the
port and weak carriers and on equality of their normalized load traces. It is not the
\(SU(2)\) beta-function coefficient, and its physical load interpretation
cannot be inferred from the dimension count alone.
The hierarchy certificate records the interval
\[
I_U=[0.041123336195630494,\;0.041125336195630496],
\]
the Krawczyk inclusion
\[
K(I_U)\subset
[0.0411243357185544983,\;0.0411243366727064662]
\subset \mathrm{int}(I_U),
\]
and a strictly negative derivative enclosure
\([-10.995768,-10.985284]\). Thus the local source equation has a unique zero in
that enclosure. On the public endpoint branch, at the CODATA-located
comparison pixel \(P_C\),
\[
\begin{aligned}
\alpha_U(P_C)&=0.041124336195630495,\\
\frac{v}{E_\star}&=2.0199803239725553\times10^{-17}.
\end{aligned}
\]
The source-only branch excludes the public Thomson endpoint upstream and has certified forward
pixel \(P_{\mathrm{fwd}}=1.630972095858897\ldots\). These two named branches
must not be merged. The dimensionless ratio is generated by the ordinary
exponential transmutation factor controlled by the unified diffusion coupling;
a weak scale in GeV additionally requires an independently source-closed
\(E_\star\).

This is also the OPH reading of Higgs naturalness on the declared
electroweak/criticality surface. The
observer-visible scalar mass is the normal-form readout
\[
m_H=H_{\mathrm{OPH}}(P_\star)
\]
on the selected branch. The split into a bare Higgs mass plus a cutoff correction is a regulator
coordinate split outside the observer-facing physical variable. A scalar relevant deformation absent
from the retained branch constraints is off-branch; a scalar deformation present on the selected
branch has its value fixed by the branch equations. The claim inherits the
electroweak/criticality quantitative
surface and the \(\mathcal R_U\) precision policy. The selected exact source-to-Higgs
coarse-graining square has \(\epsilon_H=0\), with
\(\epsilon_H\in[0,0]\), as stated below.

\subsection{Local/global resonance continuation for the hierarchy}

The local/global hierarchy construction compares the electroweak transmutation
exponent with a proposed
global capacity branch at the same joint pair \((P_\star,N_{\mathrm{CRC}})\). Conditional on
an executed direct public-record fixed point and the common
screen/electroweak load-carrier identification, the target relation is
\[
t_{\mathrm{tr}}(P_\star)
=
\frac{P_\star}{12}
\log\!\left(\frac{N_{\mathrm{CRC}}}{\pi}\right),
\]
or equivalently
\[
\frac{v}{E_{\mathrm{cell}}}
=
\left(\frac{N_{\mathrm{CRC}}}{\pi}\right)^{-P_\star/12}.
\]
This packages the mathematical electroweak hierarchy bridge as a resonance between the local
declared-map root and a conditional global screen-capacity coordinate. That screen branch supplies
twelve curvature ports. The exterior package supplies four weak-doublet
copies, and the common screen/electroweak load-carrier map is the physical identification
that makes their additive load the screen load. The oriented 24-slot register
and product-adjoint count are independent bookkeeping facts and carry no load
or clock implication.

For comparison, the observer-visible product adjoint has the independent count. On the realized
product branch,
\[
m_{\rm rep}
=2\dim(\mathfrak{su}(3)\oplus\mathfrak{su}(2)\oplus\mathfrak u(1))
=2(8+3+1)
=24.
\]
The factor \(2\) orients each product-adjoint channel and does not follow from
the screen register. This count does not enter the direct capacity or
common-load proof. The
\(\mathfrak{su}(5)\) adjoint has the same single-orientation integer for a different support:
its \(X/Y\) mixed gauge channels are absent from the OPH product branch.
This channel exclusion is narrower than proton stability. Conditional on the
declared one-generation matter table and baryon and lepton labels, the exact
dimension-six gauge and Lorentz census admits \(QQQL\), \(QQ\,u_Re_R\),
\(QLu_Rd_R\), and \(u_Ru_Rd_Re_R\). Their exterior-algebra
representatives are nonzero. The finite gauge and diagonal \(\mathbb Z_6\)
data supply no Wilson coefficient, physical decay amplitude, QCD matrix
element, or lifetime.

After whole-fiber scalarization and faithful capacity-carrier representation,
the direct active-capacity map is deflationary. It therefore has no positive
two-sided Banach attractor on a domain containing smaller capacities. Under
confusability-reflecting capacity extension, however, iteration from the declared
maximum boundary dimension stabilizes finitely at the greatest fixed point.
This order-theoretic result selects the greatest fixed point only after a
scalar continuation map and its domain have been supplied. It does not choose
the physical continuation map. A bounded counterfamily with shared base,
positivity, and carrier controls does not entail a finite-size slack law with
one physical zero. Universal all-rung membership of its countermodels in the
complete A1--A3 capacity-source contract and the executable-to-Lean bridge are
open. Direct \(N\) is not evaluable on that incomplete antecedent. An
independently produced \(\rho_{\rm op}\) is available as a
commuting-square test against \(\log M_0\); no derivative is intrinsic
at finite integer dimension.

\begin{theorem}[Conditional common-load electroweak bridge]\label{thm:ew-tick-projection}
Assume \(N_\star=\log M_0(\mathfrak U_{N_\star})\) is supplied by the
independent correctable-public-record closure. Assume the screen-sieve branch independently supplies
\[
\Gamma_{\rm scr}=\frac{P}{12}\log\!\left(\frac{N_\star}{\pi}\right),
\]
the common screen/electroweak load-carrier map identifies
\(\Gamma_{\rm scr}=\log(E_{\rm cell}/v)\), and the independent electroweak source
relation gives
\(\log(E_{\rm cell}/v)=\pi/[2\alpha_U(P)]\). Then the bridge residual
\[
\mathcal B_{\mathrm{EW}}(P,N_\star)
:=
\alpha_U(P)\log(N_\star/\pi)-\frac{6\pi}{P}
=0.
\]
Equivalently, the bridge coordinate is
\[
N_{\mathrm{EW}}(P)
=
\pi\exp\!\left[\frac{6\pi}{P\alpha_U(P)}\right].
\]
For the measured-endpoint comparison coordinate \(P_C\),
\[
N_{\mathrm{EW}}(P_C)
=
3.5323546226929906511187512962330547600462\times10^{122}.
\]
For the independently evaluated source-forward coordinate \(P_{\rm fwd}\),
\[
N_0:=N_{\mathrm{EW}}(P_{\rm fwd})
=
3.532131543418936\ldots\times10^{122}.
\]
For either declared input \(P\), the bridge-map coordinate is conventionally labeled
\[
N_{\mathrm{CRC}}^{\mathrm{EW}}(P)
=
\pi\exp\!\left[\frac{6\pi}{P\alpha_U(P)}\right],
\]
with \(\mathcal B_{\mathrm{EW}}(P,N_{\mathrm{CRC}}^{\mathrm{EW}}(P))=0\), and hence
\[
\frac{v}{E_{\mathrm{cell}}}
=
\left(\frac{N_{\mathrm{CRC}}^{\mathrm{EW}}(P)}{\pi}\right)^{-P/12}.
\]
The exterior representation witness gives the exact weak-doublet multiplicity
\(3+1=4\), and every additive isomorphism-invariant load normalized by
\(L_P(\mathbb C)=P\) gives \(L_P(\mathbb C^4)=4P\). Identifying this
abstract weak load with the public screen load is exactly the common-carrier
receipt. The notation
\(N_{\mathrm{CRC}}^{\mathrm{EW}}(P)\) does not identify this mathematical coordinate with a cosmic
readback fixed point. That identification requires the source-derived direct
public-record producer and the common screen/electroweak load-carrier map. The
weighted Planck base-\(\Lambda\)CDM comparison coordinate
\(N_\Lambda=3.31292709806038\ldots\times10^{122}\) is about \(6.6\) percent
lower than \(N_0\).
\end{theorem}

\begin{proof}
Equating the screen and electroweak load expressions gives
\(P\log(N_\star/\pi)/12=\pi/[2\alpha_U(P)]\). Rearrangement gives
\(\mathcal B_{\mathrm{EW}}(P,N_\star)=0\). Substitution gives the displayed
\(N_{\mathrm{EW}}(P)\) and the hierarchy identity.
\end{proof}

\begin{remark}[Finite-presence and mean-count reserve branches]
\label{rem:n-reserve-branches}
On the protected-collar branch in Ref.~\cite{ophmicrophysics}, declared total
reserve expectation \(P/4\) and six-class equidistribution give finite
presence probability \(P/24\) for each declared class. If one class is
physically selected as the blocked event, its scalar-weighted presence receipt
is discharged, and a separately proved global attachment maps its normalized
collar-survival factor to the capacity coordinate, the result is
\[
N_{\rm pres}=N_0\left(1-\frac{P_{\rm fwd}}{24}\right).
\]
On the source-forward numerical branch this is
\(3.292097877326465\ldots\times10^{122}\), about \(0.63\) percent below the
weighted Planck base-\(\Lambda\)CDM comparison coordinate. The exponential expression
\[
N_{\rm Pois}=N_0e^{-P_{\rm fwd}/24}
\]
is a distinct mean-count or continuum branch and needs a carrier beyond the
finite projection trace. It gives
\(3.300072225377652\ldots\times10^{122}\), about \(0.39\)
percent below the same coordinate. The inherited screen/electroweak bridge
premises, common-load identity, physical seam action, one-class selection,
and scalar-weighted receipt do not provide a global action. Exact neutral and
multiplicative completions of the same local survival datum are positive and
obey disconnected-cut composition and finite cut-count regrouping, yet they
disagree after one cut. The finite source also selects no blocked-event
semantics. The declared branch therefore supplies no global-capacity object
for a horizon identification. A positive construction requires a stronger
source-derived global action. The target was known before these candidates were examined, so
neither comparison has predictive standing.
\end{remark}

\begin{theorem}[Icosahedral screen sieve and gated capacity composition]\label{thm:screen-sieve-capacity-projection}
On the certified echosahedral carrier lineage of
Ref.~\cite{ophmicrophysics}, assume twelve primitive equal-trace port-center
atoms, a declared integer atom-counting grammar on the total-twelve fiber,
the normalized central-readback Hilbert--Schmidt cost, oriented
\((12,30,20)\) incidence, and refinement-cocycle data. Then the counting
theorem gives twelve unit ports. The incidence theorem gives their
fixed-point-free inverse pairing, the proper \(A_5\) action, and the regular rank-three icosahedral
frame, naturally under refinement and relabeling. An equal-weight invariant
load \(X\) is therefore read locally as \(X/12\). For
\(X=\log(N/\pi)\), local cell entropy \(P/4\), and electroweak multiplicity
\(\beta_{\mathrm{EW}}=4\), define the screen composition coordinate
\[
\Gamma_{\mathrm{scr}}
=
\beta_{\mathrm{EW}}\frac{P}{4}\frac{1}{12}\log(N/\pi)
=
\frac{P}{12}\log(N/\pi).
\]
If, in addition, the common screen/electroweak load-carrier map identifies
\(\Gamma_{\mathrm{scr}}=\log(E_{\mathrm{cell}}/v)\) and matches it to the
source transmutation exponent \(\pi/(2\alpha_U(P))\), then
\[
\log(E_{\mathrm{cell}}/v)
=\frac{P}{12}\log(N/\pi),
\qquad
\alpha_U(P)^{-1}=\frac{P}{6\pi}\log(N/\pi),
\]
which is equivalent to the electroweak bridge residual
\(\mathcal B_{\mathrm{EW}}(P,N)=0\).
\end{theorem}

\begin{proof}
The declared counting-realization identity
\[
H(q)=12+\sum_p(q_p-1)^2
\]
gives the unique all-one allocation with exact gap two. The oriented
incidence has distance profile \((1,5,5,1)\), positive automorphism group
\(A_5\), and exact Gram matrix \(G^2=4G\) of rank three, as proved and checked
in Ref.~\cite{ophmicrophysics}. Edge-center collars expose the twelve primitive
atoms as central ports. Equal-weight additivity gives \(X/12\). Multiplying by \(P/4\) and
\(\beta_{\mathrm{EW}}=4\) gives
\(\Gamma_{\mathrm{scr}}=(P/12)\log(N/\pi)\). The additional identification
receipt and the source transmutation equation give the remaining displays.
\end{proof}

The all-one split and gap are exact within the declared integer counting and
normalized readback-cost realization. The integer loads also have an
operational finite realization. Append-only signed atomic events generate the
register values, and an admissible repair transfers one whole unit across a
seam with load difference \(d\ge2\). For
\[
V(q)=\sum_p q_p^2
\]
each such move gives the exact decrement
\(\Delta V=-2(d-1)<0\). Hence \(V\) proves termination. The minimum number of
admissible moves is a derived settling cost, natural under every declared
carrier rotation and refinement. A half-unit display rescales the readback
units of the same event graph; it does not implement half of an atomic event.

The load-square Lyapunov function \(V\), the seam quadratic, the normalized
readback cost \(H\), and the Hessian associated with the third axiom are
distinct objects. The positive readback scale remains a unit convention.
The proper-\(A_5\) result follows independently from oriented incidence on
the declared echosahedral lineage. The port/weak-carrier intertwiner,
normalized port-load trace equality, and common screen/electroweak
load-carrier map remain physical branch premises. The exact exterior
multiplicity alone does not discharge the physical load identity.

This theorem is local to one certified carrier lineage. It does not identify
that carrier with a primitive observer, construct the federation nerve, or
derive a global \(S^2\) support screen. Those claims require the independent
access-and-record, overlap, carrier-to-screen, and refinement receipts stated
upstream.

\begin{theorem}[EW bridge-map fixed point]\label{thm:ew-exact-capacity}
Define the log-capacity map
\[
\mathcal C_{\mathrm{EW}}(P,x)
=
(1-\lambda)x+\lambda\frac{6\pi}{P\alpha_U(P)},
\qquad 0<\lambda\le1.
\]
For fixed declared \(P\) and \(\lambda=1/2\), this is a contraction with
Lipschitz constant \(1/2\). Its unique fixed point is
\[
x_\star(P)=\frac{6\pi}{P\alpha_U(P)},\qquad
N_{\mathrm{CRC}}^{\mathrm{EW}}(P)=\pi e^{x_\star(P)}.
\]
Therefore
\[
\mathcal B_{\mathrm{EW}}(P,N_{\mathrm{CRC}}^{\mathrm{EW}}(P))=0.
\]
This theorem concerns the deliberately defined map \(\mathcal C_{\mathrm{EW}}\), not the physical
direct public-record map. Identifying the output with cosmic capacity
requires an additional source law selecting one physical fixed point, a
universe-level carrier attachment, and the common screen/electroweak
load-carrier map.
At \(P_{\rm fwd}\) its value is
\(N_0=3.532131543418936\ldots\times10^{122}\), about \(6.6\) percent above
the weighted Planck base-\(\Lambda\)CDM coordinate. Applying the finite
presence or Poisson reserve factor requires the distinct premises of
Remark~\ref{rem:n-reserve-branches}; neither factor follows from this
contraction.
\end{theorem}

\begin{proof}
The map has derivative \(1-\lambda\) in \(x\), hence is a contraction for
\(0<\lambda\le1\). Solving \(x=\mathcal C_{\mathrm{EW}}(P,x)\) gives the
displayed \(x_\star(P)\). Substitution into \(\mathcal B_{\mathrm{EW}}\) gives
zero.
\end{proof}

\begin{theorem}[Conditional RG/Higgs naturality identity]\label{thm:rg-higgs-naturality}
On a declared hierarchy branch satisfying the selected-surface premises, let \(Q_s\) be the source hierarchy quotient and
\(Q_H\) the Higgs/electroweak quotient. Define
\[
\rho_{sH}([x]_s)=
\left[
P(x),N(x),
\Theta(P(x),N(x)),
\Pi_{HT}F_{\mathrm{crit}}F_{\mathrm{EW}}(P(x),N(x)),0
\right]_H,
\]
with
\[
\Theta(P,N)=
\left(\frac{N}{\pi}\right)^{-P/12}.
\]
Here \(F_{\mathrm{EW}}\) denotes the declared electroweak map and
\(F_{\mathrm{crit}}\) denotes the criticality map.
For the corresponding normal-form maps and obstruction maps,
\[
\rho_{sH}n_s=n_H\rho_{sH},
\qquad
\chi_{sH}h_s=h_H\rho_{sH}.
\]
Consequently,
\[
\varepsilon_H
=
\max\{\varepsilon^n_{sH},\varepsilon^h_{sH}\}
=0,
\qquad
\varepsilon_H\in[0,0].
\]
\end{theorem}

\begin{proof}
The two routes through the source and Higgs quotients preserve the same tuple
\[
(P,N,\Theta(P,N),\Pi_{HT}F_{\mathrm{crit}}F_{\mathrm{EW}}(P,N)).
\]
A hidden scalar relevant deformation is not an independent
physical input on the selected branch: if absent from the retained finite constraint family, it
is off-branch by refinement stability; if retained, its value is fixed by the branch equations
and belongs to the displayed Higgs/top coordinate. Both controlled squares commute exactly, so
the product defect is zero.
\end{proof}

The conditional RG/Higgs naturality map does not itself use the measured weak scale. Measured \(W/Z\),
Higgs/top, \(G\), Planck-area, and \(\Lambda\) data are excluded by its declared input contract. The local/global package
records the direct-capacity/common-load electroweak projection bridge,
mathematical EW bridge-map fixed point, conditional operational-readback implications, and RG naturality
square. Conditional on the declared branch, this certifies the dimensionless electroweak hierarchy/naturality
identities; it does not supply the additional capacity-selecting source law,
the universe-level carrier attachment, or the common screen/electroweak
load-carrier identification. The strict
source-root certificate is open, and the result does not independently attach
\(E_\star\) to a physical GeV scale. The finite producer imported from the
synthesis paper is
\[
\mathfrak F_{r,\varepsilon}(D)=
\{M_\varepsilon(q):q\in\widetilde\Omega_{r,D}\},
\qquad M_0(q)=\alpha(G_q),
\]
with the public checkpoint packet, capacity-carrier representation,
whole-fiber scalarization, confusability-reflecting extension/refinement,
sewing, and fixed-point receipts kept explicit. The compatible all-rung
counterfamily completions have different zero sets under shared base,
positivity, carrier, and executable finite controls. The complete A1--A3
source contract has not been shown to contain those completions at every
rung, and no executable-to-Lean bridge establishes that membership. The
bounded completion class is nonidentifiable; the stronger source-class
verdict is open. Direct \(N\) is not evaluable on the incomplete capacity
source antecedent. A positive result requires a complete source antecedent and
one physical zero. The operational
resolution is an independent test. The conditional local carrier has twelve ports
and an oriented 24-slot register; the latter supplies no four-load conclusion.
The
mathematical local/global hierarchy-resonance package is a conditional
consistency surface. It does not construct the physical cosmic readback producer.
The hierarchy result does not use the
local/global resonance as an input.

\subsection{Higgs/top critical stage}

The Higgs/top critical stage is conditional on the electroweak gauge core. It
is described by a downstream split law together with a Jacobian readout map on
the declared running, matching, and threshold surface. Write the coupling
asymmetry as \(\rho_{\mathrm{EW}}\), distinct from the transmutation
coefficient \(b_{\mathrm{tr}}=4\). The candidate electroweak tuple
\[
(\eta_{\mathrm{source}},\rho_{\mathrm{EW}},\lambda_{EW},\tau_{2,\mathrm{tree}}^{\mathrm{exact}},\delta n_{\mathrm{tree}}^{\mathrm{exact}})
\]
emits the shared scalar
\[
\rho_{HT}=\log\!\bigl(1+\tau_{2,\mathrm{tree}}^{\mathrm{exact}}\bigr)
\]
and the declared residual formulas
\[
R_T=
-\tau_{2,\mathrm{tree}}^{\mathrm{exact}}\eta_{\mathrm{source}}^2
+\Bigl(1+\frac{\rho_{\mathrm{EW}}}{28}\Bigr)\eta_{\mathrm{source}}^6
+\frac{\eta_{\mathrm{source}}^8}{14}
+\frac{\eta_{\mathrm{source}}^9}{27},
\]
\[
R_H=
\eta_{\mathrm{source}}^5
-\frac{3}{25}\eta_{\mathrm{source}}^6
+\frac{\lambda_{EW}\eta_{\mathrm{source}}^6}{18}
+\frac{\eta_{\mathrm{source}}^8}{2\rho_{\mathrm{EW}}}.
\]
The split coordinates are
\[
\pi_y=
\frac{\eta_{\mathrm{source}}+\left(\frac32+\frac{\rho_{\mathrm{EW}}}4\right)\rho_{HT}+R_T}{\sqrt{\pi}},
\qquad
\pi_\lambda=
\frac{\eta_{\mathrm{source}}-\left(\frac43-\frac{\rho_{\mathrm{EW}}}{54}\right)\rho_{HT}+R_H}{\sqrt{\pi}},
\]
so the declared criticality Jacobian reads out
\[
\delta y_t(\mu_t)=\pi_y\,y_t^{\mathrm{core}}(\mu_t),
\qquad
\delta\lambda(\mu_t)=-\frac{16}{9}\pi_\lambda\,\lambda^{\mathrm{core}}(\mu_t).
\]
These equations define a physical readout only if a finite split-character
certificate supplies one strict source branch and a fixed pre-split carrier,
an exact two-coordinate quotient response, explicit characters deriving
\(\rho_{HT},R_T,R_H,\pi_y,\pi_\lambda\) and \(-16/9\), oriented mismatch
descent with canonical normalization and positivity, and an exhaustive
deformation quotient proving rigidity or a target-free positive-gap selector.
The displayed formulas do not supply that certificate.
This gives
\[
m_H = 125.1995304097179~\mathrm{GeV},
\qquad
m_t^{\mathrm{crit}} = 172.3523553288312~\mathrm{GeV}.
\]
This is a coordinate on the declared electroweak/criticality Jacobian surface,
obtained by
back-solving from the measured pair through the synchronization-scale scan, so
it is an exact interpolation of the PDG pair, a target-anchored fit; it
carries no certified Higgs or top complex pole. The
target-free headline is the double-criticality family evaluated below. The same surface emits a companion top coordinate.
An exact negative result bounds this branch from below: at fixed \(P\), with
\(u=1-P/24\), the exposed target-free reduct admits both the linear
(\(y_t=u\)) and Born (\(y_t=\sqrt u\)) probability-to-amplitude completions,
which agree on every reduct field and differ in the leading Higgs and top pole
ratios by \(2u(1-u)\) and \(u(1-u)/2\). The reduct is therefore not
\(P\)-complete for the Higgs/top readout. A physical conclusion requires an
explicit source object that selects the amplitude lift. The Standard Model gauge paper
states and proves this two-completion theorem.
The selected-class quark support witness carries a
target-anchored running-top comparison row using the PDG 2025 cross-section entry; it is not a separate
public source-only top prediction. The bridge to the auxiliary direct-top PDG row is closed as a
source-domain codomain no-go; the auxiliary row is compare-only.
The one-scalar companion seed
\[
\sigma_{\mathrm{crit}}=\frac{\alpha_U\cos(2\theta_{W0})}{\sqrt{\pi}}
\]
defines the fixed-ray companion branch with
\(\pi_y=\pi_\lambda=\sigma_{\mathrm{crit}}\).
Separately, the same criticality Jacobian admits a compare-only inverse slice that hits the canonical
Higgs/top reference pair exactly: it is an exact interpolation of the PDG pair, and the landing
is the defining condition of the slice rather than a derived result. That exact comparison is
useful for bookkeeping and does not define
the forward branch used for the public rows.
The Higgs boson belongs naturally in the boson discussion. The quark-family
chapter treats the top as the third-generation up-type quark. The top coordinate reported
here is the conditional criticality coordinate on this surface, not a separate public top-mass prediction
row. The selected-class quark support wrapper carries only a target-anchored comparison witness under
the declared source-only rule.

The local implication from the candidate electroweak tuple through the displayed criticality
formulas is executable and single-valued, and this arithmetic grade is
certified from raw interval data: the input box contains the eleven declared
branch inputs (the electroweak tuple and the criticality core/Jacobian
constants, each with units or dimensionless normalization and provenance),
the outward-rounded interval extension of every displayed
node, the Jacobian interval enclosure over the full box, and a non-singular
diagonal readout block with determinant bounded away from zero; the certified statement is scoped to
the declared surface and makes no criticality-system existence or uniqueness
claim. A physical Higgs mass requires a source root, independent scale, finite
quotient-path selection, same-scheme running, split rigidity, top and
threshold transport, a complex pole, an uncertainty enclosure, and
target-independent provenance. The exact inverse slice is a compare-only
validation surface.

\begin{table}[h]
\centering
\scriptsize
\begin{tabular}{@{}>{\raggedright\arraybackslash}p{0.13\textwidth}>{\raggedright\arraybackslash}p{0.21\textwidth}>{\raggedright\arraybackslash}p{0.34\textwidth}>{\raggedright\arraybackslash}p{0.21\textwidth}@{}}
\toprule
Particle & Calculation & Theorem scope & Value \\
\midrule
\(W\) boson & electroweak calibration & no source-only physical mass & no OPH-native pole \\
\(Z\) boson & electroweak calibration & no source-only physical mass & no OPH-native pole \\
Higgs boson & Higgs/top criticality & no source-only physical mass & no OPH-native pole \\
Top quark & Higgs/top criticality & companion coordinate on declared surface; direct-top auxiliary codomain no-go & not a separate source-only prediction \\
\bottomrule
\end{tabular}
\end{table}

\subsection{Source-closure and physical-pole envelope}\label{wzh-source-closure-envelope}

The preceding formulas end at running or declared-surface coordinates. A
physical \(W/Z/H\) theorem requires the following additional composition.

\begin{proposition}[Conditional physical clock attachment]
If the same strict source branch emits the dimensionless energy
\(\varepsilon_{\rm clk}>0\) of a declared clock transition whose physical
frequency \(\nu_{\rm clk}\) defines the operational unit, then, with
\(h_{\rm P}=2\pi\hbar\) Planck's constant,
\[
E_\star=\frac{h_{\rm P}\nu_{\rm clk}}{\varepsilon_{\rm clk}}.
\]
Equivalently, \(\gamma_\star=\ell_\star\nu_{\rm clk}/c\) gives
\(E_\star=\hbar\nu_{\rm clk}/\gamma_\star\). An operator-norm perturbation
of at most \(\epsilon\) changes an isolated selected clock gap by at most
\(2\epsilon\). This is a conversion and stability theorem; the required
physical reference transition and SI binding are absent.
\end{proposition}

\begin{proposition}[Exact finite source-Hamiltonian gap]
\label{prop:finite-source-hamiltonian-gap}
Let \(D_\sigma\) be the sign-twisted seam derivative on the connected
observer-visible finite local domain, and let
\(H_\sigma=D_\sigma^\ast D_\sigma\) use the declared unit-counting measure.
The signed-incidence formula makes \(H_\sigma\) integer symmetric and positive
semidefinite. A kernel vector would supply a sign-consistent coloring on each
connected component. The declared reversing seam transport frustrates every
triangle, so the connected complex has no such coloring and the signed
incidence rank is \(8{,}662\). Therefore \(\ker H_\sigma=0\). With maximum
degree \(12\), the determinant bound gives
\[
  \lambda_{\min}(H_\sigma)\geq 24^{-8661}>0 .
\]
A sparse numerical refinement gives
\(\lambda_{\min}=0.11753670336118684\) with relative residual
\(1.49\times10^{-14}\). The decimal is measured rather than certified. The
exact statement is strict positivity and the displayed floor.
\end{proposition}

This source-derived dimensionless gap does not select a physical field fiber
or laboratory transition. It therefore supplies no \(\nu_{\rm clk}\), SI
energy, particle mass, or determinant-line attachment. A sign-consistent
four-cycle control has an exact zero mode, so positivity is not hardwired into
the instrument~\cite{ophfpe}.

The atomic component of that packet has an exact scalar reduction. Let
\(P=P_3+P_4\) project onto the isolated zero-field \(6S_{1/2}\), \(I=7/2\)
cesium ground manifold, \(\dim V_3=7\), \(\dim V_4=9\), and let \(Q=1-P\).
On every real interval where \(Q(\widehat H-z)Q\) is invertible, define the
Feshbach--Schur operator
\(\mathcal F_P(z)=P(\widehat H-z)P-P\widehat HQ\,[Q(\widehat H-z)Q]^{-1}Q\widehat HP\).

\begin{proposition}[Cesium channel scalarization and monotone clock roots]
\(z\in\operatorname{spec}\widehat H\) exactly when
\(0\in\operatorname{spec}\mathcal F_P(z)\), with matching multiplicities.
Rotational invariance and the multiplicity-free decomposition
\(V_3\oplus V_4\) force \(\mathcal F_P(z)=f_3(z)P_3+f_4(z)P_4\), and
\(\mathrm d\mathcal F_P/\mathrm dz\preceq-P\) gives \(f_F'(z)\le-1\). One
sign bracket per channel therefore isolates exactly one level, and the clock
normal form is
\(H^{\mathrm{eff}}_{\mathrm{clock}}=E_0P+\widehat a_{\mathrm{Cs}}\,\mathbf I\cdot\mathbf J\)
with \(E_0=(9E_4+7E_3)/16\) and
\(\varepsilon_{\mathrm{Cs}}=E_4-E_3=4\widehat a_{\mathrm{Cs}}\).
\end{proposition}

The public atomic packet may therefore export two monotone scalar interval
evaluators with resolvent certificates and sign brackets in place of a
correlated 55-electron matrix. The open clock parents are the atomic-scheme
electromagnetic convention, the absolute electron ratio \(m_ec^2/E_\star\),
the cesium nuclear mass/spin/current/Compton packet, the two scalar
evaluators, summable refinement tails, and the no-target provenance freeze.
A non-entailment theorem fixes the direction of effort: two source extensions
that agree on every named object and differ in one unfixed parent
produce different unique gaps, so no unique \(\varepsilon_{\mathrm{Cs}}\) is a
theorem emitted by the source, and further evaluation of the existing
formulas cannot close the clock branch.

\begin{proposition}[Operational scale metrology]
In SI units
\(E/\mathrm{GeV}=[E/(h\nu_{\mathrm{Cs}})]\,[h\nu_{\mathrm{Cs}}/\mathrm{GeV}]\),
and the second factor is an exact unit conversion. A source-only numerical
GeV mass is therefore equivalent to a source-only dimensionless operational
clock ratio. Replacing cesium by another clock species adds the burden of
predicting that clock's frequency ratio to the defining cesium transition.
\end{proposition}

The stored clock-gap candidate reproduces the displayed gravitational
coupling to a relative agreement of order \(10^{-49}\) through the
clock--gravity identity, so the decimal is a calibration checksum rather than a prediction. The selection
mechanism for source laws themselves is closed separately by a source-action
rigidity theorem: on a finite quotient with a faithful reference measure, an
affinely independent feature basis, and a source-emitted moment vector, the
maximum-entropy law is unique, every feasible competitor pays a positive
relative-entropy gap with the Pinsker lower bound, and the surviving action
freedom consists of additive constants, exact feature redundancies, and
BRST-exact terms. The physical
Standard-Model operator basis and its moment vector
\(c_r(P_\star,N_\star)\) are the open inputs of that mechanism.

\begin{proposition}[Unique frozen RG and threshold transport]
Let the canonically normalized coefficient vector \(X(\mu)\) obey a locally
Lipschitz beta system on each interval of a frozen finite threshold list.
Assume deterministic matching maps, fixed loop order, field content,
renormalization scheme, threshold locations, and decoupling order. Then one
source initial condition determines one low-energy coefficient vector. If
\(L_j\) bounds the beta-system Lipschitz constant and \(K_j\) bounds threshold
map \(j\), perturbations obey the corresponding product of
\(e^{L_j\Delta t_j}\) and \(K_j\), plus the declared truncation and matching
remainders.
\end{proposition}

\begin{proof}
Picard--Lindel\"of gives existence and uniqueness between thresholds.
Gr\"onwall bounds perturbations on each interval; applying each deterministic
Lipschitz threshold map and inducting through the ordered list gives the
stated composition.
\end{proof}

The printed running/matching packet is a declared-convention contract. It does
not satisfy this source theorem because its beta provenance, threshold
origins, matching interval composition, and truncation enclosure are open.

\begin{theorem}[Conditional complex-pole theorem]
For \(B=W,Z,H\), let \(\Gamma_B^{\rm phys}(s,\xi)\) be the BRST-complete
physical inverse two-point block, including every declared mixing field after
the Ward/Slavnov--Taylor projection, analytically continued to a declared
Riemann sheet, and set
\[
D_B(s,\xi)=\det\Gamma_B^{\rm phys}(s,\xi).
\]
Assume a reference determinant has exactly one simple physical zero inside a
certified contour, no zero lies on the contour, and the full determinant obeys
\(|D_B-D_{B,0}|<|D_{B,0}|\) there. In the neutral sector also assume the
Ward-protected photon line has been separated from the massive eigenvalue.
Then the enclosed pole is unique and stable. With the exact convention
\[
s_B=\left(M_B-\frac{i}{2}\Gamma_B\right)^2,
\qquad M_B>0,\quad\Gamma_B\ge0,
\]
it determines one mass and width. If a Nielsen identity has the form
\(\partial_\xi D_B=C_B^{(\xi)}D_B\), the simple pole is gauge-parameter
independent~\cite{gambino1999,grassiKniehlSirlin2001}.
For \(W\) and \(Z\), interpretation as a physical resonance additionally requires
a dressed BRST-invariant current amplitude whose Laurent coefficient at that
pole is nonzero; for \(H\), the analogous scalar-amplitude coupling is a
separate condition. A simple zero of the inverse-propagator determinant alone
does not establish either coupling.
\end{theorem}

\begin{proof}
Rouch\'e's theorem preserves the number of enclosed zeros under the certified
perturbation. Simplicity gives the implicit-function stability bound and a
locally simple determinant zero. Differentiating \(D_B(s_B(\xi),\xi)=0\) and using the Nielsen
identity gives \(ds_B/d\xi=0\). An invertible analytic field redefinition
multiplies \(D_B\) by a nonzero analytic factor and therefore preserves the
pole set.
\end{proof}

No source-emitted \(W/Z/H\) self-energy kernels, analytic-sheet receipt,
contour enclosures, widths, residues, or pole-convention uncertainty map are
present on the displayed mass-chart surface. Schema fields named
for the \(W\), \(Z\), and top pole masses are labels only and do not
supply this theorem.

\subsection*{Physical scope of the icosahedral Standard Model and \(W/Z\) results}

The finite icosahedral package is an exact conditional recognition result.
Incidence determines the antipode \(J\). Under the explicit contract that an
admissible response is a signed central involutive graph automorphism
implementing inverse-port readback, the responses are exactly \(\pm J\);
their common sign is conventional. Conditional also on the rank-15 matter
contract with its unique charge-conjugate projector pair, the port-current,
anomaly, Spin, and tensor-descent certificates fix the \(3+2\) block
structure, hypercharge lattice, and \(\mathbb Z_6\) quotient from the stated
premises alone. The finite source packet therefore selects the odd-Weyl Spin
typing and the maximal faithful \(\mathbb Z_6\) quotient on its declared
carrier. The kernel is computed on every realized tensor and is insensitive
to the projector representative. This finite selection supplies no laboratory
identification of its current or flux sectors and no physical seam action.
Equality with the independently reconstructed Tannaka current, attachment of
the band to three physical chiral families, exclusion of extra light sectors,
four-dimensional topological attachment, scalar attachment and dynamics, and
construction of a chiral quantum field theory remain work in progress.

The conditional field-theory implications are explicit. A finite local action
gives an exact finite gauge-invariance and locality theorem, and the familiar
electroweak tree kernel is conditional on a separate canonical continuum and
action-normalization bridge. An exact finite measure criterion and an exact
finite Hamiltonian criterion are two parallel branches over that action. A
separate formal perturbative branch carries the strict finite-order \(W/Z\)
pole theorem. A nonperturbative continuum completion gives an observable-sector
reconstruction implication and a distinct continued-sheet resonance-stability
implication. The measure and perturbative branches are parallel descendants of
the finite local action. Neither implies the other, and a perturbative pole
does not imply the continuum completion.

These quantum field theory (QFT) theorems state what follows from typed action
and quantization packets. They do not show that the target-free source emits
those packets. The source-selected action and normalization, complete measure,
target-clean perturbative matching, physical current amplitudes, source law
and covariance, uncertainty enclosure, operational clock, and continuum tower
with its continued-sheet packet are work in progress. Two external Standard
Model calculation engines provide a bounded validation surface for the
perturbative matching protocol. They are validation inputs rather than OPH
source producers. No source-native dimensionless or physical-unit \(W/Z\)
pole follows from them.


\subsection*{Exact source-side frontier}

The finite fifteen-state matter representation fixes the one-copy quadratic
indices
\[
\sum_{\rm Weyl}T_{\SU(3)}=2,\qquad
\sum_{\rm Weyl}T_{\SU(2)}=2,\qquad
\sum_{\rm Weyl}d_3d_2Y^2=\frac{10}{3}.
\]
After importing the standard four-dimensional one-loop functional in the
convention
\[
\frac{\mathrm d g_i}{\mathrm d\ln\mu}
=\frac{b_i g_i^3}{16\pi^2},
\]
these indices give
\[
b_Y=\frac{20}{9}N_g+\frac16N_H,\qquad
b_2=-\frac{22}{3}+\frac43N_g+\frac16N_H,\qquad
b_3=-11+\frac43N_g.
\]
Here \(N_g\) counts matter families and \(N_H\) counts scalar doublets. The
declared choice \((N_g,N_H)=(3,1)\) gives
\((b_Y,b_2,b_3)=(41/6,-19/6,-7)\). The finite source model does not select
that physical multiplicity pair.

On the declared one-doublet, rank-three family branch, four-dimensional power
counting classifies the renormalizable scalar and Yukawa coefficient space as
\[
\mathbb R_{>0}\times\mathbb R^2\times
\operatorname{Mat}_3(\mathbb C)^3.
\]
The three Yukawa matrices have complex dimension \(27\). This is an exact
operator-basis theorem, with no selected coefficient, vacuum coordinate, or
map from the source chart to a renormalized vacuum expectation value. A
bounded check of the selected charged-response and pole-residue simulator
pair finds no explicit scalar-completion input in those two artifacts. That
check is not an exhaustive producer classification and does not prove
physical scalar non-identifiability.

The invariant screen order-unit line and the central order-unit line of the
four-copy weak multiplicity algebra are each one-dimensional. There is a
unique positive unital normalized-trace isomorphism between them, natural
under the declared refinements. This abstract line isomorphism does not
identify a physical common load, choose the scalar sector, or normalize a
laboratory current.

A source clock would convert a dimensionless gap
\(\epsilon_{\rm clk}=\Delta E_{\rm clk}/E_\star\) through
\[
\Delta E_{\rm clk}=h\nu_{\rm clk},\qquad
E_\star=\frac{h\nu_{\rm clk}}{\epsilon_{\rm clk}}.
\]
For \(0<\epsilon_{\rm lo}\le\epsilon_{\rm clk}\le\epsilon_{\rm hi}\), the
exact interval is
\[
\frac{h\nu_{\rm clk}}{\epsilon_{\rm hi}}
\le E_\star\le
\frac{h\nu_{\rm clk}}{\epsilon_{\rm lo}}.
\]
The interval inversion is exact. Proposition
\ref{prop:finite-source-hamiltonian-gap} supplies a source-derived
dimensionless spectral gap on the finite local domain. It supplies neither a
physical reference transition nor its frequency, so the SI interval has no
source value.

\subsubsection*{Quantization steps and the \(W/Z\) landing}

The table below names the Standard-Model field-theory steps used in this
section, together with what each one implies, what its construction requires,
and its present status. The labels \(\mathrm{QFT}\text{-}\mathrm{Q0}\) through
\(\mathrm{QFT}\text{-}\mathrm{Q4}\) are local to this section and stay distinct
from the particle-receipt clauses \((Q1)\)--\((Q3)\). Their dependency
structure is a graph rather than a single ordered ladder:
\begingroup\small
\[
\begin{array}{rcl}
\mathsf{declared\ finite\ action}&\longrightarrow&\mathsf{QFT\!-\!Q1}
 \longrightarrow
 \left\{\begin{array}{l}
 \mathsf{QFT\!-\!Q2E/QFT\!-\!Q2H},\\
 \mathsf{QFT\!-\!Q3\ BV/ST}\longrightarrow
 \mathsf{strict\ finite\!-\!order\ }W/Z;
 \end{array}\right.\\[3pt]
\mathsf{nonperturbative\ observable\ tower}&\longrightarrow&
\mathsf{QFT\!-\!Q4\ OS}\\
&&\longrightarrow\mathsf{QFT\!-\!Q4\ resonance}.
\end{array}
\]
\endgroup
Native provenance on QFT-Q1 additionally requires QFT-Q0 and a
source/action-identity receipt.  No QFT-Q2-to-QFT-Q3,
QFT-Q3-to-QFT-Q2, or QFT-Q3-\(W/Z\)-to-QFT-Q4 arrow is automatic.

\begin{center}
\small
\begin{tabular}{@{}p{0.10\textwidth}p{0.31\textwidth}p{0.25\textwidth}p{0.24\textwidth}@{}}
\toprule
Tier & Mathematical implication & Required construction & Present status \\
\midrule
QFT-Q0 & finite representation, charge, anomaly, lattice, and selector consequences on the declared packet
& physical source selection and attachment
& conditional finite pass; producer incomplete \\
QFT-Q1 & finite local classical \(G_6\) action is gauge invariant and local; the standard tree kernel also needs canonical continuum normalization
& source-selected action, coefficients, regulator, normalization, and ancestry
& implication specified; producer open \\
QFT-Q2-E/H & an equivariant determinant-line section or a noncollapsing constrained Hamiltonian supplies an exact finite quantum object
& full operator, measure/current or Hamiltonian/nonvacuum packet, and refinement controls
& criteria specified; constructions open \\
QFT-Q3 & stable anomaly-free BV/ST theory is formally restorable order by order; strict \(W/Z\) poles land here
& counterterm basis, matching and Faddeev--Jackiw (FJ) engines, identities, currents, and numerical freeze
& implication specified; imported validation possible; OPH producer open \\
QFT-Q4 & OS reconstruction and a separate continued-sheet resonance theorem
& reflection-positive tower and analytic-continuation packet
& implications specified; constructions open \\
\bottomrule
\end{tabular}
\end{center}

\begin{theorem}[Finite local classical \(G_6\) action at QFT-Q1]
\label{thm:qft-q1-finite-local-g6}
Let \(K_r\) be a finite oriented spin four-complex with bounded incidence,
positive cell weights, declared boundary conditions, paired edge orientations,
and declared spin transports. Put \(U_e\in
G_6=S(U(3)\times U(2))\) on oriented edges and the declared Higgs and
left-handed matter variables on vertices. Require every matter action to
descend to a well-defined representation of \(G_6\), every plaquette term to be a
declared class function of its \(G_6\) holonomy, and every finite difference
to be gauge covariant with defined endpoint data. Then a finite sum of these
plaquette, Higgs, fermion, and invariant Yukawa terms is local and exactly
gauge invariant.
\end{theorem}

\begin{proof}
Plaquette holonomies transform by conjugation, covariant edge differences
transform at their endpoints, and class functions and invariant contractions
remove those transformations. In integer hypercharge normalization the three
Yukawa sums are
\[
1+3-4=0,\qquad 1-3+2=0,\qquad -3-3+6=0.
\]
Every term has bounded cell support.
\end{proof}

\begin{corollary}[Canonically normalized electroweak tree kernel]
\label{cor:qft-q1-ew-tree-kernel}
If, in addition, the long-wavelength map sends the finite Higgs kinetic form
to \((D_\mu H)^\dagger D^\mu H\) with canonical generator normalization and
broken background \(H_0=2^{-1/2}(0,v)^T\), then
\[
w=\frac{g^2v^2}{4},\qquad
\mathcal M_N^2=\frac{v^2}{4}
\begin{pmatrix}g^2&-gg'\\-gg'&g'^2\end{pmatrix},
\qquad z=\frac{(g^2+g'^2)v^2}{4},
\]
and the other neutral eigenvalue is zero.
\end{corollary}

\begin{remark}[QFT-Q1 boundary]
This is a classical existence template after the complex, fields,
coefficients, boundary data, and normalization bridge are supplied. It does
not show that OPH selects them, construct a chiral measure, remove mirrors or
doublers, or produce a complex pole.
\end{remark}

\begin{theorem}[QFT-Q2-E equivariant determinant-line criterion]
\label{thm:qft-q2e-determinant-line}
At a fixed finite stage, let \(D_r(U)\) be a local gauge-covariant,
\(\gamma_5\)-Hermitian Ginsparg--Wilson operator on a connected admissible
gauge-field component on which the chiral-projector rank is constant, with
the declared spectral gap and locality bounds. A basis-independent,
gauge-invariant finite chiral measure exists precisely when the determinant
line admits a nowhere-zero gauge-equivariant section in the required locality
and smoothness class. For a nonfree gauge action this is an equivariant
statement over the action groupoid, including stabilizer actions. In local
connection form the measure current must reproduce the projector curvature
and obey global loop integrability. Flatness with trivial holonomy is only a
sufficient special case.
\end{theorem}

\begin{proof}
Changes of Weyl basis are transition functions of the determinant line. A
nowhere-zero equivariant section cancels them and descends to the gauge
quotient; conversely a basis-independent gauge-invariant phase supplies
compatible nonzero fiber vectors. The curvature equation and loop condition
are the local and global integrability conditions for that section.
\end{proof}

\begin{theorem}[QFT-Q2-H finite Hamiltonian soundness]
\label{thm:qft-q2h-hamiltonian-soundness}
Let \(\mathcal H_{\rm kin}\) be finite dimensional and let local Gauss
generators exponentiate to the complete local \(G_6\) action. Define
\[
C_G=\sum_{x,a}(G_x^a)^\dagger G_x^a,\qquad
P_{\rm phys}=\mathbf1_{\{0\}}(C_G),
\qquad
1<\operatorname{rank}P_{\rm phys}<\dim\mathcal H_{\rm kin}.
\]
Suppose a bounded-range self-adjoint \(H_r\) commutes with this action and has
a unique certified physical ground state \(\Omega_r\) with a positive gap.
If a bounded self-adjoint gauge-invariant observable \(O_r\) has strictly
positive variance in \(\Omega_r\), and the packet also supplies its claimed
chiral index, positive mirror gap, primitive completeness, and
complement-complete refinement controls, then
\(P_{\rm phys}\mathcal H_{\rm kin}\) is a nonvacuous finite unitary gauge
theory with a positive-energy nonvacuum physical excitation and the declared
chiral/mirror properties.
\end{theorem}

\begin{proof}
The vector
\[
(O_r-\langle\Omega_r,O_r\Omega_r\rangle)\Omega_r
\]
is physical, orthogonal to the ground state, and nonzero by the variance
condition. The remaining conclusions are exactly the separately supplied
index, gap, completeness, and refinement clauses. Thus neither an identity
projector nor a vacuum-only physical sector passes.
\end{proof}

\begin{remark}[QFT-Q2 boundary]
The QFT-Q2-E and QFT-Q2-H results are criteria and soundness theorems. The
current OPH corpus supplies neither full-\(G_6\) construction. Anomaly
arithmetic alone instantiates neither theorem.
\end{remark}

\begin{theorem}[Formal QFT-Q3 Slavnov--Taylor restoration]
\label{thm:qft-q3-formal-st-restoration}
Let \(S_0\) be the complete canonically normalized Standard-Model BV action,
including the gauge-fixing and antifield sectors, and suppose
\((S_0,S_0)=0\). Fix a regulator/scheme satisfying the quantum action
principle, locality, and the declared power counting. Assume stability under
renormalization in a complete permitted counterterm basis with fixed
normalization conditions, classification of the applicable local
ghost-number-one BRST cohomology, vanishing of the declared perturbative
anomaly class, and a separate check of all applicable global anomalies. Then
finite local counterterms may be chosen recursively so that the formal series
\[
\Gamma=S_0+\sum_{n\ge1}\kappa^n\Gamma_n,\qquad
\kappa=(16\pi^2)^{-1},
\]
satisfies the renormalized Slavnov--Taylor identity order by order.
\end{theorem}

\begin{proof}
At order \(n\), the quantum action principle makes the breaking a local
ghost-number-one functional \(\Delta_n\). Wess--Zumino consistency makes it
closed under the linearized Slavnov--Taylor operator. The cohomology
classification splits it into an anomaly representative plus an exact term.
Anomaly clearance removes the first, and an allowed finite counterterm
cancels the second. Stability and the normalization conditions fix the
remaining invariant freedom, so induction proves the formal statement.
\end{proof}

\begin{remark}[QFT-Q3 boundary]
This is a formal power-series theorem, not a convergence theorem, QFT-Q2
construction, or QFT-Q4 Wightman construction. A numerical implementation
must produce its regulator-specific restoration transcript and verify
the Ward, Slavnov--Taylor, and Nielsen identities.
\end{remark}

\begin{lemma}[Conditional first-order FJ coordinate change]
\label{lem:qft-q3-fj-reparametrization}
Freeze the potential normalization, bare VEV shift, tadpole prescription,
field and parameter counterterms, mass arguments, mixing coordinates, and
gauge-fixing convention. If the complete finite change is
\[
p_L=p_F+\kappa\,\delta p^{(1)}+O(\kappa^2)
\]
and \(s(p)=s_0(p)+\kappa s_1(p)+O(\kappa^2)\), equality of the exact pole in
the two coordinates implies
\[
s_{1,F}=s_{1,L}+\delta p^{a(1)}\partial_as_0.
\]
\end{lemma}

\begin{proof}
Substitute \(p_L(p_F)\) and Taylor expand through first order. The sum must
include every transformed parameter, normalization, mass argument,
counterterm, and mixing coordinate; a VEV-only substitution is insufficient.
\end{proof}

\begin{theorem}[Strict charged and neutral pole coefficients]
\label{thm:qft-q3-strict-wz-coefficients}
On one scheme, contribution mask, and resonance sheet, write
\[
\Gamma_W^T=s-w+\kappa\Pi_{WW}^{(1)}
+\kappa^2\Pi_{WW}^{(2)}+O(\kappa^3).
\]
For \(s_W=w+\kappa s_{W,1}+\kappa^2s_{W,2}+O(\kappa^3)\),
\[
s_{W,1}=-\Pi_{WW}^{(1)}(w),\qquad
s_{W,2}=\Pi_{WW}^{(1)}(w)\Pi_{WW}^{(1)\prime}(w)
-\Pi_{WW}^{(2)}(w).
\]
For the massive root of the full photon--\(Z\) matrix with tree value
\(z\ne0\),
\[
s_{Z,1}=-\Pi_{ZZ}^{(1)}(z),
\]
\[
s_{Z,2}=\Pi_{ZZ}^{(1)}(z)\Pi_{ZZ}^{(1)\prime}(z)
-\Pi_{ZZ}^{(2)}(z)
+\frac{\Pi_{ZA}^{(1)}(z)\Pi_{AZ}^{(1)}(z)}{z}.
\]
Thus the one-loop-squared neutral mixing product is excluded at strict one
loop and is one mandatory term of the complete strict-two-loop mask.
\end{theorem}

\begin{proof}
Insert the root series and compare powers of \(\kappa\). In the neutral
sector use the Schur complement of the photon block; both off-diagonal
entries begin at order \(\kappa\).
\end{proof}

\begin{theorem}[Nielsen control and physical current pole at QFT-Q3]
\label{thm:qft-q3-nielsen-current-pole}
Suppose the inverse matrix and Nielsen insertions are holomorphic near a
simple massive root on the frozen sheet and, through retained order \(N\),
\[
\partial_\eta\Gamma^T
=\Lambda_\eta\Gamma^T+\Gamma^T\widetilde\Lambda_\eta
+O(\kappa^{N+1}).
\]
Then
\[
\partial_\eta\det\Gamma^T
=\operatorname{tr}(\Lambda_\eta+\widetilde\Lambda_\eta)
\det\Gamma^T+O(\kappa^{N+1}),
\qquad
\partial_\eta s_p=O(\kappa^{N+1}).
\]
If the simple left/right kernel vectors \(\ell,r\) have
\(\ell^\dagger\Gamma^{T\prime}(s_p)r\ne0\), and the dressed renormalized
BRST-invariant current vertices obey
\(J_L^\dagger r\ne0\) and \(\ell^\dagger J_R\ne0\), then the
gauge-invariant current amplitude has Laurent residue
\[
\frac{(J_L^\dagger r)(\ell^\dagger J_R)}
{\ell^\dagger\Gamma^{T\prime}(s_p)r}.
\]
\end{theorem}

\begin{proof}
Use
\(\partial_\eta\det\Gamma
=\operatorname{tr}(\operatorname{adj}\Gamma\,\partial_\eta\Gamma)\);
the adjugate identity remains valid at a singular matrix and avoids dividing
by the determinant. The simple-root implicit equation gives the
omitted-order gauge variation. The rank-one Laurent expansion of
\(\Gamma^{-1}\), contracted with the dressed current vertices, gives the
amplitude pole. This is not a positivity claim for an unstable elementary
field.
\end{proof}

\begin{theorem}[Gauge-invariant QFT-Q4 reconstruction]
\label{thm:qft-q4-gauge-invariant-os}
A compatible cofinal Schwinger family for a complete declared
gauge-invariant observable algebra that satisfies distributional convergence,
Euclidean covariance, graded symmetry/locality, reflection positivity,
clustering, growth/regularity, noncollapse, and refinement Cauchy control
reconstructs a positive Hilbert space, cyclic vacuum, positive-energy
translations, and the corresponding observable-sector Wightman distributions
and local graded net.
\end{theorem}

\begin{remark}
This observable-sector implication does not by itself construct colored local
fields, charged infrared sectors, confinement, asymptotic completeness, an
\(S\)-matrix, or a second-sheet resonance.
\end{remark}

\begin{theorem}[QFT-Q4 resonance and residue stability]
\label{thm:qft-q4-resonance-stability}
On one common continued sheet, write
\[
G_r(s)=\frac{N_r(s)}{D_r(s)}+G_{r,\rm reg}(s).
\]
Let a Jordan contour and its interior lie in a common domain on which
\(D_r,N_r,G_{r,\rm reg}\) and their limits are holomorphic. Assume locally
uniform convergence of these data and \(D_r'\), one simple enclosed zero
\(s_r\), uniform nonzero contour and derivative bounds, and the Rouch\'e
inequality
\[
\sup_C|D-D_r|<\inf_C|D_r|.
\]
Then the limiting denominator has one simple zero \(s_*\), \(s_r\to s_*\),
and, if \(N(s_*)\ne0\),
\[
\operatorname*{Res}_{s=s_r}G_r(s)
=\frac{N_r(s_r)}{D_r'(s_r)}
\longrightarrow
\frac{N(s_*)}{D'(s_*)}\ne0.
\]
\end{theorem}

\begin{proof}
Rouch\'e preserves the zero count for the holomorphic denominators.
Compactness and uniqueness give root convergence. Uniform numerator and
derivative convergence with the lower derivative bound gives residue
convergence. Ordinary uniform convergence of the meromorphic quotient through
its own poles is neither assumed nor valid.
\end{proof}

\begin{remark}[Producer boundary]
These results specify conditional implications only. The current corpus does
not instantiate the source-selected QFT-Q1 action, either full QFT-Q2 object,
the QFT-Q3 matching/FJ/two-engine/current and uncertainty packet, or the
QFT-Q4 tower and continued-sheet data.
\end{remark}


\subsubsection*{Strict-one-loop W/Z pole-map kernel}

The finite-order theory map is explicit and machine checked.  Let the
renormalized one-doublet electroweak input at scale \(Q\) be
\(\theta(Q)=(g,g',v_F,\ldots)\), with canonical Higgs kinetic term and
\(v_F>0\), and set
\[
w=\frac{g^2v_F^2}{4},\qquad
z=\frac{(g^2+g'^2)v_F^2}{4}.
\]
Use the inverse-propagator convention
\[
\Gamma^T(s)=s-m_0^2-\Delta^T(s)
=s-m_0^2+\Pi^T(s),\qquad \Delta^T=-\Pi^T,
\]
where each \(\Delta^{(1)}\) includes its one-loop factor,
counterterms, tadpoles, and the complete declared strict-one-loop mask.
Relative to the coefficient convention of
Theorem~\ref{thm:qft-q3-strict-wz-coefficients},
\[
\Delta_{ij}^{(1)}(s)=-\kappa\Pi_{ij}^{(1)}(s),
\qquad \kappa=(16\pi^2)^{-1}.
\]
The two notations are translations of one pole equation, not independent
results.

\begin{proposition}[Strict-one-loop charged and neutral pole map]
\label{prop:strict-one-loop-wz-map}
Assume the tree roots \(w,z>0\) are simple and the declared one-loop entries
are holomorphic near them on the frozen analytic sheet.  Then
\[
s_W^{[1]}=w+\Delta_{WW}^{(1)}(w),\qquad
s_Z^{[1]}=z+\Delta_{ZZ}^{(1)}(z).
\]
In the neutral tree-level photon--\(Z\) basis,
\[
\Gamma_N^T(s)=
\begin{pmatrix}
s-\Delta_{AA}^{(1)}(s)&-\Delta_{AZ}^{(1)}(s)\\
-\Delta_{ZA}^{(1)}(s)&s-z-\Delta_{ZZ}^{(1)}(s)
\end{pmatrix}+O(\epsilon^2).
\]
The product
\(\Delta_{ZA}^{(1)}\Delta_{AZ}^{(1)}\) has loop power two and is excluded
from a strict-one-loop root.  Its leading Schur-complement contribution is
\[
-\frac{\Delta_{ZA}^{(1)}(s)\Delta_{AZ}^{(1)}(s)}{s},
\]
which belongs only in a separately complete two-loop map together with the
genuine two-loop entries and pole-iteration derivatives.
\end{proposition}

\begin{proof}
Write \(s_W=w+\epsilon\sigma_W+O(\epsilon^2)\) in the charged inverse
entry.  Its order-\(\epsilon\) coefficient is
\(\sigma_W-\delta_{WW}^{(1)}(w)\).  For the neutral determinant, write
\(s=z+\epsilon\sigma_Z+O(\epsilon^2)\).  The order-\(\epsilon\) coefficient
is \(z[\sigma_Z-\delta_{ZZ}^{(1)}(z)]\); both off-diagonal entries start at
order \(\epsilon\), so their product starts at order \(\epsilon^2\).
\end{proof}

Put \(S=gv_F/2\), \(t=g'/g\),
\(d_W=\Delta_{WW}^{(1)}(w)/w\), and
\(d_Z=\Delta_{ZZ}^{(1)}(z)/z\). The strict scalar consumer then obeys
\[
s_W=S^2(1+d_W),\qquad
s_Z=S^2(1+t^2)(1+d_Z).
\]
Consequently, whenever the denominator is nonzero,
\[
\frac{s_W}{s_Z}
=\frac{1+d_W}{(1+t^2)(1+d_Z)}.
\]
This is the exact quotient of the two one-loop-truncated pole coordinates.
Re-expansion to strict one-loop order gives
\[
\frac{s_W}{s_Z}
=\frac{1}{1+t^2}\bigl(1+d_W-d_Z\bigr)+O(\epsilon^2).
\]
The explicit common scale cancels under a passive rescaling at fixed
normalized corrections. An active change of \(v_F\) can move thresholds and
therefore \(d_W,d_Z\). The identity leaves \(t,d_W,d_Z\) unselected and gives
no numerical pole ratio by itself. Exact counterfamilies in the factored
consumer coordinates show that these normalized inputs can change the
dimensionless mass and width readouts while the common scale remains
irrelevant.

For \(s_V=m_{V,0}^2+\Delta_V^{(1)}\), the strict energy-pole coefficients are
\[
\delta M_V^{(1)}=\frac{\operatorname{Re}\Delta_V^{(1)}}{2m_{V,0}},
\qquad
\Gamma_V^{(1)}=-\frac{\operatorname{Im}\Delta_V^{(1)}}{m_{V,0}}.
\]
They are distinct from the exact coordinate transform of the truncated
complex number.  On the lower-half-plane branch,
\[
M_V=\sqrt{\frac{|s_V|+\operatorname{Re}s_V}{2}},\qquad
\Gamma_V=\sqrt{2\bigl(|s_V|-\operatorname{Re}s_V\bigr)},
\qquad s_V=(M_V-i\Gamma_V/2)^2.
\]
Applying this nonlinear square root exactly resums kinematic powers of the
one-loop coefficient.  It is a useful display coordinate, not a strict
two-loop calculation and not the object to compare in a finite-order Nielsen
test.

\begin{proposition}[Evidence cannot self-attest]
\label{prop:wz-evidence-nonself-attestation}
An untrusted input boolean asserting an external Fleischer--Jegerlehner, matching, source-law,
gauge/BRST, clock, or ancestry property cannot certify that property.  A
physical-claim verifier must resolve an independent hash-bound witness, validate
it, and bind it to the exact numerical subject, order, mask, scheme, and
analytic sheet.
\end{proposition}

\begin{proof}
Choose a subject for which the external property is false and set the
untrusted boolean to true while preserving every relation recomputed by the
verifier.  If the verifier resolves no independent witness, it follows the
same accepting path.  Hence acceptance would admit a false instance.
\end{proof}

The fail-closed consumer receipt implements these rules and rejects
self-attested physical flags, unrelated-but-self-consistent poles, substituted empty
fixtures, inflated tolerances, corrupted redundant fields, and altered
neutral diagnostics. Its archived Standard Model dimensional-regularization
order-one fixture at \(Q=160\) GeV
evaluates to
\[
\begin{aligned}
s_W^{[1]}&=(6459.842027569383-160.532752773045i)\;\mathrm{GeV}^2,\\
s_Z^{[1]}&=(8222.835212344102-218.292761806439i)\;\mathrm{GeV}^2.
\end{aligned}
\]
The corresponding strict readouts are
\((M_W,\Gamma_W)=(80.374161202712,2.007425074735)\) GeV and
\((M_Z,\Gamma_Z)=(90.680036075608,2.402420059845)\) GeV.  These numbers
reconstruct an archived backend row. They carry target ancestry and supply no
independent self-energy evaluation, no source-selected vacuum normalization, no
complete neutral matrix, no source covariance, no independent gauge or
Becchi--Rouet--Stora--Tyutin receipt, and no source clock. The strict one-loop
pole map is therefore conditional, and it is not source-native physical, so no
physical conclusion follows. What is proved is the implication from a
complete declared renormalized strict-one-loop packet to the separated pole and
mass and width readouts. The construction of that antecedent from the source is
work in progress.

A separate target-free shared-\(\xi\) Standard Model harness binds a declared
action and active census to a frozen effective-field-theory interval. It
contains two independent rule tables, direct and converted
Fleischer--Jegerlehner blocks, and payload-replayed Ward and
Slavnov--Taylor controls. The rule tables are coordinate variants inside one
harness rather than two independently implemented raw loop engines. Separate
interval receipts exclude scalar zeros on declared principal-sheet boxes and
isolate, for each of \(W\) and \(Z\), one simple scalar zero with derivative
and scalar-residue balls in its declared lower-half pole box on a
channel-specific algebraic chart. They identify neither chart with the
physical resonance sheet, prove no unique continuation from the principal
chart, and do not supply the sign bridge, full-matrix
Laurent data, dressed-current amplitude, or independent numerical replay.
Proof-bearing physical validation also requires the second engine,
counterterms generated from the complete action, BRST-generated ghost
vertices, finite symmetry restoration, symbolic Nielsen identities, an
artifact-resolving third verifier, a certified continuation identity,
full-matrix contour and Laurent data, dressed-current residues, and complete
independent replay. The external fixture is not composed with the OPH
electroweak chart. The harness therefore establishes no physical \(W/Z\) pole
and supplies no OPH-native source packet.


\begin{theorem}[Rigidity, dependency, and W/Z/H composition criterion]
Let \(\mathcal F_{\mathrm{EW}}\) be the fully enumerated class of source maps
that preserve the declared locality, symmetry, dimensional, refinement, and
provenance constraints, modulo gauge and trivial reparameterization. Suppose
either the pole readout is constant on this class or a target-independent
selector has one minimizer with a positive winner gap. Suppose further that
a hash-bound acyclic dependency graph records the source objects, selectors,
conventions, uncertainty model, and outputs, and measured \(W/Z/H\) data and
calibrated proxies are absent from every source ancestor set.

If, in addition, a unique source root and independently source-closed physical
\(E_\star\) emit the canonical electroweak and criticality coefficients, the
finite quotient-path and split-character certificates hold, and the declared
renormalization-group and complex-pole hypotheses hold, then OPH determines one source-separated pole
triple \((s_W,s_Z,s_H)\) with a certified uncertainty enclosure.
\end{theorem}

\begin{proof}
Rigidity or the positive-gap selector removes alternative source maps. A
deterministic leaf depends only on its ancestors, so the dependency graph
gives formal source separation. The unique root and scale emit one coefficient
packet, declared running and matching transport it uniquely, and the pole theorem
gives one stable pole per particle. Composition proves the claim.
\end{proof}

This criterion is not satisfied by the displayed numbers. The finite
quotient-path and split-character certificates are assumptions, the public
and source-only pixel branches differ, and the physical scale,
renormalization-group transport, complex pole, and target-independence proof
are open.

The missing object is a source-emitted certificate.
The structural theory admits target-free analytic counterfamilies
\(\tau_2=-c\eta^2\), \(\delta n=d(1-\rho_{\rm EW})\eta^2\) and
\((\pi_y,\pi_\lambda)\mapsto(\pi_y+a\eta^N,
\pi_\lambda+b\eta^N)\), all on the same open physical domain but with
different mass-chart outputs. Likewise, a fixed running mass is compatible
with both $s-m_R^2$ and $s-m_R^2-\epsilon$, which have different poles.
These countermodels prove that two-channel exhaustion and analyticity do not
emit the electroweak character, criticality split, or physical kernel. Nor is a formal path
realization sufficient: any finite polynomial can be encoded by assigning one
weighted path to each monomial. A non-vacuous certificate must independently
fix its transitions, path measure, signs, quotient action, exhaustive census,
rigidity gap, and target-independence proof.

If the source selector is proved deterministic and globally unique, the source
law is a delta measure and no Monte Carlo source propagation is required. If
the source emits a nondegenerate law, the simulator must propagate that law
through the pole map. Scale variation and truncation envelopes are not source
covariance, and neither a point estimate nor componentwise intervals determine
the covariance of \((\Re s_W,\Im s_W,\Re s_Z,\Im s_Z)\).
Summable clock-Hamiltonian refinements give a unique limiting gap; fixed
Lipschitz RG segments compose uniquely; Rouch\'e, Nielsen, and analytic
conjugacy isolate gauge-independent poles. If a determinant defect is bounded
by \(\epsilon\), \(a=D'(s_0)\ne0\), and \(|D''|\le K\), every $r$ with
\(\epsilon<|a|r-Kr^2/2\) encloses the unique displaced pole and propagates to
mass/width bounds through the selected square root. The four absent source
objects are therefore the factorized clock packet, independently weighted
electroweak carrier, criticality split-character carrier, and
Becchi--Rouet--Stora--Tyutin (BRST)-complete pole-kernel packet. A runtime
dependency graph proves computational separation only. A source-only claim
also requires a target-independent specification.

The electroweak chapter therefore presents four mathematically distinct
surfaces: the selected-carrier chart, the candidate value law, the
conditional finite quotient-path theorem, and the
measured-reference inverse surface used for diagnostic \(W/Z\) values. A physical
pole surface requires the additional source-root, scale, RG/scheme,
two-point-kernel, rigidity, uncertainty, and provenance receipts.

\subsection{Repair-tuple selection and the color amplitude/loop split}\label{ew-repair-tuple-selection}

The candidate electroweak tuple sits inside a two-parameter freedom that the
stated premises do not remove. With
\(\rho_{\mathrm{EW}}=(\alpha_2-\alpha_Y)/(\alpha_2+\alpha_Y)\), the coherent
quadratic family is
\[
\tau_{2,\mathrm{tree}}^{\mathrm{exact}}=-c\,\eta_{\mathrm{source}}^2,
\qquad
\delta n_{\mathrm{tree}}^{\mathrm{exact}}=d\,(1-\rho_{\mathrm{EW}})\,\eta_{\mathrm{source}}^2,
\]
and an open neighborhood of distinct \((c,d)\) values preserves the positive
mass-chart domain and defines consistent electroweak repairs. The charged leg is the
\(\SU(2)_L\) coupling correction \(\delta\alpha_2=\alpha_2\tau_2\); the neutral leg is the
hypercharge screening correction carried by \(\delta n\).

The compact-gauge reconstruction fixes the color count but does not fix the
channel assignment or normalization that would close this freedom.
Doplicher--Roberts/Tannaka reconstruction returns the color triplet sector with statistical
dimension \(d(\rho_3)=N_c=3\), and its standard conjugate intertwiner satisfies
\(R^*R=d(\rho_3)=N_c\), so \(R\) has norm \(\sqrt{N_c}\) while
\(R/\sqrt{N_c}\) is the isometry; a closed color loop carries the full
dimension \(N_c\). The two repair legs sit at different levels. The charged leg is the broken,
mass-generating \(\SU(2)_L\) channel, driven by the color-singlet condensate amplitude whose
large-\(N_c\) scaling is the decay-constant scaling \(\sqrt{N_c}\) assigned to one standard
conjugate intertwiner. The
neutral leg would be the unbroken hypercharge screening channel, a
vacuum-polarization loop that closes the color line and carries raw trace
weight \(N_c\). Under these extra hypotheses the alternative quadratic model
has
\[
c=\frac{\sqrt{N_c}}{2}=\frac{\sqrt3}{2},
\qquad
d=\frac{N_c}{2}=\frac32,
\]
where the chart coefficient \(d=N_c/2\) represents the raw neutral trace
because
\((\alpha_2+\alpha_Y)(1-\rho_{\mathrm{EW}})=2\alpha_Y\).

The resulting comparison is a target-informed, fixed-slice diagnostic of that
alternative model. Profiling the charged coefficient against
the running-tree \(W\) coordinate is independent of the criticality Jacobian core
and inherits the declared \(\alpha_2\), \(v\), running, and scheme surface. It selects
\(c=0.8670\), while profiling the neutral coefficient returns \(d=1.461\). Both values are
target-conditioned fit outputs. Their proximity to \(\sqrt3/2\) and \(3/2\) is an internal
calibration observation, not physical \(W/Z\) evidence, and no experimental confidence band is
assigned across the noncommensurate chart and mass conventions. That slice is not a test of the complete candidate
value law, which changes both coordinates and includes higher path characters;
it therefore neither excludes that law nor proves or excludes a separate
running-tree companion.

The color-balanced rule does not equal the complete candidate value law in
Theorem~\ref{thm:particle-conditional-d10-qt}; it gives a different \(W\)
coordinate. A proof of its amplitude and trace hypotheses would establish
that alternative quadratic model, not the finite quotient-path law. A source
derivation requires the two-channel quotient, exhaustive charged and neutral
path lists, the
\((1,2,1)\) and \((1,4,2)\) incidences with their \(1/3\) color measure and
\(1/6\) central trace, the fibre Gram form and residual pairing, an exhaustive
deformation class, and a positive selector gap. Measured \(W/Z\) agreement
cannot replace that certificate. The pre-repair carrier must be specified
independently of the desired outputs as an observer-like self-reading patch
with bounded local state, typed ports, readback records, and admissible
feedback or repair moves. It may contain neither the target masses nor the
desired incidence and central-trace coefficients. Independent verification
must recompute path legality, exhaustion, exact response rank and nullspace,
central action, oriented mismatch first variation, the fibre Gram form, and
the deformation census. The order of
\(\mathbb Z_6\) supplies the averaging coefficient \(1/6\), not automatically
a normalized trace \(1/6\): the checker must separate the trivial physical
matter-kernel action from the declared center-label/transport trace space
and derive the trace, source amplitude, and subtraction sign in the latter.

\subsection{Selector rigidity and the discrete repair-law boundary}\label{ew-selector-rigidity}

The candidate electroweak selector closes the continuous \((c,d)\) freedom of the
preceding subsection. Writing \(x=\tau_{2,\mathrm{tree}}^{\mathrm{exact}}\),
\(y=\delta n_{\mathrm{tree}}^{\mathrm{exact}}\), and
\(\kappa=(\alpha_Y+\alpha_2)/\alpha_Y\), the selector
\[
J_{\mathrm{EW}}(x,y)=\frac{x^2\left[1+(2\eta_{\mathrm{source}}+x)^2+4x^2\right]}{1+4x^2}
+\frac{\kappa^2}{4}(1+4x^2)y^2
\]
admits the exact factorization
\[
J_{\mathrm{EW}}-\frac34x^2-\frac{\kappa^2}{4}y^2
=\frac{x^2}{4(1+4x^2)}
\left[8(x+\eta_{\mathrm{source}})^2+8\eta_{\mathrm{source}}^2+1
+4\kappa^2(1+4x^2)y^2\right],
\]
so \(J_{\mathrm{EW}}\ge\frac34x^2+\frac{\kappa^2}{4}y^2\) with equality exactly at the
origin, and \(x^2+y^2\ge r^2\) forces
\(J_{\mathrm{EW}}\ge\min\{3/4,\kappa^2/4\}r^2\). The selector picks the zero
deformation with a quantitative gap. The candidate value-law tuple has
\(J_{\mathrm{EW}}\approx3.12\times10^{-7}>0\), so it is incompatible with the selector,
and the residual electroweak freedom is
one discrete choice between two incompatible laws rather than a continuous
chart family.

On the source-only branch the two branch points give the tree/chart
coordinates
\[
\begin{aligned}
\text{zero-selector law:}\quad
&M_W/E_\star=6.579631\times10^{-18}, &
M_Z/E_\star&=7.463335\times10^{-18},\\
\text{carrier value law:}\quad
&M_W/E_\star=6.578870\times10^{-18}, &
M_Z/E_\star&=7.463750\times10^{-18},
\end{aligned}
\]
with unclosed-clock displays \((80.3301,\,91.1191)\) and
\((80.3208,\,91.1242)\)~GeV and discrete ambiguity widths of \(9.3\)~MeV on
\(M_W\) and \(5.1\)~MeV on \(M_Z\). These are running/chart
coordinates. The PDG mass targets use a mass-dependent-width Breit--Wigner convention, and
complex-pole masses use another convention. The exact convention map distinguishes
\(M=\operatorname{Re}\sqrt{s}\) from the legacy \(\sqrt{\operatorname{Re}s}\)
coordinate. The running/chart coordinates are not physical observables and
carry no theory covariance. They cannot decide between the laws; the decision
object is a source-law
selection principle, and the quantitative mechanism of the source-action
rigidity theorem applies once the electroweak feature basis and moment vector
are emitted.

The two-loop comparison gives \((79.115335,89.802735)~\mathrm{GeV}\). It
combines a minimal supersymmetric Standard Model one-loop baseline with a
Standard Model two-loop-minus-one-loop increment, an inconsistent hybrid
prescription. A partial Feynman-gauge pole calculation gives approximately
\((79.53284,89.71232)~\mathrm{GeV}\). Its definition of \(v\), tadpole treatment, field-content and
threshold matching, scale choice, and higher-order corrections are open. The outputs establish
neither a unique scheme conversion nor a unique \(1\)--\(2\%\) defect, and they do not exhaust
the physical prescription family.

A conditional finite electroweak root/Cartan carrier emits the value-law coefficients
exactly: the transitive \(C_3\) color action has invariant measure \(1/3\),
the regular \(\mathbb Z_6\) register has rank-one trace \(1/6\), the
transitive four-slot register has slot measure \(1/4\), the fixed
independent-product source law gives
\(\lambda_{\mathrm{EW}}=\eta_{\mathrm{source}}^2/(4\rho_{\mathrm{EW}})\), and
the depth-two path census reproduces the candidate quadratic response
polynomials. A companion \(\SU(3)\)-Casimir carrier emits the declared
criticality normalizations, \(\kappa_\lambda=C_F^2=16/9\), the response coefficients
\(3/2+\rho_{\mathrm{EW}}/4\) and \(4/3-\rho_{\mathrm{EW}}/54\), and the
logarithmic character \(\rho_{HT}=\log(1+\tau_2)\) from multiplicativity with
unit derivative. Both carriers are target-exposed candidates without a
source-independent quotient-path selection or census certificate.

The declared criticality Jacobian entries follow exactly from the core,
\(\partial m_t/\partial y_t=m_{t,0}/y_0\) and
\(\partial m_H/\partial\lambda=m_{H,0}/(2\lambda_0)\), and the exact readout
is \(m_t=m_{t,0}(1+r_y)\) and \(m_H=m_{H,0}\sqrt{1-r_\lambda}\) with the
exact linearization error \(m_{H,0}r^2/(2(1+\sqrt{1-r})^2)\). The
synchronized criticality core is target-conditioned: its synchronization scale
minimizes an objective containing the Higgs and top comparison values.
Without that switch, the literal candidate one-loop source equations give the target-free tree coordinates
\(m_t^{\overline{\mathrm{MS}}}/E_\star=1.267758\times10^{-17}\) and
\(m_H^{\mathrm{tree}}/E_\star=9.427653\times10^{-18}\), with unclosed-clock
displays \(154.78\)~GeV and \(115.10\)~GeV and a QCD-converted top display of
\(164.13\)~GeV. These are one-loop tree coordinates on the declared surface;
they are not physical poles and do not satisfy the pole conditions of
\S\ref{wzh-source-closure-envelope}.

The candidate one-loop coordinates form a zero-continuous-%
parameter family. The literal core derives its top boundary from the gauge
sector through the double-criticality condition
\(\lambda(\mu_b)=0\), \(\beta_\lambda(\mu_b)=0\), which fixes
\(y_t(\mu_b)=[(2g_2^4+(g_2^2+g_Y^2)^2)/16]^{1/4}\). The candidate branch
imposes this at the gauge-unification scale \(\mu_U\), while the electroweak
transmutation in the same model anchors \(v\) at the pixel energy
\(E_{\mathrm{cell}}=E_\star/\sqrt P\), so the candidate model carries two
different high-scale anchors. Evaluating the same law at the named source
scales gives, at one loop, \((m_t,m_H)=(164.1,\,115.1)\)~GeV at \(\mu_U\),
\((170.4,\,127.8)\) at \(E_{\mathrm{cell}}\), and \((170.7,\,128.3)\) at
\(E_\star\); a benchmark-validated two-loop upgrade shifts these to
\((169.4,\,119.4)\), \((175.7,\,131.8)\), and \((175.9,\,132.3)\). The
family brackets the measured pair in both channels at both loop orders, so
the candidate deficit decomposes into the boundary-scale choice plus loop
truncation, with no continuous freedom anywhere. Along the fit-free curve
the Higgs coordinate at the measured top mass is \(125.7\)~GeV at two loops,
within \(0.5\) percent of measurement and inside the declared tree-to-pole
matching band; the implied boundary scale is \(4.8\times10^{17}\)~GeV,
between \(\mu_U\) and \(E_{\mathrm{cell}}\). No boundary-scale selection theorem is supplied; numerical agreement cannot
select the scale, and the declared-surface values obtained by back-solving
from the measured pair stay classified as target-anchored fits.

The triple-criticality equation
\(\mathrm d\beta_\lambda/\mathrm d\ln\mu=0\) has no solution along the flow:
the family derivative is dominated by
\(-24y_t^3\beta_{y_t}>0\) and stays bounded away from zero across the whole
window, so every family point is a clean \(\lambda\) minimum and the
boundary scale requires source structure. A conditional variational
selection principle uses three premises. The criticality boundary is a record
that reconciles the model's two anchor
records, the gauge-unification record at \(\mu_U\) and the transmutation
record at \(E_{\mathrm{cell}}\). The reconciliation cost is quadratic in
renormalization time, and the two anchor records carry equal capacity. These
premises give a strictly convex cost with a unique minimizer
at the log-midpoint \(\sqrt{\mu_U E_{\mathrm{cell}}}=E_\star e^{-\pi}P^{-1/6}\),
hence \((m_t,m_H)=(172.63,\,125.77)\)~GeV at two loops; the implication is
proved exactly on rational sample points.

The quadratic-cost premise is exact under the declared Gaussian record
model: the one-loop chart is
inverse-affine, so \(1/\alpha(t)\) is exactly linear in renormalization time,
and the Kullback--Leibler divergence between fixed-resolution Gaussian records
with an affine stored coordinate is exactly
\((s^2/2\sigma^2)(t-t_i)^2\), a quadratic cost with no leading-order
approximation. Axiom 3 supplies relative-entropy minimization after the
reference family, affine observable, trace convention, variance, and
resolution have been declared. The Gaussian model and its normalization are
separately declared, and their derivation from the complete three-axiom schema
is open. The reconciliation placement is exact within a
separately declared port-additive mismatch and cost-minimizing repair model.
In that model, a record with two scale-parents settles at the
capacity-weighted log-mean. Equal capacity follows if the
anchor registers have equal slope and readback resolution, which is another
declared condition. These record and repair choices are declared inputs, and
their derivation from the three axioms is open. The finite criticality carrier
must establish that the boundary record carries exactly two parent ports, one
to the gauge-unification register and one to the transmutation register, and
that those anchor registers belong to the same register class at equal
refinement depth. Given the declared Gaussian record model and the stated
repair premises, these two census facts select the boundary scale,
and with it
\((m_t,m_H)=(172.63,\,125.77)\)~GeV, with no remaining continuous parameter.
The \(W/Z\) value law requires the same finite carrier census.
The equal-capacity assumption is measurable: a capacity asymmetry moves the
\(m_H\) coordinate by about \(2.1\)~GeV per \(e\)-fold. A three-loop implied
scale therefore provides a discriminating measurement of that assumption.

A conditional selection theorem follows. If the boundary
record reconciles the two anchor records with port-additive quadratic cost
in renormalization time and equal anchor capacities, the boundary scale is
exactly the log-midpoint. The inverse-affine one-loop chart and the declared
Gaussian maximum-entropy record model give the exact quadratic reconciliation
cost. The declared port-additive repair model places a two-parent record at
the capacity-weighted log-mean. The theorem therefore consumes those record
and repair models together with two finite carrier facts: the two-parent port
structure of the boundary record and the equal-class equal-depth relation of
the anchor registers. A carrier certificate of the census class supplies
those two facts.

\subsection{Target-conditioned electroweak chart envelope}\label{ew-conditional-envelope}

Evaluating two target-conditioned repair candidates at the endpoint pixel and across
the empirical-closure interval gives the following comparison envelope.
It is not a source interval: its candidate set is not an exhaustive
source-derived deformation class, and it mixes the public endpoint branch with
the source-only question.
\[
\begin{aligned}
m_H &\in [125.183,\,125.232]~\mathrm{GeV}, & &\text{measured } 125.13\pm0.11,\\
m_t &\in [172.278,\,172.352]~\mathrm{GeV}, & &\text{measured } 172.1\pm0.6,\\
c_W &\in [80.3692,\,80.3774]~\mathrm{GeV}, & &\text{PDG BW coordinate } 80.3692\pm0.0133,\\
c_Z &\in [91.1880,\,91.1983]~\mathrm{GeV}, & &\text{PDG BW coordinate } 91.1880\pm0.0020.
\end{aligned}
\]
The Higgs and top entries are target-conditioned comparisons on their declared mass scheme. The
\(W/Z\) entries are chart coordinates. The PDG targets use mass-dependent-width Breit--Wigner
parameters, and complex-pole masses use another convention, so no physical coverage or pull is
assigned. At the endpoint pixel the two repair selections give \(Z\)-chart coordinates
\(91.18798\) and \(91.18801\,\mathrm{GeV}\). The selection spread between the two displayed laws is
roughly eight \(\mathrm{MeV}\) in the \(W\) chart, thirty \(\mathrm{keV}\) in the \(Z\) chart, sixteen
\(\mathrm{MeV}\) in \(m_H\), and thirty \(\mathrm{MeV}\) in \(m_t\). These
figures compare two ans\"atze; they do not bound the full admissible deformation
class.

\subsection{Residual attribution across the electroweak rows}\label{ew-residual-attribution}

On this comparison envelope, the Higgs and top residuals are target-conditioned mass-scheme
diagnostics. No \(W/Z\) standard-deviation assignment is valid because the chart coordinates,
PDG mass-dependent-width parameters, and complex poles use different conventions.
The \(W\)-chart spread is repair-selection freedom: the two candidate repair laws differ by about
eight \(\mathrm{MeV}\). The \(m_H\) and \(m_t\) offsets track the same repair-selection width folded through the
declared Higgs and top Jacobian cores.

The \(Z\)-chart row is the most sensitive adapter case. The empirical-closure pixel inherits
the frozen electromagnetic anchor deficit through the fine-structure endpoint, and \(c_Z\) tracks
the pixel with \(\mathrm{d}c_Z/\mathrm{d}P\approx123~\mathrm{GeV}\) per unit \(P\). Two
propagation branches separate the mechanism. The direct branch injects the certified anchor shift
into the running couplings and re-evaluates the candidate \(W/Z\) chart. That branch shifts
\(c_Z\) by \(-66\) to \(-87~\mathrm{MeV}\) on the hypercharge line and by roughly \(-230\) to
\(-302~\mathrm{MeV}\) on the proportional line, exceeding the adapter displacement by an order of magnitude,
so it is excluded. The pixel branch moves the endpoint to the measured fine-structure value and
propagates the induced pixel shift through \(\mathrm{d}P/\mathrm{d}A_{\mathrm{Th}}\). That branch
moves the empirical pixel onto the calibration pixel to within \(4\times10^{-7}\) in \(P\),
equivalently \(-0.05~\mathrm{MeV}\) in the \(Z\) chart. This is an internal adapter closure check;
it supplies no physical mass or pole.

For this comparison surface, a closed fine-structure anchor bridge would
collapse the endpoint pixel interval onto the calibration pixel and remove
this particular chart displacement. It would not supply the finite
quotient-path certificate, the independent
scale, the RG/scheme certificate, or the complex-pole gate. The anchor bridge is the object developed
in the fine-structure companion, and its source branch is the subject of Section~\ref{ew-anchor-source-branch}.

\subsection{Source branch of the fine-structure anchor bridge}\label{ew-anchor-source-branch}

The certified anchor gap \([0.620,\,0.651]\) in inverse fine-structure units names the shift the
source side supplies at the electroweak anchor for the empirical-closure endpoint to reach the
measured fine-structure value. The gap is defined by that measured requirement, so inserting it
back is circular. The non-circular source route continues the declared running convention to the
next order and tests whether the induced anchor shift lands inside the certified gap.

That test is executed. The one-loop anchor is reproduced from the declared coefficients and the
transmutation-certificate boundary data to a residual of \(10^{-11}\). The standard two-loop
running system, integrated down the same \(33.2\) e-folds with a top-Yukawa convention scan, shifts
the anchor by \([+1.62,\,+2.14]\) in inverse fine-structure units for every convention. The shift
has the correct sign and overshoots the certified gap by a factor of two and a half to
three and a half. Restoring the
balance requires the threshold and scheme-conversion data, which the matching contract classifies
as hidden fit parameters when they are undeclared. The balance requirement is quantitative: the
threshold map removes about \(2.14\) inverse-alpha, equivalent under naive coefficient accounting
to a single effective threshold near \(458~\mathrm{GeV}\). The anchor bridge has no source
scheme-lock or threshold map; its endpoint is empirical.

\section{Flavor Transport, Generation Structure, and the Yukawa Dictionary}

The gauge branch supplies a canonical rank-three response band under its
separate premises. It does not prove physical family attachment or produce the
full flavor dictionary. This paper therefore needs a separate flavor chapter.
The imported input here is the color carrier \(N_c=3\) together with the
declared physical identification of the rank-three candidate. This chapter is
the bridge between that conditional input and the matter-family chapters for
quarks, charged leptons, and neutrinos.

The flavor derivation is deliberately constructive. It does not claim that the full
OPH flavor observable is theorem-level. Instead, it builds the object chain that a
closed theorem would have to pass through. The derivation chain has the following main
architecture: start from a refinement-indexed family transport kernel, derive a centered
generation-bundle branch generator, lift that to same-label transport data, derive the induced
overlap-edge transport cocycle, reduce the cocycle to a persistent flavor observable, and then
push that observable into common sector-response objects for the downstream quark, charged-lepton,
and neutrino derivations.

Mathematically, this derivation is where ``why this family exists'' becomes a concrete technical
question. The finite Standard Model packet gives a rank-three candidate; the flavor
derivation is where one tries to turn that candidate and the declared count into a physical attachment, transport,
splitting, suppression, phase, and excitation data. The relevant objects are the intermediate
transport and spectral structures, not the final fermion masses themselves, that the mass
readouts consume.

The active chain, in slightly compressed form, is:
\begin{enumerate}[leftmargin=2em]
\item normalize a refinement-indexed family transport kernel;
\item derive a centered generation-bundle branch generator on the conditionally declared
three-generation charged bundle;
\item lift this to same-label edge-line transport and then to an overlap-edge cocycle with
explicit defect and gap bookkeeping;
\item reduce those data to projectors, spectral gaps, pair suppressions, and cycle phases;
\item push the resulting family object into sector-response objects and then into
suppression/phase tensors for the downstream matter derivations.
\end{enumerate}

Three branch-specific claim boundaries are attached to this region of the derivation. The shared
excitation dictionary is the common proof-facing base. Above that base, the charged branch carries
exact centered readback, a closed common-shift no-go, the declared same-label \(q_e\) readback, a
source-side determinant character for any fixed formal exponent vector
\[
S_M=\sum_e M_e^{\mathrm{ch}}\log q_e,
\]
a determinant-line lift on theorem-grade physical charged data, and an algebraic mass
readout from theorem-grade \(A_{\mathrm{ch}}(P)\). The theorem branch does not emit a theorem-grade
sector-isolated charged determinant exponent vector, and it does not identify a source-side
determinant character with the physical charged determinant line. The charged-lepton theorem gap is the
determinant trace-lift attachment \(3\mu(r)=S_M(r)\) on the charged determinant channel. The charged determinant channel has a source-domain no-go boundary: the uncentered trace lift is not emitted. The neutrino branch carries a target-informed weighted-cycle
candidate above a template family kernel, with compare-only bridge and absolute-attachment diagnostics. The quark
branch carries a theorem-grade source-spread obstruction: its ordered profile shapes leave a free
\((\mathbb R_{>0})^2\) fiber. Selected-class descent proves representative independence but does
not choose either positive modulus. Target-anchored mixed-convention mass coordinates and
GeV-valued mass textures are comparison data; they are neither a source-only running sextet nor
physical dimensionless Yukawa matrices. A separate compare-only microphysical bridge records
edge-statistics transport on a diagnostic surface. Together these boundaries mark where the flavor
derivation leaves the common transport backbone and passes to branch-specific closure contracts.

\section{Quark Family Derivation}

The quark branch carries a source-only non-identifiability theorem on the public physical quark frame
class \(f_P\) chosen by \(P\). The source equations determine the ordered up- and down-sector
profile rays but leave their endpoint spans as two independent positive moduli. Selected-class
descent does not remove this freedom, so no numeric mass row is emitted. Same-family and restricted
common-refinement artifacts reproduce their chosen targets only after those moduli are obtained by
target inversion. Their rows also mix renormalization conventions, while their GeV-valued matrices
are mass textures rather than physical dimensionless Yukawa matrices. This theorem does not classify
all public quark frame classes.

\subsection{What quarks are in this derivation}

Quarks are the color-charged elementary constituents of hadronic matter. In nature, the up quark
and down quark dominate protons and neutrons, while the strange quark, charm quark, bottom
quark, and top quark appear in progressively heavier and more unstable sectors. In the OPH
particle derivation, the quark branch has a closed downstream algebraic readout conditional on a
source-only physical spread datum. The target-free source equations do not emit that datum: they
fix two ordered profile rays and leave their positive endpoint spans independent. The compatible
fiber is \((\mathbb R_{>0})^2\), so the six numeric rows are not source-only predictions.

The quark route starts from the shared flavor excitation dictionary,
internalizes the target-free mass bridge on the quark mass-profile ray, and
proves selected-bridge-fiber representative independence for the
attached physical spread datum on \(f_P\), and applies the affine mean law once both positive
moduli are supplied. The stored target-anchored matrices have GeV-valued singular values drawn
from several comparison conventions. They are mass textures, not physical dimensionless Yukawa
matrices.
The same-label left-handed selector surface closes to the singleton
\(\sigma_{\mathrm{ref}}\). That lower object is a negative sheet-selector statement
on the selected quark sheet. It does not break the independent positive-rescaling action on the two
sector profiles.

\subsection{Emitted quark rows}

\begin{table}[h]
\centering
\small
\begin{tabular}{@{}llll@{}}
\toprule
Quark & Theorem scope & Public value & Mass convention \\
\midrule
Up quark & source spread non-identifiable & withheld & \(\overline{\mathrm{MS}}\), \(\mu=2\,\mathrm{GeV}\), comparison chart \\
Down quark & source spread non-identifiable & withheld & \(\overline{\mathrm{MS}}\), \(\mu=2\,\mathrm{GeV}\), comparison chart \\
Strange quark & source spread non-identifiable & withheld & \(\overline{\mathrm{MS}}\), \(\mu=2\,\mathrm{GeV}\), comparison chart \\
Charm quark & source spread non-identifiable & withheld & \(\overline{\mathrm{MS}}\), \(\mu=m_c(\mu)\), comparison chart \\
Bottom quark & source spread non-identifiable & withheld & \(\overline{\mathrm{MS}}\), \(\mu=m_b(\mu)\), comparison chart \\
Top quark & separate extraction coordinate & withheld & cross-section pole-mass chart \\
\bottomrule
\end{tabular}
\end{table}

Running quark masses are coordinates on a renormalization-group trajectory. Finite
renormalizations and scale changes alter those coordinates while preserving physical amplitudes,
so OPH can emit an RG-covariant trajectory or invariant but cannot derive the human choice of an
\(\overline{\mathrm{MS}}\) chart. The chart must be declared after source emission. The top row is
a pole extraction coordinate rather than a sixth member of one common running-mass packet. All six
rows are withheld because the spread pair is non-identifiable from the source domain. A physical
Yukawa construction would additionally transport every running coordinate to one common scale,
convert the top coordinate, emit the running Higgs expectation value in the same scheme, and apply
\(y_q(\mu)=\sqrt2\,m_q(\mu)/v(\mu)\).

\subsection{Target-free mass bridge and selected-class theorem}\label{quark-family-strongest-overlap-branch-primitives}

The same-label left-handed solver surface closes to the singleton
\(\sigma_{\mathrm{ref}}\). This is a negative same-sheet selector statement:
same-sheet rephasing preserves CKM moduli, so it does not move that selected sheet to the
physical CKM shell. Same-sheet overlap scans, chirality-swapped basis diagnostics, and other
non-sector-attached orbit improvements are compare-only and do not change the selected-class
theorem.

A separate target-free mass bridge is internalized on the emitted quark
mass-profile ray. On the minimal light branch
\[
y_u=c_u\,\varepsilon^6,
\qquad
y_d=c_d\,\varepsilon^6,
\qquad
\varepsilon=\frac16,
\]
the light-quark overlap-defect theorem emits
\[
\Delta_{ud}^{\mathrm{overlap}}
=
\frac16\log\frac{c_d}{c_u}.
\]
The emitted same-family mass object is the ray
\[
\mathcal R_{ud}^{\mathrm{mass}}
:=
\{(c_u,c_d)\in\mathbb R_{>0}^2:c_d/c_u\ \text{fixed}\},
\]
and on that ray the same scalar is equivalently the one-scalar law
\[
\Theta_{ud}^{\mathrm{mass}}
:=
\frac16\log\frac{c_d}{c_u}.
\]
The exact scalar identities are
\[
\Delta_{ud}^{\mathrm{overlap}}=\frac{t_1}{5},
\qquad
\log\frac{c_d}{c_u}=\frac65\,t_1,
\qquad
t_1=5\,\Delta_{ud}^{\mathrm{overlap}}=\frac56\log\frac{c_d}{c_u}.
\]
The mass-profile package is therefore functorial:
\[
\begin{aligned}
\eta_Q^{\mathrm{centered}}&=-\frac{1-x_2^2}{27}\,t_1,\\
\kappa_Q&=-\frac{t_1}{54},\\
x_2&=-0.5175863354681689.
\end{aligned}
\]
The odd source package is also forced:
\[
\begin{aligned}
\beta_{u,\mathrm{diag},B}^{\mathrm{source}}&=\frac{t_1}{10},\\
\beta_{d,\mathrm{diag},B}^{\mathrm{source}}&=-\frac{t_1}{10},
\end{aligned}
\]
\[
\begin{aligned}
\mathbf r_u
&=
\left(-\frac{t_1}{10},\,0,\,+\frac{t_1}{10}\right),\\
\mathbf r_d
&=
\left(+\frac{t_1}{10},\,0,\,-\frac{t_1}{10}\right),
\end{aligned}
\]
For any supplied positive spread pair one obtains
\[
\tau_u=\frac{\sigma_d}{10(\sigma_u+\sigma_d)}\,t_1,
\qquad
\tau_d=\frac{\sigma_u}{10(\sigma_u+\sigma_d)}\,t_1.
\]
These identities do not select \(\sigma_u\) or \(\sigma_d\). The mass bridge
lies outside the selected-class theorem boundary. The pure-\(B\) payload pair
is a lower realization object beneath the exact theorem, and
\(\Theta_{ud}^{\mathrm{mass}}\) is its derived scalar wrapper.

Write the same-label left-handed physical carrier as
\[
\Sigma_{ud}^{\mathrm{phys}}
:=
\left\{
(\sigma_{\mathrm{id}},\tau,U_{u,L},U_{d,L},V_{\mathrm{CKM}},I_{\mathrm{CKM}})
:
V_{\mathrm{CKM}}=U_{u,L}^\dagger U_{d,L}
\right\}/\!\sim,
\]
where
\[
(U_{u,L},U_{d,L},V)\sim
(U_{u,L}D_u,\ U_{d,L}D_d,\ D_u^\dagger V D_d)
\]
for diagonal \(D_u,D_d\in U(1)^3\). On the selected public quark frame class \(f_P\),
represented on the realized branch by the explicit common-refinement transport-frame class
\([F_0^\dagger F_1]\), the exact \(\Sigma_{ud}^{\mathrm{phys}}\) element and the attached exact
sigma datum are independent of the declared representative inside the selected bridge fiber.
Therefore the exact sigma readout descends uniquely to the selected comparison domain on
\(f_P\). It is not a public source-only mass prediction because the physical sigma datum is
target-derived.

\paragraph{Sigma descent non-selection.}
Let \(R_{\mathrm{decl}}(f_P)\) be the declared selected bridge fiber over the public quark
frame class. If a map
\[
\Sigma:R_{\mathrm{decl}}(f_P)\to\mathbb R^4
\]
is constant on that fiber, then it descends to a well-defined public datum
\(\overline\Sigma(f_P)\). This proves representative independence only. It does not select the
value of \(\overline\Sigma(f_P)\), because every constant vector in \(\mathbb R^4\) would also
descend. Therefore a value obtained from the declared running-quark target surface is
target-derived after descent.

\paragraph{Common-scale physical rejection of the reciprocal-ray candidate.}
A physical comparison first converts all six quark eigenvalues to dimensionless Yukawa singular
values in one effective theory, scheme, and scale. Using the Standard Model
\(\overline{\mathrm{MS}}\) values at \(M_Z\) tabulated by Antusch, Hinze, and
Saad~\cite{antusch2025running}, the two ordered log-gap ratios are
\[
\rho_u=1.1108888543,
\qquad
\rho_d=0.7519410008,
\qquad
\rho_u\rho_d=0.8353228768.
\]
The reciprocal-ray law requires the last product to equal one. It is therefore falsified on the
properly typed target surface. Even after granting the four endpoint Yukawas
\((y_u,y_t,y_d,y_b)\), the minimax reciprocal shape misses the held-out \(y_c\) and \(y_s\) by
\(21.556\%\) at \(M_Z\). The result persists under running: the products are
\(0.836092\), \(0.837752\), \(0.845227\), and \(0.843611\) at
\(10^3\), \(10^5\), \(10^{12}\), and \(10^{16}\,\mathrm{GeV}\), with best held-out errors
between \(19.93\%\) and \(21.42\%\). The stored \(\rho=1.294285\) formula misses the charm
coordinate by \(56.5\%\) at \(M_Z\), even under the endpoint grant.

The generic physical interface is therefore
\[
(\mu_u,\sigma_u,\rho_u,\mu_d,\sigma_d,\rho_d)\in\mathbb R^6,
\]
with an exact inverse to the six ordered positive Yukawas. The three-scalar theorem is exact only
inside its imposed reciprocal-ray, affine-mean, fixed-shared-scale subfamily. The two-spread
counterfamily below is a valid restricted non-identifiability lower bound, but it
is not a parameterization of the whole physical interface.

A separate selector analysis sharpens the missing source object. Generation-blind flavor-singlet
scalars cannot emit a fixed nonzero bifundamental Yukawa matrix. Simultaneous \(S_3\) invariance
restricts a fixed matrix to the span of \(I\) and \(J\), which has a degenerate eigenspace and
aligned up/down frames. A physical construction therefore needs a source-derived flavor-orbit
selector, such as a spontaneous invariant functional or genuine bifundamental boundary data,
together with a quark--Higgs carrier. The calibration simulator contains no quark, Higgs, or
Yukawa coupling and its enumerated observables have zero Fisher information for these coordinates.
Its \(64\mathrm{k}\) runs are at the calibration null: the two fusion-tower readings
are \(1.495245\) and \(1.484301\), while the dense first/full readings are \(1.498275\) and
\(1.497823\). Those values are not quark-mass evidence.

\paragraph{Restricted source-spread non-identifiability theorem.}
Remove every running-mass target, exact-witness, fitted-spread, and compare-only ancestor. Grant
the remaining source packet its strongest ordered three-point shape law. For
\(\rho_{\mathrm{ord}}>0\), define
\[
v_u=\frac{1}{3(1+\rho_{\mathrm{ord}})}
\bigl(-(2\rho_{\mathrm{ord}}+1),\rho_{\mathrm{ord}}-1,
\rho_{\mathrm{ord}}+2\bigr),
\]
\[
v_d=\frac{1}{3(1+\rho_{\mathrm{ord}})}
\bigl(-(\rho_{\mathrm{ord}}+2),1-\rho_{\mathrm{ord}},
2\rho_{\mathrm{ord}}+1\bigr).
\]
Both vectors have zero trace and unit endpoint span. Their adjacent-gap ratios are
\(\rho_{\mathrm{ord}}\) and \(\rho_{\mathrm{ord}}^{-1}\), respectively. Conversely, those
three conditions determine each profile up to one positive endpoint span. Hence every compatible
pair is
\[
E_u=\sigma_u v_u,\qquad E_d=\sigma_d v_d,
\qquad (\sigma_u,\sigma_d)\in(\mathbb R_{>0})^2.
\]
The group \((\mathbb R_{>0})^2\) acts freely and transitively by independent rescaling of the two
spans while fixing the source shape data. The requested four-tuple
\[
\left(\frac{\sigma_u+\sigma_d}{2},
\frac{\sigma_u-\sigma_d}{2},\sigma_u,\sigma_d\right)
\]
is not invariant under this action. No unique source-only spread package follows from the stated
source data.

\paragraph{Non-definability of the quark spectrum from the stated data.}
This ambiguity is a mathematical non-definability result. Fix any generic admissible one-Higgs
three-generation package and write
\[
Y_q=U_{q,L}\,\operatorname{diag}
\!\left(e^{\mu_q+\sigma_qv_{q,1}},e^{\mu_q+\sigma_qv_{q,2}},
e^{\mu_q+\sigma_qv_{q,3}}\right)U_{q,R}^{\dagger},
\qquad q\in\{u,d\},
\]
with \(\sum_i v_{q,i}=0\). For every \(\lambda_u,\lambda_d>0\), replace
\(\sigma_q\) by \(\lambda_q\sigma_q\) while holding the frames fixed. Gauge representations,
anomaly cancellation, hypercharges, one-Higgs Yukawa completability, the CKM matrix, intrinsic CP
capability, and the weak-sector counting clause are unchanged. In addition
\[
\det\exp(\lambda_q\sigma_qv_q)
=\exp\!\left(\lambda_q\sigma_q\sum_i v_{q,i}\right)=1,
\]
so determinant normalization does not select either modulus.

Every member of the family carries the same gauge group, matter content,
color count, and generation count, so no stated structural datum separates
them.
which contains no Yukawa eigenvalue. The members are nevertheless physically inequivalent because
the declared physical-equivalence relation preserves Yukawa invariants. Axiom~3 also does not
remove the family: its gauge-invariant local constraint values and associated multipliers are not
numerically specified and no map from \(P\) to quark Yukawa multipliers is emitted. Thus the stated
axioms and declared completions plus fixed \(P\) admit a free
\((\mathbb R_{>0})^2\) family of equally admissible quark spectra.
There is no smallest positive rescaling; its infimum is zero, which would give the massless or
degenerate boundary rather than the observed spectrum. Therefore the stated axioms do not define
a unique quark mass spectrum. This conclusion uses no additional axiom and can be overturned only
by deriving, from the existing axioms, a source functional that is nonconstant on this explicit
counterfamily.

\paragraph{The kernel interface: the exact content of the missing derivation.}
Two companion theorems make that boundary executable. First, the admissibility battery recorded
for the family-transport branch (positive-semidefinite hermitian descendants, open three-cluster
gaps at every level, a simple centered spectrum, the conjugacy-Riesz margin, persistent projector
labeling, and overlap-edge amplitudes above the floor) places no constraint on the spectrum of
the centered compressed branch generator: for every pair \((r,s)\in(\mathbb R_{>0})^2\) there is a
certificate-passing two-level kernel whose generator has raw gap ratio exactly \(r\) and spectral
span exactly \(s\). The construction places the target spectrum by unitary conjugation,
\(T=Q\,\operatorname{diag}(\sqrt{\mu})\,Q^{\dagger}\), so the descendant spectrum is exact, and
shrinks the refinement drift geometrically until the Riesz margin passes. The persistence
certificates are therefore a filter, not a generator: a kernel derivation that merely passes them
cannot emit the ordered ratio constant, the mean-law coordinate, or the spans.

Second, the three-scalar interface on record is exact only on its imposed reciprocal-ray
subfamily. At fixed shared scale \(g_{\mathrm{ch}}\), that subfamily factors through
\[
(r,\sigma_u,\sigma_d)\in(\mathbb R_{>0})^3,
\qquad
\rho_{\mathrm{ord}}=\frac{3}{2+r},
\qquad
x_2=\frac{r-1}{r+1},
\]
with the rays, the affine mean coefficients \(A_{ud},B_{ud}\), the sector means, and the six
coordinates closed-form in the triple, and with the explicit left inverse
\(\sigma_u=\tfrac12\ln(m_t/m_u)\), \(\sigma_d=\tfrac12\ln(m_b/m_d)\), and \(r=3/\rho-2\) at
\(\rho=\ln(m_c/m_u)/\ln(m_t/m_c)\); the forward map is injective and the executed round trip
closes at machine precision inside that subfamily. It is not the general interface of two
ordered three-point spectra.

Indeed, for any \(x\ne\pm1\), let
\[
L(x)=\operatorname{ctr}(-1,x,1),
\qquad
Q(x)=\operatorname{ctr}(1,x^2,1).
\]
Their Gram determinant is
\[
\det\operatorname{Gram}(L,Q)=\frac43(1-x^2)^2.
\]
Thus \(L,Q\) form a basis of the centered three-vector plane, and every centered sector spectrum
has unique coordinates \(E_q=a_qL+b_qQ\). Once both \(a_q\) and \(b_q\) are free, \(x\), and
hence \(r\), is a basis choice rather than an additional invariant of the spectrum. The general
common-scale eigenvalue interface has six scalar coordinates: two centered coordinates and one
mean for each sector. A proposed \(r\) plus four centered coordinates plus two means is a
redundant seven-coordinate chart. The reciprocal-ray law removes one centered coordinate per sector and
thereby recovers its special three-scalar chart.

By the freedom theorem and the non-definability theorem above, the stated source data select neither
the ray triple nor the six-scalar physical interface. The executable acceptance harness is
only a fail-closed score of the reciprocal-ray subfamily.

The edge data do not remove the ambiguity. Even after granting source values
\(S_{13},S_{23},\delta_{21}>0\), every positive pair can be written as
\[
\sigma_u=S_{13}+c_u\delta_{21},\qquad
\sigma_d=S_{23}+c_d\delta_{21}
\]
for suitable \((c_u,c_d)\). The source equations emit no rule fixing these two coefficients. The
specific coefficients used by the numerical candidate are therefore a declared ansatz from a
hand-written family-transport template, not an OPH-derived kernel.

Define
\[
\sigma_{\mathrm{seed}}^{ud}:=\frac{\sigma_u+\sigma_d}{2},
\qquad
\eta_{ud}:=\frac{\sigma_u-\sigma_d}{2},
\]
\[
\begin{aligned}
A_{ud}&:=\frac{1}{2(1+\rho_{\mathrm{ord}}-x_2^2)},\\
B_{ud}&:=\frac{1}{2\!\left(1-x_2^2-\frac{x_2^2}{1+\rho_{\mathrm{ord}}}\right)},\\
\rho_{\mathrm{ord}}&=1.294284936377706.
\end{aligned}
\]
Then
\[
g_u
=
g_{\mathrm{ch}}
\exp\!\bigl(-(A_{ud}\sigma_{\mathrm{seed}}^{ud}-B_{ud}\eta_{ud})\bigr),
\qquad
g_d
=
g_{\mathrm{ch}}
\exp\!\bigl(-(A_{ud}\sigma_{\mathrm{seed}}^{ud}+B_{ud}\eta_{ud})\bigr).
\]
The Jacobian from \((\sigma_u,\sigma_d)\) to
\((\log(g_u/g_{\mathrm{ch}}),\log(g_d/g_{\mathrm{ch}}))\) has determinant
\(-A_{ud}B_{ud}\neq0\). The two free moduli therefore change the mass readout; they are not a
gauge redundancy. Given a source-only physical spread datum, the affine mean law emits
\((g_u,g_d)\) algebraically on \(f_P\). The target-derived packet is retained as a mixed-convention
comparison witness. Its dimensionful matrices are mass textures. The conditional selected-class route
can be written as
\[
\!f_P+\Sigma_{ud}^{\mathrm{source}}(P)
\Longrightarrow
\bigl(\sigma_u,\sigma_d,\sigma_{\mathrm{seed}}^{ud},\eta_{ud}\bigr)_{\mathrm{phys}}
\Longrightarrow
(g_u,g_d)
\Longrightarrow
\bigl(m_u,m_d,m_s,m_c,m_b,m_t\bigr),
\]
followed, only after common-scale RG transport and division by the running Higgs expectation
value, by dimensionless Yukawa matrices. The first arrow is obstructed by the free
\((\mathbb R_{>0})^2\) action on the stated source domain.

\paragraph{Target-anchored \(S_3\) two-mode witness.}
An ansatz built with visible target coordinates reproduces the six mixed-convention comparison
coordinates numerically. Its exact mathematical component starts
from the transposition Cayley graph of \(S_3\). The adjacency spectrum is
\(3,0,-3\) with multiplicities \(1,4,1\), so the Laplacian spectrum is \(0,3,6\) and
\[
\frac{e^{-3\tau}-e^{-6\tau}}{1-e^{-3\tau}}=e^{-3\tau}.
\]
This identity is a finite-group theorem. No OPH theorem identifies the three distinct
isotypic heat values with three generations or supplies the proposed heat time
\[
\tau_f=\frac P4-\frac{\pi\alpha_U}{5}.
\]

The ansatz sets \(w=\pi\alpha_U\), \(r=e^{-3\tau_f}\),
\(\rho=3/(2+r)\), and \(x=(r-1)/(r+1)\), and then uses
\[
\begin{aligned}
e_u&=S_{13}+\frac{\rho\delta_{21}}{1+\rho},&
e_d&=S_{23}+\frac{\delta_{21}}{2(1+\rho-x^2)},\\
\bar\sigma_u&=e_u-w\Delta S_{13},&
\bar\sigma_d&=e_d-w(1-\Delta S_{13}),\\
a_u&=e_u+\frac w2,&
a_d&=e_d-\frac w5,\\
b_u&=b_u^{\mathrm{ray}}-\frac{w\rho}{10},&
b_d&=b_d^{\mathrm{ray}}-\frac w4.
\end{aligned}
\]
The ray coordinates have the symbolic form
\[
b_u^{\mathrm{ray}}
=a_u\frac{-\rho x+\rho-x-1}{(1+\rho)(x^2-1)},
\qquad
b_d^{\mathrm{ray}}
=a_d\frac{-\rho x-\rho-x+1}{(1+\rho)(x^2-1)}.
\]
The ansatz feeds \(\bar\sigma_u,\bar\sigma_d\) into the candidate affine mean law above and
sets \(E_q=a_qL+b_qQ\). With the declared comparison inputs, this gives
\[
\begin{gathered}
r=0.3179975211,\qquad \rho=1.2942205385,\qquad x=-0.5174535369,\\
(a_u,b_u)=(5.6430129216,-4.9927828217),\qquad
(a_d,b_d)=(3.3952170977,-1.8369623926).
\end{gathered}
\]

\paragraph{Conditional normalized-trace lemma.}
Let \(\ell:M_d(\mathbb C)\to\mathbb C\) be complex linear, invariant under all unitary
conjugations, and normalized by \(\ell(I)=1\). Unitary conjugacy of rank-one projectors and
additivity over an orthonormal resolution of \(I\) give
\(\ell(A)=\operatorname{Tr}(A)/d\), so every rank-one channel has weight \(1/d\).
Likewise, a width operator on \(M_d(\mathbb C)\) invariant under the full
\(U(d)\times U(d)\) left-right action is scalar by Schur's lemma; total Hilbert--Schmidt width
\(t\) therefore assigns \(t/d^2\) to each unit slot. Consequently, \emph{if} the physical heat,
up-odd, down-odd, up-even, and down-even responses are respectively identified with one isotropic
slot in \(M_5(\mathbb C)\) of total width \(5w\) and rank-one modules of dimensions
\(2,5,10,4\), the coefficients \(w/5,w/2,w/5,\rho w/10,w/4\) follow.
The full invariance and normalization of the response law, those five module assignments, their
signs and orientation, and exclusion of competing assignments are hypotheses outside OPH.
Thus this lemma closes the denominator arithmetic conditional on the physical channel
functor; it does not remove target dependence from the formula.

The per-particle comparison table is not a source-only prediction. It is a
target-anchored diagnostic: the evaluator emits dimensionless coordinates, inserts
an unproved one-GeV unit, and compares them to a table that mixes light-quark
\(\overline{\mathrm{MS}}\) values at \(2\,\mathrm{GeV}\), heavy-quark self-scale values, and a
separate top extraction.

The raw diagonal residual sum against the supplied central values and nominal standard deviations
is \(1.1653\), with a maximum relative residual of \(0.2946\%\). It is not a likelihood or
goodness-of-fit statistic: discovery used these targets, there is no source-theory covariance, and
the rows are not one common-scale spectrum. In the target-conditioned \(219{,}615\)-member denominator
grammar, the selected tuple \((5,2,5,10,4)\) is the unique minimum of that raw residual sum, but
eight formulas tie its best maximum error. The grammar contains the target-exposed inputs, graph,
entries, assignments, signs, and functional forms, so it is not a global look-elsewhere correction.

The evaluator is runtime-target-separated but target-dependent.
\(S_{13},S_{23},\delta_{21},\Delta S_{13}\), and \(g_{\mathrm{ch}}\) descend from
the explicitly hand-written family-transport template. The last quantity is the
dimensionless template-eigenvalue mean plus its minimum gap; treating it as GeV supplies the
missing unit by hand. The selected pixel branch depends on an
internal quark-spectrum continuation. The edge laws add log-overlap
suppression to a linearly scaling \(TT^\dagger\) gap, and \(\Delta S_{13}\) is a selected
basis-dependent matrix entry. The ansatz is target-informed and supplies no
physical mass law. It does not contradict the non-identifiability theorem.

\paragraph{Sufficient source conditions for flavor.}
The proof analysis isolates a sufficient conditional completion without instantiating it.
A physical theorem would require all of the following on one acyclic,
target-free source dependency graph:
\begin{enumerate}[leftmargin=2em]
\item a unique interval-certified source root for \(P,\alpha_U(P)\), with no dependency on an
internal quark continuation;
\item a physical \(S_3\) family-carrier attachment, the heat-time law
\(\tau_f=P/4-\pi\alpha_U/5\), the \(\rho,x\) dictionary, and the two edge-response laws;
\item a source-labeled simple-spectrum family generator and charged seed with exact
common-refinement intertwiners, monomial transport of the label rays, positive common
normalization, and selected edge magnitudes bounded away from zero;
\item a physical response-channel functor that proves the normalized-trace hypotheses, module
assignments, signs, orientation, competing-channel exclusion, and the affine mean law on that
carrier, on a domain where the \(A,B\) denominators do not vanish;
\item a common-scale physical readout with
\(m_q(\mu)=y_q(\mu)v(\mu)/\sqrt2\), including field normalization, units, and the running Higgs
expectation value in the same scheme;
\item a source-only RG packet with existence and non-blowup across the required intervals, a
locally Lipschitz beta system, fixed threshold ordering and matching maps, top conversion, and
the declared comparison charts.
\end{enumerate}
Given these hypotheses, composition of the single-valued maps makes the final sextet unique.
This is a sufficient proof-obligation contract, not a proved-minimal OPH theorem.
The numerical table is not its corollary: the evaluator's coordinates are dimensionless and no
instance of the final two receipts exists. What the
counterfamily proves as a necessary condition is narrower and decisive: any successful source
functional must be nonconstant on the independent
\((\lambda_u,\lambda_d)\) centered-spread rescaling orbit.

\paragraph{Cumulant specification.}
The representation-slot cumulant ansatz declares
\[
F=\mathbb C^3_{\mathrm{perm}}\oplus V_{\mathrm{std}}
\cong\mathbf1\oplus2V_{\mathrm{std}},
\qquad \dim F=5,
\qquad \dim F_0=4,
\]
and the composite response dimensions
\[
(29,432,22,32,840,1008,432,1584).
\]
These integers are arithmetically correct for the declared direct sums and tensor products. If a
continuous Gaussian source space, full unitary isotropy, reversible independent sheets, the listed
effect ranks, and the listed signs are additionally assumed, the second-cumulant identity gives
\[
\log\mathbb E[e^X]=\frac12\operatorname{Var}X
=\frac{r}{d}w^2.
\]
This conditional identity does not construct the physical effects or select the modules.
In particular,
\(\dim\operatorname{End}_{S_3}(\mathbf1\oplus2V_{\mathrm{std}})=5\), so
\(S_3\)-invariance does not force one scalar covariance on the two-copy multiplicity space.
Pooling heterogeneous sums such as \(\operatorname{End}(F)\oplus F_0\) into a single
\(1/29\) response requires an additional symmetry that mixes physically distinct blocks.
The proposed \(F\) also has heat spectrum \(0^{(1)},3^{(4)}\), not the regular representation's
additional sign-sector value \(6^{(1)}\); the physical regular-heat-to-family attachment is open.
Finally, a genuine Gaussian covariance term is nonnegative. The negative \(w^2\) entries require a
separate signed response or subtraction law, not orientation reversal alone.

With \(w=\pi\alpha_U\), the declared cumulant formulas are
\[
\begin{aligned}
\bar\sigma_u&=3P+\left(5+\frac4{15}\right)w+\frac{w^2}{29},&
\bar\sigma_d&=2P+\frac4{15}w-\frac{w^2}{432},\\
a_u&=3P+\left(5+\frac45\right)w+\frac{w^2}{22},&
a_d&=2P+\left(1+\frac1{32}\right)w+\frac{w^2}{420},
\end{aligned}
\]
and
\[
\frac{g_{\mathrm{ch}}}{v}
=2\exp\!\left[-2\pi+\frac{P}{1008}
+\frac{w^2}{432}-\frac{w^2}{1584}\right].
\]
Conditional on these formulas and the inherited heat-time, ray, even-response, affine \(A,B\),
and exponentiation laws, the downstream map is deterministic. Using the electroweak
\(v\)-display gives a maximum residual of \(0.2943587\%\) and a raw diagonal nominal-residual sum
of \(1.1613966\) against the same mixed-convention comparison table. Neither number is a
likelihood: the calculation uses visible target-derived effective coordinates,
there is no common-scale renormalization-group packet or theory covariance,
and the supplied cumulant comparison changes
the nominal \(u\)-row uncertainty relative to the preceding bundle.

Target dependence and ablation are decisive. The target-inferred denominators for four
visible effective coordinates are approximately
\[
29.19665,\qquad428.18579,\qquad21.90604,\qquad836.47734,
\]
close to the selected \(29,432,22,840\), while no target-independent grammar
of alternative module expressions is specified. The zero-\(w^2\),
zero-\(\delta_g\) ablation has the lower
maximum residual \(0.2142293\%\) and raw residual sum \(0.6694007\). This does not make
the ablation a physical theory; it shows that the detailed covariance data is not selected by
the numerical agreement. The pixel input depends on an internal quark model,
and the selector uses mixed sources.

The ansatz is an explicit, falsifiable target-conditioned specification. The
claimed rescaling-orbit statement must also be scoped correctly:
\[
\|\lambda E_q-E_q\|^2=(\lambda-1)^2\|E_q\|^2
\]
has its minimum at one only on the orbit through the declared cumulant vector. For a generic base
vector \(C_q\), the minimizer of \(\|\lambda C_q-E_q\|^2\) is
\(\langle C_q,E_q\rangle/\|C_q\|^2\), and no OPH dynamics requires minimizing this
postulated residual. The target-conditioned module ansatz does not satisfy the
source conditions above or define a physical quark-mass law.

\paragraph{Finite maximum-entropy boundary and conditional register rigidity.}
The finite-MaxEnt premise is false. On a finite spectrum, maximizing
entropy under mean and covariance constraints gives a discrete exponential-quadratic Gibbs law,
not a Gaussian density. Explicitly, on support \(\{-2,-1,0,1,2\}\),
\[
p_A=(1/12,1/6,1/2,1/6,1/12),\qquad
p_B=(1/16,1/4,3/8,1/4,1/16)
\]
both have mean zero and variance one, while
\(\kappa_4(p_A)=0\) and \(\kappa_4(p_B)=-1/2\). The actual finite-support MaxEnt law at
variance one has a nonzero fourth cumulant (approximately \(-0.46815\)). Hence neither finite
maximum entropy nor the first two moments entails the Gaussian two-cumulant
truncation. Gaussianization
can be recovered only conditionally from an exported triangular array with a Lindeberg or adequate
mixing/bounded-dependency condition, normalization, and covariance convergence; no such array is
emitted by the source.

The quark-register proposal instead states an exact primitive-path
certificate. Conditional on its nine declared path hypotheses, normalized
trace fixes each path weight as rank over the
declared register dimension, including
\[
-\frac15,\frac4{15},\frac1{29},-\frac1{432},\frac45,\frac1{22},
\frac1{32},\frac1{420},-\frac1{10},-\frac14,
\frac1{1008},\frac1{432},-\frac1{1584}.
\]
Exact record projectors also allow only the zero or unit scalar multiplier, and inert tensor
refinement preserves normalized trace weights. This is a useful rigidity theorem, but its
hypotheses contain the required physical content: the typed register construction, exhaustive
primitive-path catalogue, effect ranks, structural multiplicities, signs, winding character,
refinement intertwiners, implementation invariance, and positive-gap branch selector. A neutral
finite register can be changed while preserving the broad structural signature and changing the
normalized response, so the three axioms alone do not select this register.
That countermodel is scoped
to the broad signature; it is not a realization of every stronger carrier condition one might add.

The selector results sharpen the same boundary. An equivariant section requires a
stabilizer-fixed candidate in each orbit fiber; an ambiguous transitive fiber need not have one.
Conversely, a complete finite invariant candidate class with a unique positive-gap minimizer has a
stable selector when refinement error is below half the gap. The proposal
supplies neither a physically
complete candidate class nor a source-derived cost and gap. The absolute-scale rescaling no-go,
finite-renormalization scheme ambiguity, interval-root schema, and piecewise-RG uniqueness theorem
are likewise exact obstruction or conditional well-posedness statements. No actual target-free
source map with strict interval inclusion, operational clock, or fixed
beta/threshold/matching packet accompanies the proposal. The register
composes algebraically with the target-conditioned cumulant ansatz and
authorizes no numerical quark-mass claim.

\paragraph{Oriented family-transport carrier test.}
The primitive oriented family-transport proposal supplies explicit complex
\(3\times3\) matrices \(T_0,T_1\), but the simulator test must not
manufacture them through a fitted readout. The fixed
direct observable is therefore the mean of the natural three-label permutation matrices carried by
the saved oriented $S_3$ edges. Its singular-value ratios are invariant under family permutations,
unitary row/column phases, and an overall scale, so they give a necessary comparison before any
entrywise match is considered. The proposal requires
\[
s(T_0)/s_1=(1,0.55147\ldots,0.25517\ldots),\qquad
s(T_1)/s_1=(1,0.54625\ldots,0.27451\ldots).
\]
A $4{,}096$-patch BW run, an independent $4{,}096$-patch fusion run, a dense
$65{,}536$-patch population run, and a $65{,}536$-patch BW run contain $830{,}066$ edges in
total. Their corresponding triples are
\[
(1,0.00522,0.00134),\quad(1,0.00552,0.00361),\quad
(1,0.00179,0.00002),\quad(1,0.00071,0.00028).
\]
Source-node block bootstraps retain the same near-rank-one result. Every state
lies within the fixed \(0.02\) Haar-null tolerance and more than \(0.5407\)
from both proposed targets under the fixed \(0.05\) tolerance. The recorded
state contains only edge endpoints,
integer $S_3$ labels, and geometry: it exports no complex oriented family amplitude and no paired
coarse/fine edge intertwiner. Hence all direct $T_0$, refined $T_1$, and joint physical-emission
receipts are false. The conclusion is scoped to the natural direct carrier. A source-derived
complex lift could define a different carrier. Choosing its phases or paths to reconstruct the
proposed matrices would make the simulator test circular.

\paragraph{Maximal source-derived quark package.}
Let \(\mathcal P_q\) denote the three axioms together with the declared
finite matter, transport, and quark-source premises used in this section.
The quark-side result has four layers:
\begin{enumerate}[leftmargin=2em]
\item the emitted mass-profile ray
\[
\mathcal R_{ud}^{\mathrm{mass}}
=
\left\{
\lambda\left(\frac15,-\frac{1-x_2^2}{27}\right):\lambda\ge0
\right\},
\qquad
\lambda=t_1;
\]
\item the same-label left-handed selector value
\[
\sigma_{ud}=\sigma_{\mathrm{ref}},
\]
which is a negative sheet-selector statement;
\item the separate target-free mass bridge
\[
\mathcal P_q\vdash
\Delta_{ud}^{\mathrm{overlap}}=\frac16\log\frac{c_d}{c_u},
\qquad
\mathcal P_q\vdash
\Theta_{ud}^{\mathrm{mass}}=\frac16\log\frac{c_d}{c_u};
\]
\item the selected-class support wrapper on \(f_P\), which proves representative independence for
the attached physical spread datum. It does not select either modulus. Conditional mass readout
would also require an RG-covariant trajectory and a declared comparison chart; physical Yukawa
matrices require the additional common-scale dimensionless conversion.
\end{enumerate}

Separate same-family and common-refinement artifacts reproduce their chosen target rows exactly.
Their sigma datum is obtained by inversion of those rows. The associated
GeV-valued matrices are target-anchored mass textures. A separate calculation also backreads a
mass-side scalar after identifying coefficient ratios with target mass ratios. None of these
surfaces breaks the free source-spread action.

\subsection{Quark Theorem Boundary}\label{quark-family-boundary}

The conditional downstream closure route is
\[
f_P
\quad+\quad
\Sigma_{ud}^{\mathrm{source}}(P)
\Longrightarrow
\bigl(\sigma_u,\sigma_d,\sigma_{\mathrm{seed}}^{ud},\eta_{ud}\bigr)_{\mathrm{phys}}
\Longrightarrow
\bigl(g_u,g_d\bigr)
\Longrightarrow
\bigl(m_u,m_d,m_s,m_c,m_b,m_t\bigr),
\]
The first arrow is not identified by the source theory. The last arrow denotes scheme-labelled
mass coordinates only after an RG trajectory and comparison chart are specified. Physical
dimensionless Yukawa matrices lie one common-scale normalization beyond it. Selected-fiber
descent and global frame classification are logically separate from both obstructions.

The conditional matter lift and the conditional port-current algebra do not cut this fiber.
Both certificates are coefficient-blind: the lift fixes the gauge representations, the
chirality certificate, the anomaly traces, and exactly one Yukawa invariant line per declared
channel, the port-current receipt carries no Yukawa-adjacent datum, and independent positive
rescaling of the coefficients along the two hadronic invariant lines fixes every certified
conclusion of both. The spread fiber survives the certified structure set with its coordinates
realized as the free scalar coefficients along those lines. A machine-checked transport
certificate records the survival together with the exclusion of the twelve frozen orbit-selector
candidates, and flips fail-closed if any certified datum ever moves under the rescaling. A
cut of this fiber requires a physical binding of the response representation, the attachment
of the screen action to three physical families, or a selector admitted under the frozen
single-comparison discipline.

The same certificates permit and type a possible register relation at the unification scale;
they supply neither a physical equality between independent Yukawa coefficients nor a physical
generation order. The declared one-scalar
package has one invariant line on each of the down-conjugate and charged-lepton-conjugate
channels through the same scalar, together with a certified zero line on the channel that would
use the up scalar for the down-type field. Thus any separately supplied relation of this form
uses the down-type and charged-lepton channel pair rather than the forbidden up-scalar channel.
Conditional on the declared alphabet and the two transitive-color-orbit
constraints of measure balance and register faithfulness, exhaustive
enumeration selects the unordered weight multiset
\(\{1/3,\,1,\,3\}\) with no measured mass or angle in its solve path. A separate
comparison-only calculation of all six assignments ranks
\[
  (y_b/y_\tau,\;y_s/y_\mu,\;y_d/y_e)=(1,\;1/3,\;3)
\]
as uniquely least discrepant. That target-informed empirical ranking does not derive the
generation order from the source axioms. Conditional on the adopted assignment and the
charged-lepton triple, common
down-sector running cancels exactly on the declared one-loop branch, so the sharp
RG-protected result is
\[
  \frac{m_s}{m_d}=\frac{1}{9}
  \left.\frac{m_\mu}{m_e}\right|_{\mu_U}=22.9743.
\]
FLAG 2024 Tables~11--12 give \(m_s/m_{ud}=27.227(81)\) and
\(m_u/m_d=0.465(24)\) for \(N_f=2+1+1\), and \(27.42(12)\) and \(0.485(19)\)
for \(N_f=2+1\)~\cite{flag2024}. Their derived central values
\(m_s/m_d=19.9438\) and \(20.3594\) make the conditional route value \(15.2\%\) and
\(12.8\%\) high. Across all six assignments, the distinct light-family
coefficient-ratio menu is \(\{1/9,\,1/3,\,3,\,9\}\), and both FLAG rows
reject every assignment under the conservative
experimental-only comparison gate. The
unavailable input covariance and absent OPH theory uncertainty preclude a covariance-aware
significance or a theory-wide falsification conclusion. This result
applies only to the declared common-transport assignment family. Different coefficient relations,
coefficient alphabets, charged-family attachments, or generation-dependent threshold transport
define other hypothesis classes. The retained conclusions are the conditional pairing theorem,
the target-free unordered multiset under the declared measure-balance and
register-faithfulness constraints, and the exact
positive-chamber Koide identity. The stipulated charged-lepton response model
is neither blind nor source-derived.
The same conditional chart gives
\(m_b(m_b)=6.03\,\mathrm{GeV}\), \(m_s(2\,\mathrm{GeV})=140\,\mathrm{MeV}\), and
\(m_d(2\,\mathrm{GeV})=6.1\,\mathrm{MeV}\), respectively \(44.2\%\), \(50.3\%\), and
\(30.1\%\) above the recorded comparison coordinates. These absolute values
are not source-only results, and the third-generation normalization and
threshold packet are open.

\paragraph{Scoped residual-axis no-go.}
Hermitian circulants on the same \(C_3\) face fiber commute and have the same
discrete-Fourier eigenbasis. A direct same-fiber residual construction
therefore gives identity mixing. For a direct construction from two distinct
real three-dimensional residual axes, the six fivefold, ten threefold, and
fifteen twofold unoriented icosahedral axes give the exhaustive acute-angle
spectrum
\[
\begin{array}{@{}ll@{}}
\toprule
\text{axis families} & \text{acute angles in degrees}\\
\midrule
5\times5 & 63.4349\\
5\times3 & 37.3774,\ 79.1877\\
5\times2 & 31.7175,\ 58.2825,\ 90\\
3\times3 & 41.8103,\ 70.5288\\
3\times2 & 20.9052,\ 54.7356,\ 69.0948,\ 90\\
2\times2 & 36,\ 60,\ 72,\ 90\\
\bottomrule
\end{array}
\]
The smallest nonzero entry is \(20.9052^\circ\), with sine \(0.356822\).
The compare-only PDG 2024 \(K_{\mu2}\) coordinate
\(|V_{us}|=0.2250(4)\)~\cite{pdg2024ckm} corresponds to
\(\arcsin|V_{us}|=13.0029^\circ\) at its central value, which is absent. The enumeration therefore
excludes direct equality between the Cabibbo angle and an acute angle between
two of these 31 real axes. It does not exclude spinorial representations,
higher-order symmetry breaking, additional dynamical corrections, or general
overlap models. The \(31.7175^\circ=\arctan(1/\varphi)\) entry provides an
exact geometry self-check.

The register-Clebsch branch's \(\sqrt{m_d/m_s}=0.2086\) display is the
Gatto--Sartori--Tonin texture estimate formed from the same rejected mass
ratio. The branch supplies no up/down Yukawa-matrix pair or relative
left-handed eigenbasis, so the display is not an independent Cabibbo
prediction.

\subsection{Strong-CP branch}\label{quark-family-strong-cp-boundary}

The selected-class quark wrapper carries target-anchored mass textures on the public quark frame
class \(f_P\). Strong CP is a separate phase-side invariant. The stated source domain does not derive the
bare QCD angle \(\theta_{\mathrm{QCD}}\), does not emit the physical anomaly-invariant
combination \(\bar\theta\), and does not prove that the physical strong-CP phase vanishes.

This phase problem is independent of the source-spread and common-scale Yukawa obstructions. A
closure would require a source-emitted quark mass matrix at one declared scale, its physical
determinant-line phase contribution, and a theorem fixing the topological-angle contribution on the
realized branch. The present GeV mass textures do not supply that input.

\section{Charged-Lepton Family Derivation}

The charged-lepton derivation differs from the quark derivation in a precise way. It does not emit public charged masses from \(P\). It emits an exact same-family witness, a closed common-shift no-go, the declared same-label \(q_e\) readback, a determinant-line lift on theorem-grade physical charged data, and a downstream algebraic mass readout from theorem-grade \(A_{\mathrm{ch}}(P)\). For any fixed formal source exponent vector \(M_\bullet^{\mathrm{ch}}\), the same-label readback defines a source-side determinant character
\[
S_M=\sum_e M_e^{\mathrm{ch}}\log q_e.
\]
The theorem branch does not emit a theorem-grade sector-isolated charged determinant exponent vector, and it does not identify a source-side determinant character with the physical charged determinant line. The source-only theorem therefore supplies no electron, muon, or tau mass.

The mathematical split is equally precise. The charged-lepton derivation starts from
the ordered charged package, proves that the realized support is a one-dimensional linear subray,
exposes the canonical quadratic support-extension direction, maps that into the charged
excitation gaps, closes a two-scalar support-extension law shell, isolates the smaller eta
source-readback primitive on that same carrier, and then builds the log-spectrum and forward
shape/scale surface. Those centered objects are common-shift invariant. The
determinant line fixes the physical affine scalar once theorem-grade physical
charged data are present, and the mass readout is then algebraic. The charged
theorem boundary has two distinct pieces: physical selection of the latent
charged sector-response candidate as \(\widehat C_e\), and
source-to-physical determinant attachment. For a fixed formal source exponent
vector \(M_\bullet^{\mathrm{ch}}\), that attachment is the identity
\[
3\mu(r)=\sum_e M_e^{\mathrm{ch}}\log q_e(r),
\]
equivalently zero determinant-normalization defect
\[
N_{\det}(P)=s_{\det}(P)-\sum_e M_e^{\mathrm{ch}}\log q_e(P),
\]
on the charged determinant channel.

Physically, the electron is the light charged lepton that makes atoms and chemistry possible, the
muon is its heavier unstable cousin, and the tau lepton is the heaviest charged lepton. A physical
charged-lepton derivation would have to explain how those three rows emerge from the shared flavor
dictionary without Koide-assisted fitting. The theorem-grade branch must select
the latent charged sector-response candidate \(\widehat C_e^{\mathrm{cand}}\)
as \(\widehat C_e\) through the branch-generator splitting theorem. On theorem-grade physical
\(Y_e\), a refinement-stable uncentered lift collapses the determinant-line section and affine
anchor to one descended physical affine scalar \(\mu_{\mathrm{phys}}(Y_e)\), with
\[
\widetilde C_e(Y_e)=\widehat C_e(Y_e)+\mu_{\mathrm{phys}}(Y_e)\,\mathbf 1,
\qquad
s_{\det}(Y_e)=3\mu_{\mathrm{phys}}(Y_e),
\qquad
A_{\mathrm{ch}}(Y_e)=\mu_{\mathrm{phys}}(Y_e).
\]

\begin{theorem}[Charged same-carrier source-pair readback]
Fix the ordered charged carrier
\[
\begin{gathered}
(-1,x_2,1),\\
x_2=-0.5175863354681689.
\end{gathered}
\]
If the charged source pair
\[
(\eta_{\mathrm{ext}},\sigma_{\mathrm{ext}})
=
(\eta_{\mathrm{source\_support\_extension\_log\_per\_side}},
\sigma_{\mathrm{source\_support\_extension\_total\_log\_per\_side}})
\]
is emitted on that carrier, then the centered charged logs are
\[
e_{\log,\mathrm{centered}}
=
-\frac{(3+x_2)\sigma_{\mathrm{ext}}-\eta_{\mathrm{ext}}}{6},
\]
\[
\mu_{\log,\mathrm{centered}}
=
\frac{x_2\sigma_{\mathrm{ext}}-\eta_{\mathrm{ext}}}{3},
\]
\[
\tau_{\log,\mathrm{centered}}
=
\frac{(3-x_2)\sigma_{\mathrm{ext}}+\eta_{\mathrm{ext}}}{6},
\]
and therefore the charged masses are
\[
m_e = g_e\,e^{e_{\log,\mathrm{centered}}},
\qquad
m_\mu = g_e\,e^{\mu_{\log,\mathrm{centered}}},
\qquad
m_\tau = g_e\,e^{\tau_{\log,\mathrm{centered}}}.
\]
\end{theorem}

This theorem is exact on the same-carrier shell, but the absolute values are not emitted.
The visible scalar order is
\[
\eta_{\mathrm{ext}}
\quad\text{then}\quad
\sigma_{\mathrm{ext}},
\]
and the charged absolute scale \(g_e\) is unresolved. The scalars
\(\eta_{\mathrm{ext}}\) and \(\sigma_{\mathrm{ext}}\) are the first
same-carrier residuals. If the latent candidate
\(\widehat C_e^{\mathrm{cand}}\) is selected as the physical operator, then
\(\eta_{\mathrm{ext}}\) and
\(\sigma_{\mathrm{ext}}\) become charged spectral invariants instead of separate primitive goals,
and the absolute-scale burden is pushed to one affine-covariant absolute charged anchor
\(A_{\mathrm{ch}}\). In the
local chain, \(\widehat C_e\) itself is undeclared. The available object is
only the centered compressed generation-bundle candidate
\(\widehat C_e^{\mathrm{cand}}\). Its missing theorem is the
compression-descendant commutator condition after the central split.
The centered common-shift quotient is closed negatively, so centered data alone do not
emit the affine anchor \(A_{\mathrm{ch}}\).

\begin{theorem}[Charged absolute-scale underdetermination]
Let
\[
E_e^{\mathrm{centered}}
=
\bigl(e_{\log,\mathrm{centered}},\mu_{\log,\mathrm{centered}},\tau_{\log,\mathrm{centered}}\bigr)
\]
be the centered charged log triple emitted from the charged source pair
\((\eta_{\mathrm{ext}},\sigma_{\mathrm{ext}})\). Then for every \(c\in\mathbb R\),
\[
Y_e(c):=\exp(c)\,\mathrm{diag}\!\bigl(e^{E_e^{\mathrm{centered}}}\bigr)
\]
has the same charged spectral invariants
\[
\eta_{\mathrm{ext}},\qquad \sigma_{\mathrm{ext}},\qquad \gamma_{21},\qquad \gamma_{32},
\]
the same centered log vector, and the same ratio data. Only the absolute masses scale:
\[
(m_e,m_\mu,m_\tau)\longmapsto e^c\,(m_e,m_\mu,m_\tau).
\]
Hence the charged OPH chain determines only the quotient class
\[
E_e^{\mathrm{centered}}\in \mathbb R^3/\langle(1,1,1)\rangle,
\]
not the absolute scalar \(g_e=e^c\).
\end{theorem}

The proof is immediate from the centered sum rule
\[
e_{\log,\mathrm{centered}}+\mu_{\log,\mathrm{centered}}+\tau_{\log,\mathrm{centered}}=0,
\]
which implies
\[
\det Y_e^{\mathrm{shape}}=1,
\qquad
\det Y_e = g_e^3.
\]
All emitted charged invariants depend only on differences of log entries, so a common shift in the
\((1,1,1)\) direction leaves them unchanged. The stated source domain
therefore gives a no-go for the \(P\)-driven charged mass: the theorem fixes
the centered shape, while the determinant-line landing fixes the absolute
normalization.

The charged absolute-scale
branch is explicitly typed. The charged scale is a linear quantity
\[
g_e = e^{\mu_e^{\mathrm{abs}}},
\]
so the log-coordinate \(\mu_e^{\mathrm{abs}}\) must not be mixed directly with centered log gaps.
The same-family writeback therefore records the type-consistent shell
\[
\begin{aligned}
\mu_{e,\mathrm{seed}}^{\mathrm{abs}}
&= \log(0.9231656602589082)\\
&= -0.07994658034676537,
\end{aligned}
\]
\[
\mu_{e,\mathrm{cand}}^{\mathrm{abs}}
=
\mu_{e,\mathrm{seed}}^{\mathrm{abs}}-\gamma_{\min}
=
-0.38231224060567365,
\]
\[
g_e^{\mathrm{cand}}
=
e^{\mu_{e,\mathrm{cand}}^{\mathrm{abs}}}
=
0.6822819838027987.
\]
This is a representation-consistency shell only, not a charged-mass theorem. It fixes the
linear-vs-log coordinate discipline. This surface does not emit a theorem-grade value law for \(g_e\).
The checked shortcut
\[
\begin{aligned}
\Delta_e^{\mathrm{abs}}&=0.30236566025890826,\\
g_e&=0.6822819838027987
\end{aligned}
\]
is not a theorem-grade closure. It merely chooses one representative on the common-shift orbit and
does not land on the physical charged masses.

A physical completion requires a theorem-grade affine-covariant section
\(A_{\mathrm{ch}}\) of the quotient map, satisfying
\[
A_{\mathrm{ch}}(\ell + c\,\mathbf 1)=A_{\mathrm{ch}}(\ell)+c.
\]
The clean realization of that section is an uncentered charged response lift carrying a
determinant line, in which
\[
A_{\mathrm{ch}}=\tfrac13 \log \det(Y_e)=\tfrac13 \operatorname{tr}(\log Y_e).
\]
Once such an anchor exists,
\[
g_e = e^{A_{\mathrm{ch}}(\ell)},
\qquad
\Delta_e^{\mathrm{abs}} = \log g_{\mathrm{ch}}^{\mathrm{shared}} - A_{\mathrm{ch}}(\ell).
\]

\begin{proposition}[Twenty-four-register rate no-go]
An oriented register with twenty-four slots, a twenty-four-stage automaton, or
a decomposition into twenty-four formal update steps determines no physical
frequency, Hamiltonian gap, or mass scale. Indeed, rescaling every physical
generator by \(H\mapsto\lambda H\), \(\lambda>0\), preserves the register,
transition graph, grading, and settled quotient while multiplying all
generator gaps and rates by \(\lambda\). A repair iteration counter is likewise
not an operational clock. A physical rate requires a source-derived clock
instrument, calibrated comparison line, and residual bound.
\end{proposition}

\begin{theorem}[Conditional charged determinant-clock attachment]
Let \(L_{24}\) be a one-dimensional source-derived operational clock line
whose calibrated positive gap has norm \(\Delta_{24}\). Suppose there is a
refinement-natural, quotient-visible norm-preserving line isomorphism
\[
\Theta_e:L_{24}^{\otimes3}\longrightarrow\det M_e
\]
to the physical charged-lepton determinant line. Then the common charged
mass scale is fixed by
\[
g_e:=|\det M_e|^{1/3}=\Delta_{24}.
\]
More generally, a separately derived normalization factor in \(\Theta_e\)
propagates as its cube root. Thus the implication from a calibrated clock-line
attachment to the determinant scale is exact.
\end{theorem}

\begin{proof}
The tensor-cube norm is \(\|L_{24}^{\otimes3}\|=\Delta_{24}^3\).
Norm preservation gives \(|\det M_e|=\Delta_{24}^3\); take the positive cube
root.
\end{proof}

The register count does not construct \(L_{24}\), calibrate
\(\Delta_{24}\), or produce \(\Theta_e\). Those are the removable physical
clock, determinant-descent, and normalization gates. The determinant scale
is work in progress. Koide balance fixes a shape condition and is
independent of this affine scale attachment.

\paragraph{Positive-chamber face-circulant identity and conditional finite-GNS balance.}
The regular \(C_3\) face fiber has the general Hermitian carrier
\[
C=aI+\rho\left(e^{i\delta}R+e^{-i\delta}R^2\right),
\qquad a,\rho\geq0,
\]
where \(R\) is the regular cyclic shift, \(R^3=I\). The commutant of the regular \(C_3\)
action is spanned by \(I,R,R^2\), which makes \(C\) the general Hermitian
\(C_3\)-equivariant carrier. Fourier diagonalization gives
\[
r_k=a+2\rho\cos\!\left(\delta+\frac{2\pi k}{3}\right),
\qquad k=0,1,2.
\]
In a positive spectral chamber, \(m_k=s r_k^2\) and
\(\sqrt{m_k}=\sqrt{s}\,r_k\). The roots-of-unity identities
\[
\sum_{k=0}^2\cos\!\left(\delta+\frac{2\pi k}{3}\right)=0,
\qquad
\sum_{k=0}^2\cos^2\!\left(\delta+\frac{2\pi k}{3}\right)=\frac32
\]
give the exact Koide identity
\[
Q=\frac{\sum_km_k}{(\sum_k\sqrt{m_k})^2}
 =\frac{1+2(\rho/a)^2}{3},
\]
so \(\rho/a=1/\sqrt2\) is exactly equivalent to \(Q=2/3\) in the positive-eigenvalue
chamber. At balance, positivity holds for
\(|\delta|\leq\pi/12\pmod{2\pi/3}\), and throughout that chamber Koide is independent of
\(\delta\). Outside it, physical square roots are absolute eigenvalues and the physical Koide
ratio differs from the signed trace expression. Equivalently, if
\(E_0=3a^2\) and \(E_c=6\rho^2\) are the singlet and charged-plane Hilbert--Schmidt powers, then
\[
Q=\frac{1+E_c/E_0}{3},
\qquad
Q=\frac23\ \Longleftrightarrow\ E_c=E_0.
\]
The positive square-root-mass vector lies on the \(45^\circ\) Koide cone about the democratic
direction. Circulant symmetry leaves the relative norm of the singlet and charged plane free.

The finite response packet sets that norm under its declared hypotheses. Let
\[
\mathcal V_K=\mathbf1\oplus\chi\oplus\bar\chi,
\qquad \mathcal H_{\rm or}=\mathbb C^2,
\qquad
\mathcal A_K=B(\mathcal V_K\otimes\mathcal H_{\rm or})\simeq M_6(\mathbb C).
\]
A one-state record cannot distinguish orientation reversal, so \(\mathbb C^2\) is the minimal
faithful orientation record. Minimality of this response-local register is an
explicit declared hypothesis. With \(q_+\) one oriented outcome, define
\[
E_+=P_0\otimes I_{\rm or}+P_c\otimes q_+,
\qquad
Z_0=P_0\otimes I_{\rm or},
\qquad
Z_c=P_c\otimes q_+.
\]
The blocks \(Z_0,Z_c\) both have rank two. Born--L\"uders conditioning of the unique MaxEnt
state \(I_6/6\) on \(E_+\) gives
\[
p_0=p_c=\frac12.
\]
The orientation-blind singlet has event action zero, while the selected charged orientation has
action \(\ln2\). Charged multiplicity two times its weight \(1/2\) equals the singlet weight.
The unique normalized trace on \(M_6(\mathbb C)\) fixes these probabilities without the free trace
parameter present on \(\mathbb C\oplus M_2(\mathbb C)\).

In \(L^2(B(\mathbb C^3),\tau_3)\), the operators \(I,R,R^2\) are orthonormal, and
\[
\mathcal J(e_0)=I,
\qquad
\mathcal J(e_+)=R,
\qquad
\mathcal J(e_-)=R^2
\]
is unitary and intertwines the \(C_3\) action, orientation reversal, and the two block projections.
The canonical GNS square-root amplitude is
\[
\xi=\sqrt{p_0}\,e_0+\sqrt{\frac{p_c}{2}}
\left(e^{i\delta}e_++e^{-i\delta}e_-\right),
\]
so its response has \(a=\sqrt{p_0}\) and \(|b|=\sqrt{p_c/2}\). Hence
\[
\frac{|b|^2}{a^2}=\frac{p_c}{2p_0}=\frac12,
\qquad
\frac{|b|}{a}=\frac1{\sqrt2},
\qquad Q=\frac23.
\]
An orientation action gap \(\Delta S\) gives the general form
\[
\frac{p_c}{p_0}=2e^{-\Delta S},
\qquad
\frac{|b|^2}{a^2}=e^{-\Delta S},
\qquad
Q(\Delta S)=\frac{1+2e^{-\Delta S}}{3}.
\]
For comparison, the target-informed minimal complete public response
coordinate has
\[
 Q_{\mathrm{MCPR}}=0.6666644634090389,\qquad
 \frac{|b|}{a}=0.7071044442750720,
\]
whose modulus is \(3.30\) ppm below \(1/\sqrt2\). The PDG 2026 central
charged-lepton masses give
\[
 Q_{\mathrm{PDG}}=0.6666644634026367.
\]
The numerical proximity is diagnostic. The response architecture and phase
were stipulated with target knowledge, so this comparison is not a
source-only prediction~\cite{pdg2026}.

A physical Koide implication is conditional on faithful trace-preserving u.c.p. maps \(\Phi\)
and \(\Psi\) from the source process to a physical response process and back, with
\(\Psi\Phi=\mathrm{id}\). Kadison--Schwarz applied to both maps puts every source element in the
multiplicative domain of \(\Phi\), so \(\Phi\) is a \(*\)-monomorphism and an exact \(L^2\)
isometry. If it also intertwines \(Z_0,Z_c\) with physical response records and sends the
conditioned source state to the normalized positive square-root-mass response
\(C=(Y_e^\dagger Y_e)^{1/4}\), then it transports \(p_0=p_c\) to \(E_0=E_c\) and proves physical
Koide. For an accepted finite chiral checkpoint with three left and three right family modes,
positive kinetic metrics, and a committed neutral Higgs direction, the reversible-checkpoint
conditions force \(\Phi=\operatorname{Ad}_{J_L\oplus J_E}\), preserving chirality and the
regular \(C_3\) action. If \(M_F=X_F^2\) is the positive source response, then
\[
\widehat Y_e=\frac{\sqrt2}{v}J_LM_FJ_E^\dagger,
\qquad
\mathcal M_L=J_LM_FJ_L^\dagger.
\]
The source and canonical physical square-root responses have equal block
powers. This closes the finite balanced shape implication on supplied
physical charged data, not the determinant scale. The extended branch \(\mathrm{OPH}^{+}_{\rm ch}\) adds graded
physical completion and quotient source-law selection. Given an exhaustive declared carrier class with
a positive winner gap, a source-closed BV/BRST continuum, and an interval or contraction
certificate for the charged QFT self-map, it selects one accepted charged sector and one dressed
mass readout. The stable electron mass is a dressed spectral lower edge; the muon and tau masses
are real parts of specified dressed resonance roots. A balanced \(C_3\) fixed point gives
\(Q=2/3\). A \(C_3\)-symmetric fixed point with charged attenuation \(\chi_\star\) gives
\[
Q=\frac{1+e^{-2\chi_\star}}{3}.
\]
An off-plane response requires the full response operator. These extended-branch principles and
the QFT certificate are additional conditions beyond the three axioms.
The relation \(\operatorname{Hom}_{C_3}(\mathbf1,\chi\oplus\bar\chi)=0\) prevents the neutral
source grammar from producing a labeled charged-family vector through symmetry. A charged source
tensor, connection, or selected orbit must supply the phase and individual ratios.

The count \((N_c+1)/(2N_cN_g)=2/9\) is arithmetic, and the value \(2/9\) is the
empirical Brannen--Koide fitted phase~\cite{koide1983,brannen2006}; the integer
factorization is conditioned on that value. A physical phase theorem requires a
regular-\(C_3\) family bundle, an oriented family connection
whose declared loop has holonomy
\(\exp[i\beta_{\rm EW}/(2N_cN_g)]\), and an attachment identifying that holonomy with the phase of
the charged Fourier coefficient. Hypercharge acts as a common scalar on the generation copies,
while the shift \(R\) permutes them. Pure finite \(A_5/C_3\) holonomy yields cube-root phases.
Construction of the required continuous family connection and its source selector is work in
progress. Under the stated connection and attachment hypotheses, \(\delta=2/9\). The phase of one
equal link differs by a factor of three from the total triangular holonomy.
With the scale-free normalization \(a=1\), the ordered roots
\[
r_k = a + 2\rho\cos\!\left(\delta + \frac{2\pi k}{3}\right)
\]
become
\[
\begin{aligned}
r_e&=0.040349908219207475,\\
r_\mu&=0.5802119201475368,\\
r_\tau&=2.3794381716332555.
\end{aligned}
\]
Under those assumptions the bridge selects the same-carrier pair
\[
\begin{aligned}
\eta_{\mathrm{ext}}&=-6.729586682888832,\\
\sigma_{\mathrm{ext}}&=8.154061112725994,
\end{aligned}
\]
with ordered gap values
\[
\begin{aligned}
\gamma_{21}&=5.33160859254774,\\
\gamma_{32}&=2.822452520178255,\\
\kappa_{\mathrm{ext}}&=-4.59605680397234,
\end{aligned}
\]
and centered logs
\[
E_{\log,\mathrm{centered}}
=
[-4.495223235091244,\ 0.836385357456495,\ 3.6588378776347503].
\]
Against the charged references, this continuation-centered shape has residual norm
\[
\left\|E_{\log,\mathrm{centered}}^{\mathrm{cont}}
-E_{\log,\mathrm{centered}}^{\mathrm{ref}}\right\|
\approx
2.13\times10^{-5},
\]
so this compare-only branch is a near-exact centered-shape closure up to one common absolute
scale. The ppm-scale shape agreement carries a large pull in measurement units: the same
\(\delta=2/9\) ansatz misses the measured \(m_\mu/m_e\) ratio by \(9.83\)~ppm, approximately
\(440\sigma\) of the measurement precision, so the row is excluded as a prediction and stands as
an approximate fit. The theorem branch is unpromoted because the public affine absolute anchor is external
to this near-exact centered charged-shape branch. The compare-only common scale
required for exact absolute masses is
\[
g_e^\star = 0.04577885783568762.
\]
Equivalently, relative to the stored shared-budget seed
\[
g_{\mathrm{ch}}^{\mathrm{shared}} = 0.9231656602589082,
\]
the missing affine absolute anchor on the determinant-line route would have to take the target value
\[
\Delta_e^{\mathrm{abs},\star}
=
\log\!\frac{g_{\mathrm{ch}}^{\mathrm{shared}}}{g_e^\star}
=
3.003986333402356.
\]
This identifies the charged theorem boundary. The compare-only branch nearly
solves the centered charged shape. The derivation does not select the latent
candidate \(\widehat C_e^{\mathrm{cand}}\) as the physical
\(\widehat C_e\), and it lacks one
affine-covariant absolute anchor \(A_{\mathrm{ch}}\) that would turn that centered readback into
public charged masses on the theorem branch.

\paragraph{Icosahedral face-incidence carrier.}
The screen geometry supplies a sharper candidate carrier than a scalar
twelve-port moment. The twenty outward-oriented icosahedral faces form
\(A_5/C_3\), and the face stabilizer cyclically permutes the three corners.
Thus each face has a canonical local regular-\(C_3\) fiber, and an equivariant
Hermitian circulant has a face-representative-independent unordered spectrum.
This is an exact geometric lemma. It is not a physical generation
attachment: the sixty face-corner flags form the regular \(A_5\) torsor, not
one canonical three-dimensional matter-family space.

A conditional theorem package declares the diagonal affine repair map
\[
T(\kappa,\chi_\rho,\zeta_\delta)
=
(s_\kappa-q_\kappa\kappa,\,
 s_\chi+q_\chi\chi_\rho,\,
 s_\zeta+q_\zeta\zeta_\delta).
\]
For that declared map it proves
\(\lVert DT\rVert_\infty=0.0015356510519\ldots<1\), so Banach's theorem closes
existence, uniqueness, and iterative stability of its fixed point. This is a
conditional closure, not a derivation of the repair dynamics. A
source-multiplier family \(T_{\boldsymbol\lambda}\) with the
same displayed symmetry, analytic degree, Jacobian, and contraction but
different masses exists. It proves scoped selector non-identifiability
under those properties, not non-entailment from every OPH axiom. Exact
zeroth-order block balance gives
\(|b|/a=1/\sqrt2\), whereas the endpoint operator uses
\(|b|/a=e^{-\chi_\rho}/\sqrt2\); a source-derived bare-to-endpoint repair bridge
is work in progress.

\begin{proposition}[Finite charged-register realization]
\label{prop:charged-cfq-digital-model}
The stipulated charged register-and-path class is nonempty. There is an explicit
finite algebraic model with register dimensions
\[
(50,31,10,512,77,21,27,5),
\]
connected transition graphs, normalized tracial states, rank-one event
projectors, and the eight declared path weights
\[
\left(\frac1{50},-\frac1{31},-\frac1{310},\frac1{512},
\frac1{77},\frac1{21},\frac1{27},\frac1{135}\right).
\]
Each noncentral event in \(B(H_r)\) admits an accepted/rejected central record
dilation in \(B(H_r)\otimes D_2\). The graph family is covariant under the
sixty proper icosahedral rotations, and normalized path weights are unchanged
under inert ancillary stabilization \(X\mapsto X\otimes I_k\).
\end{proposition}

\begin{proof}
For each declared connected graph, the diagonal matrix units and both directed
units on every edge generate \(E_{ij}\) along graph paths and hence generate
the full matrix algebra. The normalized trace of a rank-one event is the
reciprocal register dimension, and tensor products multiply the two depth-two
weights. For an event \(P\), the pinching map
\[
\mathcal E_P(X)=PXP+(I-P)X(I-P)
\]
is a unital trace-preserving conditional expectation. Its instrument writes
\(P\rho P\) and \((I-P)\rho(I-P)\) into orthogonal accepted/rejected pointers in
the central \(D_2\) factor. Explicit permutation intertwiners preserve graph
incidence and ranks for all sixty proper rotations. Finally,
\(\tau_{dk}(P\otimes I_k)=\tau_d(P)\), and the partial-trace coarse map retracts
the inert embedding. These facts construct one model of the declared schema.
\end{proof}

Proposition~\ref{prop:charged-cfq-digital-model} proves formal nonemptiness.
It supplies no derivation of the charged source. The eight dimensions, path
automaton, coupling degrees, grading, signs, clock, and scalar response are
model inputs. The verifier exhausts that automaton, not the physically
admissible OPH path space. The construction resolves the abstract
event-versus-central-record issue at fixed cutoff. Physical response
selection, cofinal screen refinement, and charged-sector attachment are work
in progress.

\begin{proposition}[Conditional charged nature and singularity transport]
\label{prop:charged-nature-pole-transport}
Let \(M_F\) be the positive face endpoint. Suppose a physical chiral
three-family carrier supplies a natural unitary \(J_L\), canonical kinetic
metrics, and
\[
 \frac{v^2}{2}\widehat Y_e\widehat Y_e^\dagger
 =J_LM_F^2J_L^\dagger.
\]
Then its positive left charged response is \(J_LM_FJ_L^\dagger\). Suppose in
addition that the exact renormalized charged kernel is supplied, its singular
lines satisfy the declared infrared and regularity conditions, and a
finite-register/Dyson map has singularity readout
\(J_LM_FJ_L^\dagger\). Regular
invertible field changes preserve the zero set, and a regular Nielsen
factorization \(\partial_\xi K=A_\xi K+KB_\xi\) makes each simple zero gauge
independent.
\end{proposition}

\begin{proof}
The first statement follows from uniqueness of the positive square root.
Multiplication of \(K\) on the left and right by analytic invertible matrices
multiplies its determinant by nonvanishing factors. Jacobi's identity gives
\[
 \partial_\xi\det K=(\operatorname{tr}A_\xi+
 \operatorname{tr}B_\xi)\det K,
\]
so a simple zero cannot move while the Nielsen factors are regular. These are
the standard mixed-fermion and gauge-independence implications used in the
pole literature~\cite{kniehl2013,gambino1999}; the stable charged-line
infrared qualification follows the infraparticle boundary
\cite{buchholz1986,duch2023}.
\end{proof}

Proposition~\ref{prop:charged-nature-pole-transport} closes the downstream
logic after its physical premises are supplied. It does not supply them.
The displayed square identity is the desired operator attachment, and the
singularity premise sets the Dyson readout equal to the face response.
The explicit kernel \(K_0(s)=sI-M_F^2\) has zero self-energy; it is a
free algebraic existence witness rather than the interacting charged-lepton
one-particle-irreducible kernel. The physical carrier attachment, interacting
kernel, and singularity transport are work in progress.
The verifier reconstructs the fixed point and response shape but does not check
\(M_F=g_{\rm end}S_F\): a coherent change
of the ordered spectrum by factors \((2,1/2,1)\), together with the matching
mass matrix, Yukawa matrix, and pole roots, preserves the determinant and
satisfies the tested mass relations. The package is therefore a conditional
conditional theorem. It does not identify physical charged matter, construct
an interacting pole, or predict a mass.

These results do not yield physical charged-lepton masses. Such a conclusion
requires a quotient-visible face-to-charged-family
intertwiner, the \(\ln2\) MaxEnt and ensemble-to-Yukawa bridge, a charged
connection deriving both the \(2/9\) base phase and its correction, a normalized
\(\mathbb Z_6\) determinant character, an OPH dynamical theorem selecting the
physical registers, state, exhaustive path category, grading, clock, and
response rather than declaring them, cofinal refinement naturality,
one target-independent coherent source tuple, and a mass-scheme map. The
icosahedral face carrier, face-incidence theorem, and conditional
finite-quotient weights supply only the stated finite ingredients.

\subsection{Absolute charged-lepton mass intervals from electromagnetic transport}\label{charged-kappa-interval}

The missing absolute anchor is one real scale. Fixing \(A_{\mathrm{ch}}\) is equivalent to fixing
the single family rescaling \(Y_e\mapsto e^{\kappa}Y_e\), under which every charged mass scales as
\(m_i\mapsto e^{\kappa}m_i\) while the centered shape and every mass ratio are invariant. The
declared charged antecedents are invariant under that rescaling, which is the content of the
common-shift no-go: the gauge representations, the anomaly structure, the complexity vector, the
determinant algebra, and the same-label readback take the same value on the whole family
\(\{e^{\kappa}Y_e\}\), so no source object built from those antecedents alone selects \(\kappa\).

The rescaling is not a symmetry of the declared electromagnetic transport. The charged leptons
enter the photon vacuum polarization, and the leptonic packet in the fine-structure endpoint map
carries
\[
P_{\mathrm{lep}}(\kappa)=\frac{1}{3\pi}\sum_{i}\left[2\log\frac{M_Z}{m_i}-\frac53\right],
\qquad
\frac{\mathrm{d}P_{\mathrm{lep}}}{\mathrm{d}\kappa}=-\frac{2}{\pi}.
\]
The declared endpoint pixel map \(P=\varphi+\sqrt\pi/A_{\mathrm{Th}}(P)\) consumes \(P_{\mathrm{lep}}\),
so the emitted endpoint and the pixel move with \(\kappa\). The common-shift no-go omits this
transport from its antecedent list. Source-only, the system is the stiff curve \((P(\kappa),\kappa)\)
with \(|\mathrm{d}P/\mathrm{d}\kappa|\approx6\times10^{-5}\), and the absolute masses are work in
progress on the source branch.

On the empirical-closure branch the transport identifies \(\kappa\). Inverting the on-shell
decomposition
\[
a_0+g=A^{-1}(0)\left(1-P_{\mathrm{lep}}(\kappa)-\Delta_{\mathrm{had}}-\Delta_{\mathrm{top}}\right)
\]
with the frozen anchor \(a_0=128.308\), the certified anchor-gap interval \(g\in[0.620,\,0.651]\),
the hadronic payload \(\Delta_{\mathrm{had}}=0.027609\pm0.000112\) pinned to the published
data-driven compilation of Keshavarzi, Nomura and Teubner (Phys.\ Rev.\ D 101, 014029), the
target-anchored mass ratios,
and the measured \(A^{-1}(0)\) supplied as a compare-only exclusion inside the interval solve path
gives
\[
\kappa\in[-0.070,\,+0.061],\qquad \kappa_{\mathrm{central}}=-0.004.
\]
An independent 100-digit decimal calculation brackets \(\pi\), pads the
transcendental evaluation, and rounds every endpoint outward. The resulting
target-anchored diagnostic intervals are
\[
\begin{aligned}
m_e&=0.5089~[0.4765,\,0.5434]~\mathrm{MeV}, & &\text{measured } 0.5110~\mathrm{MeV},\\
m_\mu&=105.22~[98.5,\,112.4]~\mathrm{MeV}, & &\text{measured } 105.66~\mathrm{MeV},\\
m_\tau&=1.7695~[1.657,\,1.890]~\mathrm{GeV}, & &\text{measured } 1.7769~\mathrm{GeV},
\end{aligned}
\]
with central values within about \(0.4\%\) of measurement and the physical triple inside every
interval. The physical on-shell anchor requirement lies inside the certified gap interval. The
solve also inverts exactly at the witness: the anchor-gap value \(g=0.6379\) closes the branch on
the measured triple, and its distance \(+0.0070\) from the standard on-shell reference deficit
\(0.631\) is the scheme correction required by the bridge. A source-emitted bridge value is a sharp
falsification target: a value at \(0.6379\) within the payload width closes the lepton branch on
the witness, and a value outside the certified interval refutes the decomposition. The
interval width reduces to the published payload uncertainty and the anchor
gap. A source derivation requires the hadron backend and the scheme bridge.

The certified anchor-gap interval is the exact affine image of the hadronic payload interval,
\(g(X)=A^{-1}(0)(1-X)-A_{\mathrm{lep}}-a_0\), so the rectangle solve above carries the one payload
uncertainty in two anti-correlated slots. On the payload-coherent reading, where the gap is
evaluated at the same payload value as the hadronic term, the payload cancels from the
decomposition and the surviving width is the higher-order leptonic remainder and the kernel
truncation. The coherent central coincides with the rectangle central, and the certified width
contracts by a further factor of \(3.8\):
\[
\begin{aligned}
m_e&=0.5089~[0.5001,\,0.5178]~\mathrm{MeV}, & &\text{measured } 0.5110~\mathrm{MeV},\\
m_\mu&=105.22~[103.40,\,107.06]~\mathrm{MeV}, & &\text{measured } 105.66~\mathrm{MeV},\\
m_\tau&=1.7695~[1.7390,\,1.8004]~\mathrm{GeV}, & &\text{measured } 1.7769~\mathrm{GeV},
\end{aligned}
\]
with the measured triple inside every coherent interval. The coherent
logarithmic half-width is \(0.01732\), corresponding to multiplicative
one-sided widths of \(-1.72\%\) and \(+1.75\%\). The rectangle interval remains the
independent-gap statement for every bridge value inside its certified range;
the coherent row is conditional on the declared payload-coherent anchor-gap
premise. Both intervals require the same hadron backend and scheme bridge for
source closure. The surviving coherent width is a
premise floor rather than a budget slack: the higher-order band matches the published per-order
structure of the leptonic running with negligible \(\kappa\)-sensitivity across the certified
interval, so the certified floor is the scheme-bridge ambiguity itself, and a tighter certified
width requires the source bridge rather than a smaller budget.

The determinant prescription is a consistency test on this scale. It reproduces
the physical leptonic packet at \(|\kappa|\approx0.006\). Its runtime has no
continuous fit parameter, while its architecture and charged ratios retain
target ancestry. The construction adds no axiom. The interval branch carries
absolute masses only on the empirical-closure branch; the source-only charged
no-go is in force.

\subsection{Conditional response candidate and symmetry-breaking boundary}\label{charged-mcpr-frontier}

A conditional package defines the charged response by
stipulating a typed, incidence-complete response architecture on an oriented
icosahedral flag: eight finite registers of dimensions
\(50,31,10,512,77,21,27,5\), eight primitive path classes with normalized
rank-one trace weights, a public-block amplitude, an oriented process phase,
and a \(\mathbb Z_6\) determinant character. Conditional on that
architecture, the three-channel response map is a Banach contraction with
contraction constant \(1.5\times10^{-3}\); its unique fixed point, the
regular-\(C_3\) shape functional, and the determinant character emit one
dimensionless charged triple on the declared-architecture branch with zero
runtime charged-mass input:
\[
\frac{(m_e,m_\mu,m_\tau)}{E_\star}
=(4.18511\times10^{-23},\,8.65348\times10^{-21},\,1.45532\times10^{-19}),
\]
with mass ratios \(m_\mu/m_e=206.7683\), \(m_\tau/m_e=3477.3655\),
\(m_\tau/m_\mu=16.8177\), and unclosed-clock displays
\((0.51096,\,105.649,\,1776.78)\)~MeV. A segregated compare-only calculation
places this coherent branch about \(84\) ppm below the comparison
triple; the candidate computation has no dependency on the comparison file.

The register dimensions, path table, signs, amplitude, phase, and determinant
exponent are declared model inputs, so the triple is a declared-model
coordinate. Those architecture choices are target-informed and stipulated
with knowledge of the charged-family target. Their executable evaluation is
runtime-reference-free, and the candidate is neither blind nor source-only.
The \(84\) ppm comparison is diagnostic. The Koide subsection derives the finite
public-block amplitude \(1/\sqrt2\) from the connected \(M_6\) response register.
The candidate supplies no recoverable attachment between that event and a physical chiral mass
channel, so its mass claim is conditional. The arithmetic identity
\((N_c+1)/(2N_cN_g)=2/9\) supplies no oriented phase; a holonomy
reading requires link variables, a loop, and a nonlocal readout map.
A product of fourteen normalized rank-one traces is neither a
\(\mathbb Z_6\) character nor a determinant, and the canonical determinant
norm carries the kinetic factor
\(\det\mathcal M_L=(v/\sqrt2)^3\,|\det Y_e|/\sqrt{\det Z_L\det Z_E}\); the
declared \(6^{-14}\) weight is a candidate positive path weight.

The source-side problem is a finite
symmetry-breaking analysis on the icosahedral carrier. The vertex permutation
module decomposes as
\(\mathbb R^{12}\cong W_1\oplus W_3\oplus W_3'\oplus W_5\); the traceless
quadrupole map \(Q(w)=\sum_iw_i(p_ip_i^T-I/3)\) satisfies \(Q=QP_5\) and
\(Q^*Q=\frac85P_5\), so it sees exactly \(W_5\), and the two adjacent gaps of
its spectrum span the centered family plane. A unique invariant MaxEnt state
has zero expectation on every nontrivial irreducible module, and the
homogeneous twelve-port branch has the unique uniform minimizer, so charged
family shape requires a source-selected \(W_5\) orbit. Because
\(\operatorname{Sym}^2_0(\mathbf3)\cong\operatorname{Sym}^2_0(\mathbf3')\cong W_5\),
the family attachment is multiplicity-one up to one sign once the chiral
family fibers carry a three-dimensional \(A_5\) representation.

Residual symmetry alone does not select the mass ratios. Realize \(W_5\) as
the traceless symmetric real \(3\times3\) matrices. Invariance under a
threefold or fivefold rotation about \(n\) forces
\[
 A=\alpha(nn^T-I/3),
\]
with spectrum \((2\alpha/3,-\alpha/3,-\alpha/3)\). Those fixed points have a
double eigenvalue. The twofold fixed locus is
\[
 A=\begin{pmatrix}u&v&0\\v&w&0\\0&0&-u-w\end{pmatrix},
\]
which is three-dimensional, hence two-dimensional after overall scaling,
and admits simple spectrum. It leaves exactly enough continuous freedom to
carry the two independent mass ratios without predicting them. The
source-side family problem is therefore selection by a specific
screen-derived \(A_5\)-invariant potential. One coefficient-free candidate
is excluded: the quartic truncation of the relative-entropy expansion
restricted to \(W_5\) at fixed band amplitude has exactly two extremal
orbit branches, the degenerate vertex-axis prolate orbit and a
simple-spectrum twofold-fixed orbit with centered eigenvalues proportional
to \((-(\varphi^4-1),-1,\varphi^4)\) and sorted-gap ratio exactly the
golden ratio \(\varphi\); the observed centered log-mass gap ratio
\(1.889\) differs from \(\varphi\) by \(16.7\) percent in the closest
orientation, so the entropy packet selects neither branch as the charged
shape. No source mechanism beyond that truncation is supplied.

The finite theorem does not supply an effective action on \(W_5\) with a
unique simple-spectrum minimizing orbit. The invariant algebra has one
quadratic, two cubic, and two quartic coefficients, so the
\(W_5\)-restricted Hessian scalar \(h_5\) at the homogeneous state is a
necessary local discriminator. A physical charged shape additionally
requires the physical \(A_5\) family lift, normed determinant-line descent
with kinetic factors, an interacting charged kernel, an operational scale,
and a source packet fixed before comparison.

The affine scale has a second admissible closure through electromagnetic
transport. With the exact one-loop Ward kernel
\(I(z)=\int_0^1x(1-x)\log(1+zx(1-x))\,\mathrm dx\) in closed form, the
charged response
\(\mathcal W_Q(\mu;\ell)=\frac2\pi\sum_iI(q^2e^{-2(\mu+\ell_i)})\) is a
smooth, strictly decreasing bijection of the common log-scale \(\mu\) onto
\((0,\infty)\) at fixed centered shape \(\ell\), with high-energy slope
\(-2/\pi\). A source-complete spacelike Ward endpoint pair with a
source-complete nonleptonic subtraction therefore fixes a unique
\(\mu_{\mathrm{ch}}\), hence the determinant line
\(|\det(M_e/E_\star)|=e^{3\mu_{\mathrm{ch}}}\) and the electron ratio
\(m_ec^2/E_\star\) consumed by the cesium clock packet. The common-shift
no-go of the preceding subsections holds on the reduct without
mass-dependent electromagnetic transport; a Ward endpoint response and a
determinant-line basepoint are the two admissible closures of that orbit.
The Ward inversion, its interval enclosure, and its fail-closed source-packet
schema are encoded in the Ward determinant-line receipt. A source result
requires the endpoint pair, nonleptonic subtraction, higher-order
monotonicity certificate, and atomic transport factor.

\section{Neutrino Family Derivation}

The neutrino branch contains one exact no-go and one failed comparison candidate. The intrinsic
isotropic branch is excluded by its spectral cap. The weighted-cycle construction descends from
two hand-written family-transport matrices and declared label, orientation,
cycle, and exponent choices selected with target information. Its bridge and
absolute attachment are compare-only.

For the isotropic matrix \(M=aI+\rho C\) with unit-modulus off-diagonal entries, the general
Gershgorin estimate applied to \(M^\dagger M\) gives
\[
\max_{i,j}|\Delta m_{ij}^2|\leq 8a\rho+4\rho^2.
\]
The certified bound is \(1.52304\times10^{-6}\,\mathrm{eV}^2\), about three orders of magnitude
below the atmospheric scale.

The weighted-cycle comparison candidate gives
\[
\begin{aligned}
\theta_{12}&=34.225904631810025^\circ,\\
\theta_{23}&=49.72282845058266^\circ,\\
\theta_{13}&=8.686355527700156^\circ,
\end{aligned}
\]
\[
\begin{aligned}
\delta_{\mathrm{PMNS}}&=305.58061231449796^\circ,\\
J&=-0.02753115613565372,
\end{aligned}
\]
and the dimensionless hierarchy invariant
\[
\begin{aligned}
\frac{\Delta m_{21}^2}{\Delta m_{32}^2}
&=0.030721110097966534.
\end{aligned}
\]
The official NuFIT 6.1 normal-ordering profile gives
\[
\begin{aligned}
\Delta\chi^2_{23,\delta}&=20.11955
\quad\text{with the tabulated atmospheric likelihood},\\
\Delta\chi^2_{23,\delta}&=18.43528
\quad\text{without it},
\end{aligned}
\]
at \((\sin^2\theta_{23},\delta_{\mathrm{CP}})=(0.582056,-54.419^\circ)\)~\cite{nufit61}.
Both values exceed the two-parameter \(3\sigma\) contour value \(11.83\). Separate marginal
intervals do not test this correlation. The published profiles overlap and are not summed;
their maximum is a lower bound on the unavailable full fixed-candidate \(\Delta\chi^2\).

The declared shared-basis construction sets
\[
M_{\mathrm{shared}}=U_{e,\mathrm{left}}^*\,M_{\mathrm{wc}}\,U_{e,\mathrm{left}}^\dagger,
\qquad
U_{\nu,\mathrm{shared}}=U_{e,\mathrm{left}}\,U_{\mathrm{wc}},
\]
the identity
\[
U_{e,\mathrm{left}}^\dagger U_{\nu,\mathrm{shared}}=U_{\mathrm{wc}}.
\]
This cancellation is a tautology because \(U_{\nu,\mathrm{shared}}\) was defined as
\(U_{e,\mathrm{left}}U_{\mathrm{wc}}\). It does not independently identify the physical
charged-lepton basis or validate the PMNS point. No theorem in the combined theorems places the
\(f\)-labelled weighted-cycle operator in the charged-lepton mass basis.

The charged-basis construction is open, and its three normalized shape
singular values differ only at about
the \(10^{-12}\) relative level, so its left singular vectors do not define a stable physical
family basis. The charged-lepton continuation is support-extension gated and does not
supply a replacement labelled basis. At source level, the physical PMNS matrix is therefore
unformed.

Taking the stored basis declarations literally gives a sharply different diagnostic. Reordering
\(M_{\mathrm{wc}}\) from \([f3,f1,f2]\) to the independently declared shared basis
\([f1,f2,f3]\), computing its Takagi matrix \(U_\nu^{(f)}\), and then forming
\(U_{e,\mathrm{left}}^\dagger U_\nu^{(f)}\) yields
\[
\theta_{12}=38.6812^\circ,
\qquad
\theta_{23}=76.2759^\circ,
\qquad
\theta_{13}=44.8283^\circ,
\]
with \(\delta=324.2140^\circ\) and \(J=-0.0233164\). In particular,
\(\sin^2\theta_{13}=0.4970\) and \(\sin^2\theta_{23}=0.9437\), both outside the corresponding
NuFIT 6.1 profile grids. This is a diagnostic conditional on the stored matrices and label
order. It is not a source-closed prediction because the charged basis and family kernel also
descend from template-level inputs.

The transported Takagi identity is
\[
U_{\nu,\mathrm{shared}}^T M_{\mathrm{shared}} U_{\nu,\mathrm{shared}}
=\operatorname{diag}(m_i)\in\mathbb R_{>0}.
\]
Here \(U_{\mathrm{wc}}\) denotes the canonical phase-fixed Takagi unitary, not the raw
arbitrary-phase eigensystem of \(M_{\mathrm{wc}}^\dagger M_{\mathrm{wc}}\).
Equivalently, \(U_{\mathrm{wc}}^T M_{\mathrm{wc}}U_{\mathrm{wc}}\) is positive diagonal in its
declared basis. Canonical Takagi readout under the
declared charged-basis assumption and the electron-row gauge \(U_{e1}\in\mathbb R_{>0}\) gives the
comparison-only Majorana pair
\[
\alpha_{21}^{(\mathrm{Maj})}=153.6185177794357^\circ,
\qquad
\alpha_{31}^{(\mathrm{Maj})}=257.0032408220805^\circ.
\]
The declared exponent law uses the positive transport-load segment between
\(\chi=1+\epsilon\) and \(1+\gamma_{1/2}\): on a one-dimensional affine segment,
the balanced selector and the least-distortion selector for any positive
translation-invariant quadratic form coincide at the midpoint, so
\[
D_\nu=\frac{\chi+(1+\gamma_{1/2})}{2},
\qquad
p=1+\gamma+\frac{\epsilon}{D_\nu}.
\]
The midpoint fact does not derive the segment endpoints, the exponent law, the cycle topology, or
their application to neutrino transport. The formula is target-informed and is not a source-side
consequence of OPH.
On the declared positive selector segment there is a stronger compare-only two-parameter adapter:
under the declared normal-ordering hypothesis, solving \(\tau_\nu\) against the representative
central ratio and then solving \(\lambda_\nu\) against \(\Delta m_{32}^2\) gives
\[
(m_1,m_2,m_3)=(0.01745663295,\ 0.01948419960,\ 0.05308139066)\ \mathrm{eV},
\]
\[
\begin{aligned}
\Delta m_{21}^2&=7.49\times10^{-5}\ \mathrm{eV}^2,\\
\Delta m_{31}^2&=2.5129\times10^{-3}\ \mathrm{eV}^2,\\
\Delta m_{32}^2&=2.438\times10^{-3}\ \mathrm{eV}^2.
\end{aligned}
\]
These exact central numbers are compare-only. The reduced invariant contains target feedback, so
the following values are diagnostic coordinates:
\[
\begin{aligned}
C_\nu&=0.9994295999075177,\\
P_\nu&=6.699825740519345,\\
B_\nu&=P_\nu C_\nu=6.696004159297337,
\end{aligned}
\]
and therefore
\[
\lambda_\nu=\frac{m_{\star,\mathrm{eV}}}{q_{\mathrm{mean}}^{p_\nu}}\,P_\nu C_\nu=1.7237014208357415,
\]
\[
m_i=\lambda_\nu \widehat m_i,
\qquad
\Delta m_{ij}^2=\lambda_\nu^2 \widehat{\Delta}_{ij}.
\]
The normalized same-label overlap-defect calculation is exact after the template and declared
selectors are supplied. The positive-segment adapter, bridge corridor, and bridge-coordinate
calculations are diagnostic. In particular, the following quantity is only a diagnostic
coordinate:
\[
C_\nu:=\frac{B_\nu}{I_\nu^{1/2}\,\widehat{\mathrm{ratio}}^{\,1/2}\,\mathrm{sum\_defect}^{-1}}
\]
It is defined from the declared proxy and normalizer. The construction also factors exactly through
\(q_e = q_{\mathrm{mean}} qbar_e\), so a \(qbar_e\)-only collapse law for the bridge factor does
not follow from the conditional algebraic surface stated here. The best algebraically complete
conditional local object beneath that bridge is the
defect-weighted same-label edge family \(q_e=\sqrt{g_e d_e}\) together with the induced
\(\mu_e\) family. That family sits below \(C_\nu\) and the induced amplitude
\(B_\nu\); it does not replace them.
No source-only PMNS, hierarchy, Majorana-phase, or absolute-mass prediction survives these gates.

\paragraph{Cosmological-neutrino pressure surface.}
The compare-only absolute attachment displays
\[
\sum_i m_{\nu_i}
=
0.09001192964464505~\mathrm{eV}.
\]
Under the declared normal-ordering hypothesis, standard relic-neutrino inheritance, and absent an
explicitly declared extra OPH relativistic coherence channel, the diagnostic
cosmological-neutrino-background branch uses
\[
N_{\mathrm{eff}}^{\mathrm{OPH}}=3.044,
\]
with the usual broad, low-amplitude normal-ordering free-streaming imprint. This is a diagnostic
cosmological pressure surface for the rejected candidate.
The displayed mass sum sits below the Planck+BAO \(0.12~\mathrm{eV}\) bound~\cite{planck18} but is
exposed to strict DESI DR2 \(\Lambda\)CDM-style bounds near \(0.0642~\mathrm{eV}\) under their
model assumptions~\cite{desidr2nu}. A cosmological upper bound below the displayed compare-only value would
exclude the displayed absolute-mass continuation under a model class that also respects
oscillation lower limits and the OPH background assumptions.

On the same-label neutrino-only branch, the centered eta-class is exactly \(S_3\)-isotropic, so
the same-label data are edge-constant and the solar \(1\!-\!2\) split cannot open there by itself.
The first solar mover therefore has to come from realized flavor-side same-label gap/defect
readback, not from another neutrino-only selector tweak.

\subsection{Intrinsic neutrino eta-chain}

The proof-facing input is smaller than the raw same-label matrix payload. It compresses to the
same-label scalar certificate
\[
\bigl(g_e,\ \omega_e\bigr)_{e\in\{12,23,31\}},
\qquad
\omega_e=\mathrm{same\text{-}label\ overlap}_e^2,
\qquad
d_e=1-\omega_e,
\]
modulo one common scale. From that certificate one forms
\[
q_e=\sqrt{g_e d_e},
\qquad
\eta_e=\log q_e-\frac13\sum_f \log q_f,
\qquad
e\in\{12,23,31\},
\]
the intrinsic selector depends only on the centered class
\[
[\eta_e]\in \mathbb{R}^3/\langle(1,1,1)\rangle.
\]
Equivalently one may work with any positive normalized family
\[
\mu_e=\frac{e^{\eta_e}}{\frac13\sum_f e^{\eta_f}}.
\]
Common scaling cancels identically, so the intrinsic selector is determined by the centered
eta-class alone.

This factorization result is conditional on the scalar inputs; it is not a source-closure result.
The scalar values are numerically complete. Their gap data inherit the
template family-transport kernel, and their overlap data inherit a
candidate-only line lift. Neither input is source-closed, so the intrinsic
spectrum is a conditional diagnostic.

On the isotropic neutrino-only branch one has
\[
(\eta_{12},\eta_{23},\eta_{31})=(0,0,0),
\]
which gives the conditional intrinsic singular values
\[
(s_0,s_1,s_2)=
(2.3986448447627196,\ 2.3986448447627196,\ 2.590074050773907)\times10^{-12}\ \mathrm{GeV},
\]
with ascending squared-mass gaps
\[
\begin{aligned}
s_1^2-s_0^2&=0,\\
s_2^2-s_0^2&=9.549864971855843\times10^{-25}\ \mathrm{GeV}^2.
\end{aligned}
\]
This is the exact neutrino-only isotropy obstruction behind the solar-splitting boundary.

\begin{theorem}[Fixed isotropic reference plus same-label scalars suffice for the intrinsic deformation]
Fix the isotropic reference \(M_0\), equivalently its parameters
\((a,\rho,\Omega)\). Assume the same-label gap and overlap scalars are supplied on all realized
arrows. Then the full
intrinsic neutrino mass-eigenstate bundle factors through the scalar certificate
\[
\bigl(g_e,\omega_e\bigr)_{e\in\{12,23,31\}}
\]
or equivalently through the centered eta-class \([\eta_e]\). No additional raw matrix payload is
needed to form the intrinsic selector, the depressed-cubic spectrum, the ascending gaps, or
the intrinsic spectral subspaces. When the singular spectrum is simple, canonical Takagi
congruence fixes an ordered column basis up to column signs. At a degeneracy, including the
isotropic \(s_0=s_1\) point, only the degenerate spectral subspace and its projector are fixed.
Ascending singular-value sorting does not select the physical normal or inverted mass-eigenstate
labels.
\end{theorem}

\begin{theorem}[Exact principal-branch selector from the centered eta-class]
Fix the cycle sum
\[
\Omega=\psi_{12}+\psi_{23}+\psi_{31},
\]
and a positive weight family \(\mu_e\). Define
\[
\mu_{\min}=\min_e\mu_e,
\qquad
F_{\max}=\sum_e\arcsin\!\left(\frac{\mu_{\min}}{\mu_e}\right).
\]
If \(|\Omega|<F_{\max}\), then on the principal branch
\(\psi_e\in(-\pi/2,\pi/2)\), the affine energy
\[
A(\psi)=\sum_e \mu_e(1-\cos\psi_e)
\]
has a unique minimizer. It is given by
\[
\psi_e^\star=\arcsin(\lambda/\mu_e),
\]
where \(\lambda\) is the unique solution of
\[
\sum_e \arcsin(\lambda/\mu_e)=\Omega,
\qquad
\lambda\in(-\min_e\mu_e,\min_e\mu_e).
\]
\end{theorem}

\begin{proof}
The Euler--Lagrange equations are \(\mu_e\sin\psi_e=\lambda\). On the principal branch the scalar
function
\[
F(\lambda)=\sum_e\arcsin(\lambda/\mu_e)
\]
is strictly increasing because
\[
F'(\lambda)=\sum_e\frac{1}{\sqrt{\mu_e^2-\lambda^2}}>0.
\]
Its range on \((-\mu_{\min},\mu_{\min})\) is \((-F_{\max},F_{\max})\), so
\(F(\lambda)=\Omega\) has a unique solution under the stated condition. Strict
convexity follows because the Hessian of \(A\) is
\[
\nabla^2A=\mathrm{diag}(\mu_e\cos\psi_e),
\]
which is positive definite on the principal branch and is positive definite after
restriction to the affine plane \(\sum_e\psi_e=\Omega\).
\end{proof}

For the normalized family
\[
(\mu_{12},\mu_{23},\mu_{31})=(1.4497003,\ 0.9968526,\ 0.5534471),
\]
one has \(F_{\max}=2.5511001\), \(|\Omega|=1.7561443\), and positive domain margin
\(0.7949558\). The numerical selector lies inside the theorem domain.

Let
\[
a=m_\star=\frac{v^2}{\mu_u},
\qquad
\rho = |(M_0)_{12}|.
\]
Then the intrinsic Majorana matrix is
\[
M(\eta)=
\begin{pmatrix}
a & \rho e^{i\psi_{12}} & \rho e^{i\psi_{31}}\\
\rho e^{i\psi_{12}} & a & \rho e^{i\psi_{23}}\\
\rho e^{i\psi_{31}} & \rho e^{i\psi_{23}} & a
\end{pmatrix}.
\]

\begin{theorem}[Exact intrinsic spectral cubic]
Let \(H=M^\dagger M\). Then
\[
H=dI+T,
\qquad
d=a^2+2\rho^2,
\]
where \(T\) has zero diagonal and off-diagonals
\[
x_{12}=2a\rho\cos\psi_{12}+\rho^2e^{i(\psi_{23}-\psi_{31})},
\]
\[
x_{23}=2a\rho\cos\psi_{23}+\rho^2e^{i(\psi_{31}-\psi_{12})},
\]
\[
x_{13}=2a\rho\cos\psi_{31}+\rho^2e^{i(\psi_{23}-\psi_{12})}.
\]
Define
\[
P=|x_{12}|^2+|x_{23}|^2+|x_{13}|^2,
\qquad
Q=\Re\!\bigl(x_{12}x_{23}\overline{x_{13}}\bigr).
\]
Then the eigenvalues \(\lambda_k\) of \(T\) are exactly the three real roots of
\[
\lambda^3-P\lambda-2Q=0,
\]
and the intrinsic squared masses are
\[
m_k^2=d+\lambda_k.
\]
Equivalently,
\[
\lambda_k=
2\sqrt{\frac{P}{3}}
\cos\!\left(
\frac13\arccos\!\left(\frac{3\sqrt3\,Q}{P^{3/2}}\right)-\frac{2\pi k}{3}
\right).
\]
\end{theorem}

\begin{corollary}[Intrinsic eta-chain spectral closure]
Once the centered eta-class is emitted at the flavor boundary, the intrinsic neutrino branch
emits
\[
s_0\leq s_1\leq s_2,\qquad
s_1^2-s_0^2,\ s_2^2-s_0^2,\ s_2^2-s_1^2,
\]
and the corresponding intrinsic spectral subspaces with no PMNS import and no flavor-label
leakage. A columnwise Takagi basis is emitted only on the simple-spectrum locus. The physical normal-ordering assignment is
\((\nu_1,\nu_2,\nu_3)=(s_0,s_1,s_2)\), while the inverted-ordering assignment is
\((\nu_3,\nu_1,\nu_2)=(s_0,s_1,s_2)\). The source does not choose between them, so
physical ordering is open.
\end{corollary}

The matrix used in this corollary is Majorana, so its physical column matrix is the Takagi matrix
\(U_\nu\) defined by
\[
U_\nu^T M U_\nu=\operatorname{diag}(s_0,s_1,s_2)\in\mathbb R_{\geq0}.
\]
The Takagi matrix, whose columns diagonalize \(M^\dagger M\), gives maximum congruence
off-diagonal residual \(5.4\times10^{-28}\,\mathrm{GeV}\) on the anisotropic matrix. On the
stored charged matrix, the intrinsic combination gives
\[
(\theta_{12},\theta_{23},\theta_{13},\delta)
=(45.0137^\circ,\ 2.50489^\circ,\ 2.53597^\circ,\ 210.6251^\circ).
\]
That diagnostic is strongly incompatible with the NuFIT angle surface, but it is not an OPH
prediction because the charged basis is open and template-derived.

\subsection{Perturbative laws around the isotropic point}

Write
\[
\psi_e=\varphi+\delta_e,
\qquad
\delta_{12}+\delta_{23}+\delta_{31}=0.
\]
At first order in the centered eta-class,
\[
\delta_e=-\tan\varphi\,\eta_e+O(\eta^2).
\]
Define
\[
\sigma^2=\frac23\sum_e \delta_e^2.
\]
Assume \(2a\cos\varphi+\rho\neq0\) and that the collective state is separated from the
isotropic doublet. Then the two ascending doublet eigenvalues satisfy
\[
s_{0,1}^2=m_d^2\mp 2a\rho|\sin\varphi|\,\sigma+O(\delta^2),
\]
so
\[
s_1^2-s_0^2 = 4a\rho|\sin\varphi|\,\sigma+O(\delta^2)
=4a\rho\frac{\sin^2\varphi}{|\cos\varphi|}
\sqrt{\frac23\sum_e\eta_e^2}+O(\eta^2).
\]

Under the declared normal-ordering hypothesis, the ascending atmospheric gaps obey
\[
\Delta m_{31}^2 = \Delta_{\mathrm{atm,iso}} + \frac12\Delta m_{21}^2 + O(\delta^2),
\qquad
\Delta m_{32}^2 = \Delta_{\mathrm{atm,iso}} - \frac12\Delta m_{21}^2 + O(\delta^2),
\]
so the first-order invariant atmospheric object is the collective-to-doublet centroid gap
\[
\Delta_{\mathrm{cent}}=
s_2^2-\frac12(s_0^2+s_1^2)
=
\Delta_{\mathrm{atm,iso}}+O(\delta^2).
\]
Its quadratic shift is
\[
\delta\Delta_{\mathrm{cent}}
=
-
a\rho\,\sigma^2
\frac{a(4\cos^2\varphi-1)+6\rho\cos\varphi}
 {2a\cos\varphi+\rho}
+O(\delta^3).
\]

The projective largest-mass right-singular direction, in the phase gauge that
approaches the real democratic vector at the isotropic point, obeys the
first-order deformation law
\[
u_3=
u+
\kappa
\begin{pmatrix}
\delta_{23}\\
\delta_{31}\\
\delta_{12}
\end{pmatrix}
+O(\delta^2),
\qquad
u=\frac1{\sqrt3}(1,1,1)^T,
\qquad
\kappa=
\frac{\sqrt3\,(2a\sin\varphi+3i\rho)}
 {9(2a\cos\varphi+\rho)}.
\]
This equation fixes only the complex line. If \(v_3\) denotes its normalized
right-hand side, the canonical Takagi representative used by the executable
construction is
\[
u_3^{\rm Tak}=e^{-\frac{i}{2}\arg(v_3^T M_\nu v_3)}v_3,
\qquad
(u_3^{\rm Tak})^T M_\nu u_3^{\rm Tak}=s_3>0,
\]
up to the unavoidable column sign. In particular, at the isotropic point the
real democratic vector is generally outside positive Takagi gauge.
Equivalently,
\[
u_3=
u-
\tan\varphi\,\kappa
\begin{pmatrix}
\eta_{23}\\
\eta_{31}\\
\eta_{12}
\end{pmatrix}
+O(\eta^2).
\]

One exact demonstration uses centered eta-class
\[
(\eta_{12},\eta_{23},\eta_{31})
=
\begin{gathered}
(0.13631072512014578,\,-0.18362987264261327,\\
0.0473191475224675),
\end{gathered}
\]
and returns the ascending intrinsic singular values
\[
(s_0,s_1,s_2)=
(2.3929601069646055,\,
2.4048200109774875,\,
2.589606227283229)\times10^{-12}\ \mathrm{GeV},
\]
with
\[
\begin{aligned}
s_1^2-s_0^2
&=5.690121167370743\times10^{-26}\ \mathrm{GeV}^2,\\
s_2^2-s_0^2
&=9.798023388600230\times10^{-25}\ \mathrm{GeV}^2.
\end{aligned}
\]
These are exact outputs of the intrinsic eta-chain once the fixed isotropic reference and centered
eta-class are supplied. They are not flavor-labeled OPH rows, and ascending sorting does not pick
a physical ordering. Under the declared normal-ordering hypothesis
\((\nu_1,\nu_2,\nu_3)=(s_0,s_1,s_2)\); under the inverted-ordering hypothesis
\((\nu_3,\nu_1,\nu_2)=(s_0,s_1,s_2)\). The source-side label rule is absent.

\subsection{Weighted-cycle bridge rigidity and absolute attachment}

The weighted-cycle route supplies a frozen PMNS/hierarchy candidate. The reduced invariant
\(C_\nu\) was selected on a compare-only correction surface, with
\(B_\nu=P_\nu C_\nu\) retained as a diagnostic amplitude parameterization:
\[
\begin{aligned}
C_\nu&=G^2Q S^{-1/2}=0.9994295999075177,\\
P_\nu&=6.699825740519345,\\
B_\nu&=P_\nu C_\nu=6.696004159297337.
\end{aligned}
\]
The blocked absolute attachment displays
\[
\begin{aligned}
\lambda_\nu&=\frac{m_{\star,\mathrm{eV}}}{q_{\mathrm{mean}}^{p_\nu}}\,P_\nu C_\nu=1.7237014208357415,\\
m_i&=\lambda_\nu \hat m_i,\\
\Delta m_{ij}^2&=\lambda_\nu^2\widehat{\Delta m_{ij}^2}.
\end{aligned}
\]
The basis permutation, holonomy orientation, and CP sign lack source-side selection. Exhaustive
enumeration of both stored orientations and all row and mass-column relabelings of the stored
candidate matrix finds no normal-ordering-consistent relabeling that passes both NuFIT 6.1
correlated \(3\sigma\) gates. This excludes a convention-only rescue of that matrix; it neither
derives nor exhausts possible source-side physical charged-basis placements.
The one-parameter atmospheric-anchor slice and the two-parameter positive-segment adapter are
compare-only continuations.

\section{Hadrons, QCD, and the Emergence of Ordinary Matter}

The hadron derivation is the most operationally demanding part of the particle construction. The
source-only artifact is not a fitted hadronic residual and not a bare scalar
\(\rho_{\mathrm{had}}\). It is a source-derived hadronic spectral backend: a QCD quotient
ensemble, source QCD parameter map, Euclidean slab/vacuum-transfer construction, hadronic
Hilbert quotient, Ward-normalized electromagnetic current accounting, positive two-current
spectral export \(d\rho_Q^{(2)}\), higher-point and transition spectral exports,
same-scheme remainder \(\Xi_Q\), and systematics accounting. A source-only hadron mass row requires
a working OPH hadron backend, such as GLORB/Echosahedron, that can emit these
objects with verifiable provenance and production systematics. Small local
surrogates cannot provide the required quantum chromodynamics. The two-current spectral measure is
the marginal needed for running-\(\alpha\) and HVP transport; it is insufficient for HLbL and
rare-decay long-distance amplitudes without the four-current and transition sectors. The
source-backend boundary has an empirical comparison contract. The empirical comparison uses a
separate \(e^+e^-\to\mathrm{hadrons}\) payload class for the hadronic spectral contribution and
stays separate from source-only OPH rows. The source-only derivation nevertheless keeps the
mathematical execution bridge, deterministic sampling law, evaluator, and
surrogate validation as an execution contract for that backend.

\subsection{Seeded \(2+1\) family and unquenched measure}

Let
\[
m_l := \frac{m_u+m_d}{2},
\qquad
\rho_l := \frac{m_u+m_d}{2\,\Lambda_{\overline{\mathrm{MS}}}^{(3)}},
\qquad
\rho_s := \frac{m_s}{\Lambda_{\overline{\mathrm{MS}}}^{(3)}}.
\]
The seeded fixed-physics family is
\[
a\Lambda_n = a\Lambda_{\mathrm{seed}}\,2^{-n},
\qquad
a\Lambda_{\mathrm{seed}} = \sqrt{\rho_l \rho_s},
\]
\[
a m_l^{(n)} = a\Lambda_n \rho_l,
\qquad
a m_s^{(n)} = a\Lambda_n \rho_s,
\]
\[
\beta_n = 6 + \frac{9}{2\pi^2}\log\!\frac{1}{a\Lambda_n},
\qquad
L_n = \left\lceil \frac{\lambda_L^{\mathrm{target}}}{a\Lambda_n}\right\rceil,
\qquad
T_n = \left\lceil \frac{\lambda_T^{\mathrm{target}}}{a\Lambda_n}\right\rceil.
\]
The unquenched ensemble measure is
\[
d\mu_n(U)
=
Z_n^{-1}
\exp[-S_g(U;\beta_n)]
\det D_l(U; a m_l^{(n)})^2
\det D_s(U; a m_s^{(n)})
\,dU.
\]
On this branch, \(N_f=2+1\) and QED is off.

\subsection{Deterministic configuration and source contract}

The emitted execution law is
\[
U_{n,c}
=
K_n^{N_{\mathrm{therm}} + cN_{\mathrm{sep}}}
(U_{\mathrm{cold}}; \operatorname{seed}_{n,c}),
\]
with stop-time formula
\[
t_{\mathrm{stop}}(n,c)
=
N_{\mathrm{therm}} + cN_{\mathrm{sep}}.
\]
The deterministic cfg seed law is
\[
\operatorname{seed}_{n,c}
=
\operatorname{Decode}\bigl(h_{n,c}\bigr),
\]
with
\[
\;h_{n,c}
=
\mathrm{SHA256}\!\Big(
\mathrm{Serialize}\!\big(
n,\beta_n,L_n,T_n,a m_l^{(n)},a m_s^{(n)},c
\big)
\Big).
\]
The source set is fixed as
\[
S_{n,c} =
\left\{
(0,0,0,0),
\left(\left\lfloor\frac{L_n}{2}\right\rfloor,
\left\lfloor\frac{L_n}{2}\right\rfloor,
\left\lfloor\frac{L_n}{2}\right\rfloor,
\left\lfloor\frac{T_n}{2}\right\rfloor\right)
\right\}.
\]
The surrogate execution uses
\[
N_{\mathrm{therm}} = 2048,
\qquad
N_{\mathrm{sep}} = 512.
\]
These values are surrogate execution inputs only; they are not claimed as theorem outputs.

The production geometry summary makes the runtime boundary concrete. On the emitted
seeded \(2+1\) family there are three ensembles and six configurations total. If one stores all
four links at every site as full double-complex \(3\times 3\) matrices, the naive raw gauge
storage estimate is
\[
576\ \mathrm{bytes/site},
\]
which gives total naive raw gauge storage
\[
2.80071464105088\times 10^{14}\ \mathrm{bytes}
\]
over the production schedule. The normalized backend correlator dump needed by the
analysis layer is tiny:
\[
195264\ \mathrm{bytes}
\]
for the full \(\pi_{\mathrm{iso}}\), \(N_{\mathrm{iso,dir}}\), and \(N_{\mathrm{iso,ex}}\)
payload on the production schedule. The hadron requirement is the backend export bundle
that would feed the normalized dump from real production execution.

\subsection{Operational reading of the required computation}

The lightweight analysis layer is separate from the physical backend. The
seeded ensemble family, deterministic configuration and source contract,
jackknife evaluator, forward-window selector, and uncertainty schema are
explicit. A physical result requires Ward-projected hadronic spectral data
from a working OPH hadron backend.

Technically, the required production computation is the standard lattice-QCD stable-channel
workflow on the emitted family: run unquenched rational or standard hybrid
Monte Carlo for the three seeded ensembles; for each realized configuration
and fixed source solve the clover-improved Wilson light/strange systems; construct the
zero-momentum \(\pi_{\mathrm{iso}}\), \(N_{\mathrm{iso,dir}}\), and \(N_{\mathrm{iso,ex}}\)
two-point sequences; write the normalized backend bundle; then apply the
jackknife and forward-window estimators to obtain
\(a m_{X,\mathrm{ground}}\), \(R_X\), \(m_X[\mathrm{GeV}]\), and the published
\(\sigma_{\mathrm{stat}}\), \(\delta_{\mathrm{cont}}\), \(\delta_{\mathrm{vol}}\), and
\(\delta_{\chi}\) fields. This gives a complete source-only execution contract and a separate
empirical display policy.

The engineering burden is asymmetric. Local execution suffices for surrogate
validation and downstream readout tests because those calculations use a tiny
correlator payload. The physical branch needs backend hardware and execution
semantics. This paper treats
GLORB/Echosahedron-class OPH hardware, or an equivalent working OPH hadron backend, as the
source-only backend requirement. Production hadron masses enter the particle construction only with a
backend-emitted Ward-projected spectral measure and its uncertainty budget. This is a
source-backend scope boundary.

\subsection{Stable-channel correlators, effective masses, and forward windows}

The pion stable channel is
\[
p_\pi^{(n,c,s)}(t)
=
\sum_x \Re\,\mathrm{tr}_{c,\mathrm{spin}}
\left[
\gamma_5 S_l(x;s)\gamma_5 S_l(s;x)
\right].
\]
The nucleon stable channel is
\[
p_{N,\mathrm{dir}}^{(n,c,s)}(t) = \sum_x G_d,
\qquad
p_{N,\mathrm{ex}}^{(n,c,s)}(t) = \sum_x G_x,
\]
\[
p_N^{(n,c,s)}(t)
=
p_{N,\mathrm{dir}}^{(n,c,s)}(t) - p_{N,\mathrm{ex}}^{(n,c,s)}(t).
\]
Cfg/source and ensemble averaging are
\[
\bar p_X^{(n,c)}(t)
=
\frac{1}{|S_{n,c}|}\sum_{s\in S_{n,c}} p_X^{(n,c,s)}(t),
\]
\[
C_X^{(n)}(t)
=
\frac{1}{|C_n|}\sum_{c\in C_n} \bar p_X^{(n,c)}(t).
\]

The effective-mass laws are
\[
a m_{\mathrm{eff},\pi}(t)
=
\log\!\frac{C_\pi(t)}{C_\pi(t+1)},
\qquad
a m_{\mathrm{eff},N}(t)
=
\log\!\frac{|C_N(t)|}{|C_N(t+1)|}.
\]
The forward window is
\[
W_n = \{ t : 1 \le t+1 < \lfloor T_n/2 \rfloor \}.
\]
The evaluator monitors log-convexity,
\[
R_{\log\mathrm{conv},\pi}(t)
=
C_\pi(t)^2 - C_\pi(t-1)C_\pi(t+1),
\]
\[
R_{\log\mathrm{conv},N}(t)
=
|C_N(t)|^2 - |C_N(t-1)|\,|C_N(t+1)|,
\]
tail-drop,
\[
D_X(t)
=
a m_{\mathrm{eff},X}(t) - a m_{\mathrm{eff},X}(t+1),
\]
and mirror suppression,
\[
M_X(t)
=
\exp[-a m_{\mathrm{eff},X}(t)(T_n - 2t)].
\]
The selected forward window is the longest contiguous run satisfying finite effective masses,
nonnegative log-convexity up to tolerance, nonnegative tail-drop up to tolerance, mirror
suppression below threshold, and local plateau flatness. The candidate ground-state mass is then
the weighted window average
\[
a m_{X,\mathrm{ground}}
=
\frac{\sum_{t\in W_X^{\mathrm{sel}}} w_t\,a m_{\mathrm{eff},X}(t)}
{\sum_{t\in W_X^{\mathrm{sel}}} w_t},
\qquad
w_t = \frac{1}{\max(\sigma_t^2,\varepsilon)}.
\]

\subsection{Statistics, systematics, and dimensional readout}

Delete-1 jackknife is performed over the cfg axis after source averaging inside each cfg. With
\(n_{\mathrm{cfg}}\) configurations and integrated autocorrelation time
\(\tau_{\mathrm{int},\mathrm{cfg}}\), the effective cfg count is
\[
n_{\mathrm{eff},\mathrm{cfg}}
=
\frac{n_{\mathrm{cfg}}}{2\,\tau_{\mathrm{int},\mathrm{cfg}}}.
\]
The published statistical error is
\[
\sigma_{\mathrm{stat},X}
=
\mathrm{JKstderr}(a m_{X,\mathrm{ground}}).
\]
The machine-readable systematics field uses
\[
\sigma_{\mathrm{sys},X}
=
\sqrt{
\delta_{\mathrm{cont},X}^2
+
\delta_{\mathrm{vol},X}^2
+
\delta_{\chi,X}^2
}.
\]
Continuum, volume, and chiral proxies are encoded as
\[
R_X^{(n)} = \frac{a m_X^{(n)}}{a\Lambda_n},
\qquad
R_X^{(n)} \approx R_X(0) + c_X (a\Lambda_n)^2,
\]
\[
\delta_{\mathrm{cont},X}^{(n)}
=
a\Lambda_n\,\left|R_X^{(n)} - R_X(0)\right|,
\]
\[
\delta_{\mathrm{vol},\pi}^{(n)}
=
a m_\pi^{(n)}
\frac{e^{-a m_\pi^{(n)}L_n}}{\max(a m_\pi^{(n)}L_n,1)},
\]
\[
\delta_{\mathrm{vol},N}^{(n)}
=
a m_N^{(n)}
\frac{e^{-a m_\pi^{(n)}L_n}}{\max(a m_\pi^{(n)}L_n,1)},
\]
\[
\delta_{\chi,\pi}^{(n)}
=
a\Lambda_n \left|R_\pi^{(n)} - \langle R_\pi\rangle_n\right|,
\]
\[
Q_N^{(n)} = \frac{R_N^{(n)}}{R_\pi^{(n)}},
\qquad
\delta_{\chi,N}^{(n)}
=
a\Lambda_n \langle R_\pi\rangle_n
\left|Q_N^{(n)} - \langle Q_N\rangle_n\right|.
\]
These are surrogate publication proxies, not production physical systematics. Given a
ground-state candidate,
\[
R_X = \frac{a m_{X,\mathrm{ground}}}{a\Lambda_{\overline{\mathrm{MS}}}^{(3)}},
\qquad
m_X[\mathrm{GeV}] = R_X\,\Lambda_{\overline{\mathrm{MS}}}^{(3)}[\mathrm{GeV}],
\]
with
\[
\Lambda_{\overline{\mathrm{MS}}}^{(3)} = 0.3344017073\ \mathrm{GeV}.
\]

\subsection{Surrogate execution bundle and frontier}

The surrogate execution evolves a latent state \(z\in\mathbb{R}^d\) with Hamiltonian
\[
H(z,p) = S(z) + \frac12 \sum_i p_i^2,
\]
where
\[
S(z)
=
\frac12 \sum_i \omega_i^2 z_i^2
+
\lambda_4 \sum_i z_i^4
+
\kappa \sum_i (z_{i+1}-z_i)^2
+
2\alpha_l \sum_i \log(\mu_l^2 + z_i^2)
+
\alpha_s \sum_i \log(\mu_s^2 + \tfrac12 z_i^2).
\]
Leapfrog integration plus Metropolis accept/reject gives the surrogate HMC update
\[
(z,p) \mapsto (z',p'),
\qquad
P_{\mathrm{acc}} = \min(1,e^{-\Delta H}).
\]
This kernel is a deterministic executable surrogate honoring the emitted receipt and seed law,
with no physical lattice-QCD RHMC/HMC claim.

For validation of the execution bridge, the surrogate locks the ground-state masses to the hadron
comparison proxy values
\[
\begin{aligned}
m_{\pi,\mathrm{proxy}}
&=0.13497682776768472\ \mathrm{GeV},\\
m_{N,\mathrm{iso,proxy}}
&=
\frac{m_p+m_n}{2}
=
0.93891875434\ \mathrm{GeV}.
\end{aligned}
\]
Thus
\[
a m_{\pi,\mathrm{proxy}}^{(n)}
=
\frac{m_{\pi,\mathrm{proxy}}}{\Lambda_{\overline{\mathrm{MS}}}^{(3)}[\mathrm{GeV}]}
\,a\Lambda_n,
\]
\[
a m_{N,\mathrm{iso,proxy}}^{(n)}
=
\frac{m_{N,\mathrm{iso,proxy}}}{\Lambda_{\overline{\mathrm{MS}}}^{(3)}[\mathrm{GeV}]}
\,a\Lambda_n.
\]
The surrogate correlator is
\[
C_X^{\mathrm{sur}}(t)
=
A_{X,0} e^{-a m_X t}
+
A_{X,1} e^{-a m_{X,\mathrm{ex}} t}
+
A_{X,\mathrm{mir}} e^{-a m_X (T_n-t)}
\]
up to a multiplicative correlated noise factor. Direct and exchange nucleon pieces are written as
\[
C_{N,\mathrm{dir}}^{\mathrm{sur}}(t)
=
f_{\mathrm{dir}}\,C_N^{\mathrm{sur}}(t),
\qquad
C_{N,\mathrm{ex}}^{\mathrm{sur}}(t)
=
(f_{\mathrm{dir}}-1)\,C_N^{\mathrm{sur}}(t),
\]
so that
\[
C_N^{\mathrm{sur}}(t)
=
C_{N,\mathrm{dir}}^{\mathrm{sur}}(t)
-
C_{N,\mathrm{ex}}^{\mathrm{sur}}(t).
\]

On the finest surrogate ensemble, the resulting diagnostic candidates are
\[
m_{\pi,\mathrm{iso}}^{\mathrm{sur}} = 0.135039383836\ \mathrm{GeV},
\qquad
m_{N,\mathrm{iso}}^{\mathrm{sur}} = 0.938960210578\ \mathrm{GeV},
\]
with worst absolute error approximately \(6.26\times10^{-5}\ \mathrm{GeV}\) across the surrogate
stable-channel family. These numbers validate the emitted execution bridge. They are not
promotable production hadron predictions.

So the hadron derivation closes the full
\[
\text{receipt} \to \text{execution} \to \text{writeback} \to \text{evaluation} \to \text{budgets} \to \text{forward-window summary}
\]
path on executed surrogate data, while the physical closure requires production
unquenched RHMC/HMC, real Dirac solves and baryon contractions, production autocorrelation
studies, and production continuum / finite-volume / chiral systematics.

\section{Observer-Centric Particle Ontology and Measurement}

OPH is observer-centric at the level of its basic ontology, so particle
physics also requires an account of what a particle \emph{is} and how a
measurement turns an excitation into an observed state. The fixed-cutoff
record and measurement theorems give that account precise mathematical
content.

\subsection{Observer patches, overlap consistency, and particle data}

Ref.~\cite{ophconsensus} formulates the basic kinematic picture in patch-net language: each observer
patch carries a local state, neighboring patches compare a shared interface alphabet, and a
global state is physically admissible exactly when neighboring projections agree on the overlap.
The algebraic OPH form of the same statement is the first and second axioms:
local physical data are carried by patch algebras \(\mathcal A(P)\), and those local
states must agree on shared subalgebras whenever patches overlap.

The particle interpretation starts from patch-local algebras, overlap-visible observables, edge
sectors, and transport data rather than a single absolute global particle basis. A particle row in
the reported derivation is a readout claim about a stable or candidate-stable excitation structure
visible to a family of observer patches and consistent on their overlaps.

For this record-algebra discussion, use the support-local algebra-state-record reduct
\[
O_{\mathrm{red}}=(P,\mathcal A(P),\rho,R),
\]
where \(P\) is the support patch, \(\mathcal A(P)\) its local algebra, \(\rho\) the local
state, and \(R\) the record algebra. The full operational observer additionally carries overlap
interface algebras and restriction maps, allowed update and repair instruments, and checkpoint data
used for continuation. For the particle derivation, \(R\) belongs to the quantum
observer surface itself. On the exact fixed-cutoff measurement
surface it is generated by central record projectors, and practical readout may use
approximately commuting projectors that are close to that central reference algebra. That is what
makes the observer-facing record surface shareable across overlapping observer descriptions
without violating the usual no-cloning constraints on generic quantum states.

\subsection{Record algebras and definite outcomes}

The integrated measurement appendices and Ref.~\cite{ophmicrophysics} make the measurement
interface precise at fixed cutoff. On the declared operational surface, the completed
write/verify slice carries a finite commutative central record algebra. Practical readout may
instead use projectors \(Q_a\) on the same declared slots with commuting central reference
projectors \(\widehat Q_a\), where
\[
\delta_{\mathrm{rec}}:=\max_a \|Q_a-\widehat Q_a\|.
\]
Then
\[
\|[Q_a,Q_b]\|\le 4\,\delta_{\mathrm{rec}},
\]
and, whenever \(\|\widetilde\rho-\rho\|_1\le\varepsilon\),
\[
\Bigl|
\operatorname{Tr}(\widetilde\rho Q_a)-\operatorname{Tr}(\rho \widehat Q_a)
\Bigr|
\le
\varepsilon+\delta_{\mathrm{rec}}.
\]
On the explicit fixed-cutoff screen architecture, the exact record algebra after a completed
write/verify cycle is
\[
\mathcal Z_{\mathrm{rec}}(t)
=
\mathrm{Alg}\Bigl(
\{P_{m_I^{(J)}(t)}\}_{I\neq J},
\{P_{r_I^{\mathrm{bulk}}(t)}\}_I,
\{\Pi_\alpha^{(IJ)}(t)\}_{I<J,\alpha}
\Bigr).
\]
Ref.~\cite{ophmicrophysics} proves that, on the declared operational measurement surface,
\(\mathcal Z_{\mathrm{rec}}(t)\) is commutative and central for the readout instrument being
used.

This fixed-cutoff centrality result is what supports definite outcomes without adding a second
ontology. The integrated measurement construction phrases the same idea through edge-center
decomposition: on a fixed collar with exact Markov structure aligned with the edge split (the
Standard Model gauge paper's Markov-split alignment hypothesis), the state splits into
superselection blocks,
\[
\rho_{ABC}
=
\bigoplus_j
q_j\,
\rho^{(j)}_{A b_L^{(j)}}\otimes \rho^{(j)}_{b_R^{(j)} C},
\]
and the label \(j\) is classical center data. The supplement then goes one step further: in the
physical algebra there are no interference observables between different sector blocks. So once a
particle detection event has been recorded in the observer-accessible center/record algebra, the
accessible physics is organized as a classical mixture over those blocks instead of as a
superposed measurement record.

The finite matrix theorem states the exact boundary of that claim. For a
supplied projective partition \(\{P_j\}\), let
\(\mathcal E(X)=\sum_jP_jXP_j\). Two matrices \(X\) and \(Y\) give equal trace
statistics against every observable commuting with all \(P_j\) if and only
if \(\mathcal E(X)=\mathcal E(Y)\). In particular,
\(\mathcal E(P_iXP_j)=0\) for \(i\neq j\), so each cross-sector corner is
invisible to the complete sector-preserving commutant. The result is relative
to the supplied partition and readout. It does not construct the physical
edge center used by a detector.

Accordingly, a measured particle family or excitation channel should be
thought of as an observer-accessible sector or record value carried by the central measurement
surface, not as a mysterious absolute collapse of the full universe-state into a preferred basis
chosen by hand.

\subsection{Born probabilities and post-measurement states}

The fixed-cutoff measurement package is also explicit about probabilities and updates. For every
event \(E\) in the sigma algebra generated by \(\mathcal Z_{\mathrm{rec}}(t)\), the measurement
probability is
\[
\mathbb P_t(E)=\operatorname{Tr}\!\bigl(\rho_t P_E\bigr),
\]
where \(P_E\in\mathcal Z_{\mathrm{rec}}(t)\) is the projector for that event. Conditioning on the
event then produces the post-measurement state
\[
\rho_t\!\mid_E
=
\frac{P_E\rho_t P_E}{\operatorname{Tr}(\rho_t P_E)}.
\]
In the supplement, the same statement appears as the L\"uders update on the record algebra. In
the flagship paper \emph{From Observer Consensus to Standard Physics}~\cite{ophsynthesis}, this fixed-cutoff
Born/L\"uders package and its Bell/CHSH extension are imported as theorem-bearing, not
non-theorem commentary.

On the declared binary-icosahedral spinor branch, an exact finite candidate
has \(\lvert S_{\mathrm{CHSH}}\rvert=1+3/\sqrt5\) for an
incidence-defined setting family. The carrier admits (960) maximizing
quadruples, while the declared family contains (120) of them. No source
mechanism selects those settings or supplies a completed two-wing record
instrument, so the candidate is not a physical Bell prediction.

This directly answers the particle-state question in the present framework. Measurements give particles
actual states by conditioning the observer-accessible state on the central record algebra. A
detector click, a stable overlap-sector readout, or a pointer value does more than announce a
pre-existing classical label; it defines the conditioned state for subsequent observer-accessible
physics. If the same event is re-read without any accepted repair move touching its support, the
result in Ref.~\cite{ophmicrophysics} shows that the re-read probability is \(1\). So the measurement interface is
well defined and operationally auditable.

\subsection{Particle identity as transport-stable structure}

Once one asks ``which event happened?'' and then ``which particle family is this?'', the
flavor derivation becomes essential. The active flavor derivation does not start with the
names electron, muon, tau lepton, or up quark built into the fundamental observable. It starts
with transport kernels, generation-bundle data, same-label eigenline transport, overlap-edge
cocycles, and a persistent flavor observable carrying intrinsic labels \(f1,f2,f3\). This is an
important discipline condition in the paper. At the deepest flavor surface used here, family
identity is transport-stable intrinsic structure first and named experimental family assignment
second.

The shared Yukawa/excitation dictionary carries substantial weight even though it is not
itself the leading missing object. The flavor construction gives an invariant base:
projectors, spectral gaps, pair suppressions, cycle phases, and common sector-response objects.
No map from that intrinsic base to the named low-energy fermion families is
derived. The quark calculation supports direct numerical comparison, while
the charged-lepton theorem emits no physical mass. The neutrino weighted-cycle
construction is a target-informed template candidate without a physical
flavor-labelled attachment.

Accordingly, particle identity is a refinement-stable
observer-accessible transport pattern whose named interpretation is inherited from the
available dictionary and declared comparison domain. Without that dictionary,
no named physical identity follows.

\section{Affine Event Records and Cross-Boundary Token Stitching}\label{sec:h3-record-worldlines}

The sector-identity analysis concerns which transport-stable excitation
type is being read. A different question is whether two located record tokens, seen near a chart
or partition boundary, are the same continuing observer-visible token. This section closes that
second question at the finite-certificate level. It does not add a particle species, mass, charge,
or scattering theorem. It says when a bounded patch system is allowed to stitch affine event
records into nonbranching record-worldlines, and when it must report ambiguity.

\subsection{Event chart, frame, and clock atlas}

The compact cap-normal theorem supplies the Lorentz-natural frame hyperboloid. It identifies sky directions with
\(q(\Omega)=(1,\Omega)\), represents oriented round caps by
\(n_C=(\cot\alpha,\csc\alpha\,\mathbf c)\), and maps each cap to a geodesic half-space in
\(H^3\). The compact event-manifold packet separately supplies affine event charts under its
population, separation, chart, cone, and reachability receipts. This section consumes event
supports, observer clocks, chart transitions, gauge transport, and worldline stitching. Frame
estimation, event location, overlap descent, and worldline stitching are separate certificates.

Fix a curvature radius \(R_H>0\) and write
\[
H^3_{R_H}=\{X\in\mathbb R^{1,3}:\langle X,X\rangle_L=-R_H^2,\ X^0>0\},
\qquad
\langle X,Y\rangle_L=-X^0Y^0+X^1Y^1+X^2Y^2+X^3Y^3.
\]
The hyperbolic distance used for conditioned frame comparison is
\[
d_H(X,Y)=R_H\,\mathrm{arcosh}\!\left(-\frac{\langle X,Y\rangle_L}{R_H^2}\right).
\]
An event-record atlas is a tuple
\[
\mathcal A_H=
\left(
R_H,\{U_i,\chi_i,u_i,J_i,\vartheta_i\},\{G_{ji},F_{ji}\},E,\nabla,
\mathsf{Prov}_{R_H},\mathsf{Prov}_u
\right).
\]
Here \(\chi_i\) are affine event charts, \(u_i\in H^3_{R_H}\) is the selected local observer-frame
anchor when such an anchor is claimed, \(J_i\) are event-chart domains for record support, and
\(\vartheta_i:J_i\to T\) are observer-clock maps into a common comparison-time line \(T\).
\(\mathsf{Prov}_{R_H}\) records whether \(R_H\) is a unit convention, an imported physical scale, or
an independently derived OPH scale. The unit-chart convention alone supplies
no independent OPH scale.
\(\mathsf{Prov}_u\) records which observer, tetrad, or clock object selected the anchors \(u_i\).
The Lorentz part \(G_{ji}\in\mathrm{SO}^{+}(1,3)\) transports frames and tetrads; the full
Poincar\'e transition transports affine event coordinates. The map
\(F_{ji}\) is the sector or gauge transport functor on the record payload. These two holonomies
are deliberately separate: atlas holonomy tests geometry, and gauge holonomy tests transported
sector data. When two charts describe the same transported observer frame, the atlas additionally
requires
\[
G_{ji}u_i=u_j.
\]

\begin{theorem}[H3 geometric naturality]\label{thm:h3-geometric-naturality}
For every \(A\in\mathrm{SO}^{+}(1,3)\), the map \(X\mapsto AX\) preserves \(H^3_{R_H}\),
\(d_H\), geodesics, exponential and logarithm maps, parallel transport, Hausdorff distances between
compact frame supports, and Fr\'echet means whenever the mean is unique.
\end{theorem}

\begin{proof}
The defining equation of \(H^3_{R_H}\) and the formula for \(d_H\) depend only on the Lorentz
inner product and the future sheet. Elements of \(\mathrm{SO}^{+}(1,3)\) preserve both. The
Levi-Civita connection, geodesic exponential, logarithm, and parallel transport are functorial for
isometries, so the listed constructions commute with \(A\). Hausdorff distance and unique
Fr\'echet means are metric constructions, hence are preserved by any isometry.
\end{proof}

\subsection{Cap-response frame balls and affine event supports}

The Standard Model gauge paper supplies the conditional inverse theorem that turns calibrated modular cap
responses into conditioned observer-frame supports. For every descended record token \(i\) and
clock slice \(t\), the frame layer consumes a certificate
\[
R_i(C,t,O)=\omega_{i,O}\!\left(\sigma_t^{C,O}(M_{C,0,O})\right),
\qquad
y_{i,j}=F_j(X_i(t))+e_{i,j},
\]
on a compact frame domain \(\Omega\subset H^3_{R_H}\), with quantitative observability
\[
\alpha\bar d(X,Y)\le |F(X)-F(Y)|_W\le L\bar d(X,Y),
\qquad |e_i|_W\le\sigma_i(t).
\]
Given an \(\varepsilon\)-net and residual tolerance \(\tau_i(t)\), the imported theorem emits
\[
S_i(t)=B_H(\widehat X_i(t),r_i(t)),
\qquad
r_i(t)=R_H\left[
\frac L\alpha\varepsilon+\frac2\alpha\sigma_i(t)+\frac1\alpha\tau_i(t)
\right].
\]
A frame estimate \(\widehat X_i(t)\) is a unique finite output only when
\(\Delta_{\mathrm{frame},i}(t)>0\). If this gap fails, the output is a frame ambiguity set or
support ball. This ball cannot serve as an event position. Event stitching consumes an affine
event support \(A_i(t)\) produced by the event-manifold chart and ancestry receipts. Distinct event
supports at a shared clock slice require a certified separation in the affine event metric. A
conditioned-frame separation may also be recorded:
\[
d_H(\widehat X_i,\widehat X_k)>r_i+r_k+m_{\mathrm{sep}}
\]
for a declared positive frame margin \(m_{\mathrm{sep}}\). It does not replace event separation.

\subsection{Proto-record extraction and descent}

A raw local proto-record is a connected component of an invariant detector scalar above a frozen
threshold, together with its certified affine event support, conditioned frame ball, local clock
interval, sector payload, and uncertainty collar.
The detector scalar is built from record-algebra data, not from implementation object IDs. Extraction is
chart-natural when transporting the raw field by \(G_{ji}\) and \(F_{ji}\) carries extracted
components in chart \(i\) to the extracted components in chart \(j\), up to the declared support
collar.

Before temporal stitching, overlapping charts are descended. The overlap-descent relation joins
two proto-records only when their transported affine event supports lie in the same connected
component on a real chart overlap and the same-component margin dominates the support collar. If distinct
component separation does not dominate the collar, the descent emits an ambiguity certificate
instead of choosing a label. On component-connected covers, this descent produces canonical
global descended records by the usual sheaf gluing argument: cocycle-compatible local components
with positive separation margins have a unique global component, and the construction is natural
under atlas restriction.

\subsection{Boundary-crossing germs and transported residuals}

A cross-boundary candidate edge is admissible only when the two descended records meet a real
patch interface in adjacent observer-clock intervals and share a common comparison chart. The
interface contact must be oriented and transverse: the signed-distance interval brackets the
interface, and the normal velocity has a positive lower bound. Tangencies, grazing contacts, and
file or shard-name boundaries without a physical interface are rejected. If the local slab contains
an interaction marker, the record is routed to the interaction solver instead of stitched by the
free-propagation rule.

For an admissible pair, the certificate transports the left record into the right comparison chart
and evaluates frozen residuals
\[
r_t,\quad r_{\partial},\quad r_x,\quad r_v,\quad r_{\mathrm{int}},
\]
for clock adjacency, boundary contact, affine event position, transported velocity or finite
difference, and interaction absence. The cost function is declared before seeing the proposed
matching. Record IDs, global IDs, shard-local IDs, and stitch keys are forbidden in admissibility,
costs, and tie-breaking.
All affine event-position and velocity residuals are evaluated as certified interval or
uncertainty-box distances. A point-estimate residual may be printed as a diagnostic; it cannot
replace the event-support radius in an admissibility or matching certificate. Conditioned
\(H^3\) frame residuals are a separate tetrad-consistency check.

\begin{lemma}[Candidate-edge naturality]\label{lem:h3-candidate-edge-naturality}
The admissible candidate-edge set and its residual intervals are invariant under common
\(\mathrm{SO}^{+}(1,3)\) chart changes and under gauge-bundle trivialization changes.
\end{lemma}

\begin{proof}
Every event term is computed from the affine overlap cocycle, oriented interface data, and the
common clock line. Frame terms use \(d_H\) and
Theorem~\ref{thm:h3-geometric-naturality}. Sector and gauge residuals compare transported payloads after applying the declared
connector, so changing a local trivialization conjugates both sides of the comparison. The
residual intervals and the resulting admissibility predicate are therefore unchanged.
\end{proof}

\subsection{Assignment gap, nonbranching paths, and ambiguity}

Let \(L\) and \(R\) be the descended record sets on two adjacent time slabs. A stitch matching is a
partial one-to-one assignment between \(L\) and \(R\), augmented by declared appearance and
disappearance interval costs. Its certified gap is
\[
\Delta_{\mathrm{cert}}
=\min_{M\ne M_*}\underline C(M)-\overline C(M_*),
\]
where \(M_*\) is the proposed matching, \(\overline C(M_*)\) is its upper cost bound, and
\(\underline C(M)\) is the lower cost bound for every competing matching.
The event-location gap \(\Delta_{\mathrm{event}}\) and stitch-assignment gap
\(\Delta_{\mathrm{cert}}\) are distinct receipts: the first certifies a unique event support
inside one clock slice, while the second certifies a unique temporal assignment across slices.

\begin{theorem}[Stable ID-independent stitch assignment]\label{thm:h3-stable-stitch-assignment}
If the candidate-edge set is natural, the costs are frozen and ID-independent, the proposed
matching is one-to-one, and \(\Delta_{\mathrm{cert}}>0\), then \(M_*\) is the unique certified
minimum-cost stitch. The certified matching is natural under chart changes, gauge
trivializations, and relabeling of implementation record IDs.
\end{theorem}

\begin{proof}
The strict gap says every competing matching has lower cost bound above the proposed matching's
upper cost bound, so no competitor can tie or beat it under any realization inside the certified
intervals. Lemma~\ref{lem:h3-candidate-edge-naturality} makes the candidate graph and cost
intervals invariant under chart and gauge presentation changes. Because IDs are excluded from
admissibility and costs, relabeling them changes only bookkeeping fields, not the optimum.
\end{proof}

\begin{theorem}[Unavoidable ambiguity]\label{thm:h3-unavoidable-ambiguity}
If a nontrivial permutation of the visible descended records preserves the atlas, clocks,
interfaces, sector/gauge payloads, candidate graph, and cost intervals, then no natural
ID-independent finite procedure can certify a unique stitch across that slab.
\end{theorem}

\begin{proof}
A natural procedure must commute with the permutation. A unique certified output would therefore
be fixed by the permutation. But the permutation sends one equally admissible stitch to a distinct
equally admissible stitch. Choosing between them would introduce a label or presentation
dependence. The only natural finite output is an ambiguity label.
\end{proof}

The certified stitch graph is then the graph obtained by chaining the unique matchings across
adjacent slabs. Since each matching is partial one-to-one, every connected component is a
nonbranching path, ray, or isolated record. These are the affine event-record worldlines of this
section.

\subsection{Refinement naturality and dynamical lift}

For a refinement \(s\to r\), let \(Q_{sr}\) contract fine descended records to coarse records and
let \(\eta_{sr}\) bound the cost distortion between the fine and coarse stitch problems. If the
coarse certified gap obeys \(\Delta_r>2\eta_{sr}\), then the contracted fine optimum is
isomorphic to the coarse optimum. This is the refinement-naturality gate: a cross-boundary stitch
is not paper-grade until the coarse/fine contraction agrees within the declared margin.

The stitch certificate lives on the conditional Lorentzian event manifold. A geodesic or
free-particle interpretation requires a separate dynamical receipt with a timelike speed margin,
connection compatibility, and an action or propagation law. The stitch theorem proves token
continuation in event charts; it does not prove geodesic motion.

\begin{theorem}[Certified affine event-record stitching]\label{thm:h3-record-worldline-stitching}
Given a valid affine event atlas satisfying The Standard Model gauge paper's event-manifold receipts, a common
operational clock atlas, affine event supports for every descended record token and clock slice,
and separately typed conditioned-frame certificates where observer frames are used,
chart-natural proto-record extraction, overlap descent with positive event-support separation
margins, real transverse interface-crossing germs, sector and gauge transport continuity, a frozen
ID-independent one-to-one assignment with \(\Delta_{\mathrm{cert}}>0\), no interaction marker in
the free-propagation slab, and a passing coarse/fine contraction gate, the emitted stitch graph is
a unique natural nonbranching affine event-record graph. If any positive event-location gap,
separation margin, or assignment gap is absent, the natural output is \(\mathrm{AMBIGUOUS}\); if an
interaction marker is present, the output is \(\mathrm{INTERACTION\ REQUIRED}\); if atlas,
interface, event-location, or transport gates fail, the stitch is rejected.
\end{theorem}

\begin{proof}
The atlas, event-location, and extraction gates make local proto-records presentation-independent.
The positive \(\Delta_{\mathrm{event}}\) gate licenses event-position outputs; otherwise the
procedure carries event boxes or ambiguity sets forward. Descent turns overlap-compatible local components into
descended records before any temporal decision is made.
The interface and transport gates define a natural candidate-edge graph by
Lemma~\ref{lem:h3-candidate-edge-naturality}. The strict assignment gap gives a unique
ID-independent matching by Theorem~\ref{thm:h3-stable-stitch-assignment}. Chaining partial
one-to-one matchings gives nonbranching components, and the refinement gate makes the result
stable under the declared coarse/fine contraction. The alternative outcomes
are exactly the negations of these conditions, with
Theorem~\ref{thm:h3-unavoidable-ambiguity} forcing ambiguity when
the visible data do not separate a unique continuation.
\end{proof}

\subsection{Scope boundary of the measurement chapter}

The fixed-cutoff consensus theorem gives the central record algebra. Given the
separately declared algebra-state and two-wing representation, the measurement
theorem gives Born probabilities, L\"uders conditioning, repeated-read
stability, and the stated Bell/CHSH bound. It does not prove global observer continuation or
strange-loop closure. Those metaphysical proposals do not enter the technical
argument.

% \section{Results at a Glance}
\section{Discussion: Matter, Antimatter, Supersymmetry, and Physical Boundaries}

The OPH theorems distinguish structural results, conditional quantitative
branches, and target-conditioned comparisons. Matter versus antimatter,
supersymmetry, and several adjacent questions lie on that boundary. The quark
calculation is an obstruction and target-conditioned surface rather than a set
of source-only numeric predictions. Its \(S_3\) table is a target-anchored
diagnostic.

\subsection{Matter and antimatter}
\label{subsec:baryogenesis-source}

The Standard Model structural branch fixes the chiral gauge architecture within which
matter and antimatter are defined. Once a chiral fermion representation is present, its conjugate
representation gives the corresponding antiparticle content. In that limited sense, OPH explains
why matter and antimatter both exist: they are paired by the realized gauge and chiral structure
of the low-energy branch.

The cosmological asymmetry requires a separate finite source object. At regulator \(r\), write
\[
\mathfrak B_r=(Q_r,C_r,E_r,\omega_{R,r},L_{r,T},p_{r,i},
\tau_r,T_r,\rho_{R,r}).
\]
Here \(Q_r\) is the physical quotient, \(C_r\) is its CP involution,
\(\omega_{R,r}\) is a CP-odd integral winding cocycle, \(L_{r,T}\) is the
physical quotient repair generator, \(p_{r,i}\) is the initial law,
\(\tau_r\) and \(T_r(\tau)\) are the physical clock and temperature map, and
\(\rho_{R,r}\) assigns record charges \(r_\psi\) to chiral matter.

For a record phase acting as
\(\psi\mapsto e^{ir_\psi\Theta_R}\psi\), the finite-quotient
anomaly-and-current theorem gives
\[
\boxed{
k_R=\sum_{\psi\,{\rm LH}}r_\psi\,2T_2(R_\psi)
}
\]
and
\[
\boxed{
\dot\Theta_R(T)=\theta_0
\sum_{q,q'\in Q_r}p_r(q,T)L_{r,T}(q,q')
\omega_{R,r}(q,q')
},
\qquad \theta_0=\frac{2\pi}{m_R}.
\]
The first equation is the mixed electroweak anomaly index. The second is an
oriented probability current in proper time. Quotient settlement supplies no
rate matrix, and a repair iteration count supplies no physical clock.

The theorem also fixes the sign boundary. If the initial law and generator are
CP symmetric,
\[
p_{r,i}(C_rq)=p_{r,i}(q),
\qquad
L_{r,T}(C_rq,C_rq')=L_{r,T}(q,q'),
\]
then every transition cancels its CP image and
\[
\boxed{\dot\Theta_R(T)=Y_B=0.}
\]
An oriented register and a normal-form map therefore select no cosmological
matter sign. A nonzero sign requires a quotient-intrinsic CP-odd boundary
condition, action term, or transition affinity.

The conditional \(\mathbb Z_6\) determinant/deck winding supplies an integral
phase coordinate with inversion under orientation reversal. Its natural
Standard Model attachment is the central hypercharge direction. The associated
coefficient vanishes generation by generation:
\[
\boxed{
k_R^{YWW}=N_g(3Y_Q+Y_L)
=3\left(3\cdot\frac16-\frac12\right)=0.
}
\]
This is the mixed \(SU(2)_L^2U(1)_Y\) anomaly-cancellation condition. The
gauge/deck phase cannot directly supply the required
\(\Theta_R W\widetilde W\) coupling. Quotient periodicity and the twelve
electroweak zero modes do not select \(k_R=1\) or any other nonzero record
attachment.

A distinct global record attachment can carry a nonzero index. On the
three-generation branch,
\[
k_R(Y)=k_R(B-L)=0,
\quad k_R(B)=k_R(L)=3,
\quad k_R(B+L)=6.
\]
A quotient-visible gauge-singlet \(B+L\) record phase is therefore a viable
conditional attachment. Its derivation from OPH repair data is open. On
that branch, the freeze-out transport functional requires
\[
\left\langle\frac{\dot\Theta_R}{T}\right\rangle_{\rm fo}
=(4.463\pm0.028)\times10^{-9}.
\]
This is a source-generator target. It may not be consumed to select
\(L_{r,T}\), the attachment, or the CP-odd affinity.

The exact source theorem and gauge/deck no-go are closed at finite quotient.
The physical baryon abundance is work in progress. Its open source
packet consists of a distinct anomalous record phase, source-only transition
rates, a proper-time and temperature calibration, a CP-odd initial or boundary
law, sign-domain coherence, and the resulting full source history. Transport,
washout, and freeze-out must then be evaluated on that emitted history.

\subsection{Why the declared branch does not require supersymmetry}

The heat-kernel appendix below contains the gauge-coupling calibration. The
declared branch needs no supersymmetric partner spectrum to produce its
unification-like running benchmark. This does not rule out supersymmetry in a
different physical completion.

At one loop, the appendix records the familiar fact that Standard Model beta-function
coefficients do not produce successful naive unification, whereas minimal
supersymmetric Standard Model coefficients do. The
conditional calibration hypothesis is that edge-sector heat-kernel weights shift the effective
beta-function coefficients in that direction through geometric or entropic sector
multiplicity instead of through an actual low-energy superpartner spectrum. In that reading,
``unification-like behavior'' and ``a supersymmetric particle zoo'' come apart.

This calibration is not a derived unification theorem. The appendix records one
calibration branch in which Peter--Weyl multiplicities, the printed one-loop running frame, the
threshold conventions, and an additional fermionic-grading restriction together produce an
supersymmetric benchmark shift. It does not prove that those effective loop multiplicities,
statistics restrictions, and decoupling conventions are uniquely forced by OPH edge sectors alone.

For this paper, the correct conclusion is therefore scope-limited and modest:
\begin{enumerate}[leftmargin=2em]
\item the conditional derivation uses the declared Standard Model branch without having to
introduce a supersymmetric partner for every known field;
\item the heat-kernel appendix contains a conditional calibration in which
unification-style running features could arise from edge-sector structure
instead of supersymmetric particle content;
\item the beta-shift calculation uses explicit calibration assumptions rather
than a closed edge-sector theorem;
\item none of this is a universal theorem that supersymmetry is impossible.
\end{enumerate}

So if the paper asks ``why is there no supersymmetry in the derived spectrum?'', the best
technical answer is: because the declared branch stops at the Standard
Model content, and the existing unification discussion is trying to reproduce the relevant
running behavior through a separate heat-kernel calibration, without
extending the particle content to a full
supersymmetric multiplet structure.

\subsection{Three Generations And Flavor Closure}

The target-free structural matter packet leaves the three values in
\(3\leq N_g\leq5\) indistinguishable. On the strengthened screen interface,
the operational Laplacian cost selects the canonical rank-three band uniquely
among the single complete faithful candidates. The declared unitary simulator
recovers the same band at its lowest positive generator frequency. This
supplies an exact finite family candidate.
The fifteen-state charge table has exact anomaly cancellation, diagonal
\(\mathbb Z_6\) action, and a nondegenerate chirality grading. Tensoring the
rank-three response band with that table gives a conditional complex
rank-\(45\) candidate. On a separate finite local domain, the receipt checks
the declared scalar signed operator tensored with \(I_{45}\); its zero kernel
and exact positive dimensionless gap are inherited with multiplicity \(45\).
The source does not select this matter action. The twelve-port Spin packet and
the 8,662-node operator domain have no certified source, domain, or transport
bridge. Matter-pole identification, the continuum Spin/locality limit,
physical seam selection, the missing persistence leg, and laboratory current
identification have no receipts~\cite{ophfpe}.
The full flavor dictionary,
excitation map, and CKM closure are not derived. The charged source landing
from \(P\) to physical charged data is a source-domain no-go, and the neutrino
branch has no source-closed flavor prediction. Its weighted-cycle candidate fails the NuFIT 6.1
correlated profile~\cite{nufit61}.
The flavor and family chapters therefore spend substantial space on
constructive object boundaries as well as final numbers.

\subsection{Ordinary matter and the hadron backend boundary}

A second distinction belongs in discussion form instead of in the results sections: the
difference between ``elementary rows exist'' and ``ordinary matter is fully derived.'' Ordinary
visible matter is dominated by hadrons, especially protons and neutrons. The OPH
derivation has a hadron construction, with source-only hadron masses gated by a production
backend export. The paper emits the eligible elementary rows and retains the weighted-cycle
coordinates only as a rejected comparison record, while the full
observed matter spectrum is not emitted from a source-only OPH backend.

\subsection{Claim boundary for conditional branches}

The reader should therefore interpret this whole discussion section with the same rule used by
\emph{Recovering Observer Spacetime and Einstein Dynamics from Overlap Consistency and Deriving Standard Model Gauge Structure from Observer Overlap Consistency}~\cite{ophEinstein,ophGauge}. The structural gauge and gravity theorems do not depend on a quark source-spread selector, charged-lepton source landing, a source-derived flavor-labelled neutrino kernel, or source-only hadron masses. The independently supported bosonic rows have the same separation.

\section{Conclusion}

Complete reversible response and endogenous proper-carrier transport force
the local Standard Model gauge Lie algebra
\[
\mathfrak{su}(3)\oplus\mathfrak{su}(2)\oplus\mathfrak u(1).
\]
Under the separately declared matrix current and rank-15 matter table, anomaly
balance and tensor descent give the exact hypercharge lattice, three-color
carrier, common \(\mathbb Z_6\) kernel, and maximal faithful image of that
representation. Selection of the physical global quotient is open. The
transportable-sector/Tannaka chain remains a separate conditional
classification. Physical source binding, laboratory current identification,
the two routes' current identity, matter-pole and continuum identification of
the conditional rank-\(45\) candidate, selection of the local matter action,
cross-domain Spin transport, scalar multiplicity/dynamics, and continuum QFT
are open. The
strengthened
finite screen interface selects its rank-three candidate under the stated
premises. Exclusion of extra light sectors is a declared completion.

On the declared quantitative branch, one local coordinate
\(P\) organizes the electromagnetic, electroweak, hierarchy, and Higgs readouts.
The common-load premise gives a dimensionless hierarchy and
\(\epsilon_H=0\). The weak-boson and Higgs numbers are chart coordinates.
The conditional perturbative implication from a complete renormalized packet
to strict charged and neutral complex-pole coefficients is specified. The OPH
source does not produce that packet: its source-selected action, complete
Yukawas, Fleischer--Jegerlehner coordinate~\cite{dudenasLoschner2020,dittmaierRzehak2022},
target-clean running/matching, two independent
engines, executable ST/Nielsen/current-pole receipts, source law/covariance,
physical reference transition, and SI clock binding are work in progress. The
finite local-domain Hamiltonian has an exact positive gap for its signed seam
operator. This finite graph operator is distinct from the compact-gauge repair
generator and the continuum Yang--Mills Hamiltonian. It supplies no physical
clock or mass scale. A nonperturbative chiral gauge construction would
additionally face the known lattice chiral-measure
program~\cite{luscher1999,luscher2000,kadohKikukawa2008} and the
\(\gamma_5\)-scheme restoration bookkeeping~\cite{beluscaMaito2020}. A finite
chiral quantum object and a nonperturbative observable reconstruction are not
prerequisites for the bounded imported-Standard-Model perturbative validation.
Both are work in progress for their stronger claims.

The flavor analysis produces exact algebraic statements and explicit negative
results. The positive-chamber face-circulant identity reduces Koide balance
to \(|b|/a=1/\sqrt2\), and equal rank-two event blocks give that modulus in
the conditional finite tracial-GNS packet. Physical chiral-mass attachment,
phase, and numerical ratios are open. The analysis proves a
two-modulus quark non-identifiability result, and rejects the tested
register-Clebsch, reciprocal-quark, and weighted-cycle neutrino routes. The
Clebsch rejection leaves a conditional channel-pairing boundary and a
target-free unordered weight set under declared constraints; neither supplies
Yukawa-coefficient equality or a source-derived family order. The exact
real-axis enumeration excludes direct identification of the Cabibbo angle
with an acute angle between two of the 31 real three-dimensional icosahedral
residual axes, while leaving spinorial, higher-order, and dynamical routes
outside its scope. The \(W_5\) stabilizer theorem shows why residual symmetry
cannot finish the charged-family problem: threefold and fivefold fixed points
are degenerate, and the twofold fixed locus retains two parameters after
scaling. Absolute
charged-lepton masses require a source-derived clock and
determinant map. Source-only hadron masses require a working hadronic production
backend. Empirical hadronic transport remains identified as empirical input.

The relation between the particle hierarchy coordinate and cosmic capacity is
conditional on two constructions in the companion papers: the direct
public-record fixed point from a complete positive source packet and the
identification of its carrier with the screen/electroweak load. The particle results are stated separately as
theorems, conditional implications, comparisons, rejected candidates, and
requirements for physical identification.

\appendix
\input{tex_fragments/PARTICLE_TECHNICAL_SUPPLEMENT_PORT.tex}

\begin{thebibliography}{9}

\bibitem{observableNormalForms}
B.\ M\"uller, D.\ Matscheko, and J.\ Hill,
\emph{Observation-Determined Normal Forms: Stability, Obstructions, and Refinement in Constraint
and Rewrite Systems},
2026.\\
Available at \url{https://github.com/FloatingPragma/observer-patch-holography/blob/main/extra/observable_normal_forms.pdf}.

\bibitem{ophsynthesis}
B.\ M\"uller, A.\ Osika, M.\ Poneder, K.\ Xue, B.\ Cassie, P.\ Nguyen, M.\ A.\ Visser, K.\ A.\ Anirudha, D.\ Matscheko, and J.\ Hill,
\emph{From Observer Consensus to Standard Physics}.
\url{https://wkaxfdgxoqmghwgshymt.supabase.co/storage/v1/object/public/papers/from_observer_consensus_to_standard_physics.pdf}

\bibitem{ophconsensus}
B.\ M\"uller, K.\ Xue, K.\ A.\ Anirudha, D.\ Matscheko, and J.\ Hill,
\emph{Reality as a Consensus Protocol: The Fixed-Point Computation That Implements Physics}.
GitHub PDF:
\url{https://github.com/FloatingPragma/observer-patch-holography/blob/main/paper/reality_as_consensus_protocol.pdf}

\bibitem{ophmicrophysics}
B.\ M\"uller, A.\ Osika, K.\ Xue, B.\ Cassie, M.\ A.\ Visser, and D.\ Matscheko,
\emph{Federated Echosahedral Screen Microphysics: Patch Hardware, Records, and Observer Synchronization in OPH}.
GitHub PDF:
\url{https://github.com/FloatingPragma/observer-patch-holography/blob/main/paper/screen_microphysics_and_observer_synchronization.pdf}

\bibitem{ophfpe}
B.\ M\"uller,
\emph{OPH-FPE: finite simulator and receipt engine for Observer-Patch
Holography physics experiments},
2026.
\url{https://github.com/muellerberndt/oph-physics-sim}

\bibitem{ophEinstein}
B.\ M\"uller, A.\ Osika, M.\ Poneder, K.\ Xue, P.\ Nguyen, M.\ A.\ Visser, and D.\ Matscheko,
\emph{Recovering Observer Spacetime and Einstein Dynamics from Overlap Consistency}.
GitHub PDF:
\url{https://github.com/FloatingPragma/observer-patch-holography/blob/main/paper/recovering_observer_spacetime_and_einstein_dynamics_from_overlap_consistency.pdf}

\bibitem{ophGauge}
B.\ M\"uller, A.\ Osika, M.\ Poneder, K.\ Xue, P.\ Nguyen, M.\ A.\ Visser, and D.\ Matscheko,
\emph{Deriving Standard Model Gauge Structure from Observer Overlap Consistency}.
GitHub PDF:
\url{https://github.com/FloatingPragma/observer-patch-holography/blob/main/paper/deriving_standard_model_gauge_structure_from_observer_overlap_consistency.pdf}

\bibitem{codata2022}
National Institute of Standards and Technology,
\emph{2022 CODATA Recommended Values of the Fundamental Constants of Physics and Chemistry},
NIST SP 959, May 2024.
\url{https://physics.nist.gov/cuu/pdf/wallet_2022.pdf}

\bibitem{koide1983}
Y. Koide,
\emph{New view of quark and lepton mass hierarchy},
Physical Review D \textbf{28}, 252 (1983).

\bibitem{brannen2006}
C. A. Brannen,
\emph{The Lepton Masses},
preprint, 2006.
\url{http://brannenworks.com/MASSES2.pdf}

\bibitem{antusch2025running}
S. Antusch, K. Hinze, and S. Saad,
\emph{Updated Running Quark and Lepton Parameters at Various Scales},
arXiv:2510.01312, 2025.
\url{https://arxiv.org/abs/2510.01312}

\bibitem{flag2024}
Y. Aoki et al. (Flavour Lattice Averaging Group),
\emph{FLAG Review 2024},
arXiv:2411.04268, 2024,
Tables 11--12.
\url{https://arxiv.org/abs/2411.04268}

\bibitem{pdg2024ckm}
S.~Navas et al. (Particle Data Group),
\emph{CKM Quark-Mixing Matrix},
Review of Particle Physics 2024, Sec.~12.2.2, revised May 31, 2024.
\url{https://pdg.lbl.gov/2024/reviews/rpp2024-rev-ckm-matrix.pdf}

\bibitem{pdg2026}
Particle Data Group,
\emph{Review of Particle Physics: 2026 particle listings},
2026.
\url{https://pdg.lbl.gov/2026/listings/particle_properties.html}

\bibitem{planck18}
Planck Collaboration,
\emph{Planck 2018 results. VI. Cosmological parameters},
Astronomy \& Astrophysics 641, A6, 2020.
\url{https://arxiv.org/abs/1807.06209}

\bibitem{desidr2nu}
DESI Collaboration,
\emph{Constraints on Neutrino Physics from DESI DR2 BAO and DR1 Full Shape},
arXiv:2503.14744, 2025.
\url{https://arxiv.org/abs/2503.14744}

\bibitem{nufit61}
I. Esteban, M. C. Gonzalez-Garcia, M. Maltoni, I. Martinez-Soler, J. P. Pinheiro, and T. Schwetz,
\emph{NuFit-6.0: Updated global analysis of three-flavor neutrino oscillations},
JHEP \textbf{12} (2024) 216, arXiv:2410.05380,
with the NuFIT 6.1 (2025) profile-table release at
\url{https://www.nu-fit.org/?q=node/309}.

\bibitem{kniehl2013}
B. A. Kniehl,
\emph{All-order renormalization of the propagator matrix for fermionic systems
with flavor mixing},
Physical Review D \textbf{89}, 096005 (2014), arXiv:1308.3140.

\bibitem{gambino1999}
P. Gambino and P. A. Grassi,
\emph{The Nielsen identities of the Standard Model and the definition of mass},
Physical Review D \textbf{62}, 076002 (2000), arXiv:hep-ph/9907254.

\bibitem{buchholz1986}
D. Buchholz,
\emph{Gauss' law and the infraparticle problem},
Physics Letters B \textbf{174}, 331--334 (1986).

\bibitem{duch2023}
P. Duch and W. Dybalski,
\emph{Infrared problem in quantum electrodynamics},
arXiv:2307.06114.

\bibitem{luscher1999}
M. L\"uscher,
\emph{Abelian chiral gauge theories on the lattice with exact gauge
invariance},
Nuclear Physics B \textbf{549}, 295--334 (1999), arXiv:hep-lat/9811032.

\bibitem{luscher2000}
M. L\"uscher,
\emph{Lattice regularization of chiral gauge theories to all orders of
perturbation theory},
Journal of High Energy Physics \textbf{06} (2000) 028, arXiv:hep-lat/0006014.

\bibitem{kadohKikukawa2008}
D. Kadoh and Y. Kikukawa,
\emph{A simple construction of fermion measure term in \(U(1)\) chiral
lattice gauge theories with exact gauge invariance},
Journal of High Energy Physics \textbf{02} (2008) 063, arXiv:0709.3658.

\bibitem{beluscaMaito2020}
H. B\'elusca-Ma\"ito, A. Ilakovac, M. Ma\dj or-Bo\v{z}inovi\'c, and
D. St\"ockinger,
\emph{Dimensional regularization and Breitenlohner--Maison/'t
Hooft--Veltman scheme for \(\gamma_5\) applied to chiral Yang--Mills theory},
Journal of High Energy Physics \textbf{08} (2020) 024, arXiv:2004.14398.

\bibitem{dudenasLoschner2020}
V. D\=ud\.enas and M. L\"oschner,
\emph{Vacuum expectation value renormalization in the Standard Model and
beyond},
Physical Review D \textbf{103}, 013003 (2021), arXiv:2010.15076.

\bibitem{dittmaierRzehak2022}
S. Dittmaier and H. Rzehak,
\emph{Electroweak renormalization based on gauge-invariant vacuum expectation
values of non-linear Higgs representations},
Journal of High Energy Physics \textbf{05} (2022) 125, arXiv:2203.07236.

\bibitem{grassiKniehlSirlin2001}
P. A. Grassi, B. A. Kniehl, and A. Sirlin,
\emph{Width and partial widths of unstable particles in the light of the
Nielsen identities},
Physical Review D \textbf{65}, 085001 (2002), arXiv:hep-ph/0109228.

\end{thebibliography}
\end{document}
